% =====================================================================
%  H. Brøchner
%  Benedict Spinoza. En Monographie.
%  Kjøbenhavn: P. G. Philipsens Forlag, 1857.   188 pp.
%  Printer: Thieles Bogtrykkeri.
%
%  DIPLOMATIC TRANSCRIPTION (Danish).
%  Copy digitised: Bayerische Staatsbibliothek, shelfmark Ph.U. 62 m,
%  digitised by Google (»BIBLIOTHECA REGIA MONACENSIS« stamp, PDF 8).
%  This is NOT a Royal Danish Library scan; the consequences are in §6.
% =====================================================================
%
%  ---------------------------------------------------------------
%  1. PAGE MAP  (verified, recon pass, 3 September 2026)
%  ---------------------------------------------------------------
%  The scan is 202 PDF pages: one 600 ppi 1-bit JBIG2 image per page,
%  plus a 1034x204 Google watermark strip at the foot of each. The map is
%  UNIFORM over the whole book, with no offset change anywhere:
%
%      printed p. N   =   PDF page N + 10        (printed 1--188 = PDF 11--198)
%
%  Pagination is CONTINUOUS from the Fortale: printed 1--14 is the Fortale
%  and 15--188 the monograph, in one arabic sequence. There is no roman
%  front-matter run.
%
%  Front matter: PDF 1 Google notice, PDF 3--4 the BSB shelfmark leaf,
%  PDF 7 title page, PDF 8 the BSB stamp, PDF 9 Indhold, PDF 10 blank.
%  PDF 199--202 are blank leaves and cover. THERE IS NO ERRATA LEAF and no
%  publisher's advertisement, so there is no authorial list of corrections:
%  every anomaly recorded below is the compositor's.
%
%  Verified by reading the printed folio off the ABBYY layer of all 188 body
%  pages; it agrees with +10 on 183 of them, and the five that yield no folio
%  are exactly the five that carry none in the print -- p. 1 (Fortale
%  opening), p. 15 (divisional half-title), p. 16 (its blank verso), p. 17
%  (Indledning opening) and p. 27 (Første Capitel opening). Each is fixed
%  from below by its neighbour. Four of the five were then confirmed on the
%  image.
%
%  Chapter openings, from the heads themselves, cross-checked against the
%  Indhold:
%      Fortale         1     Femte Capitel    82
%      Indledning     17     Sjette Capitel  108
%      Første Capitel 27     Syvende Capitel 132
%      Andet Capitel  44     Ottende Capitel 162
%      Tredie Capitel 65     Niende Capitel  171
%      Fjerde Capitel 71     Tiende Capitel  181
%  Caps 4, 5, 6, 7, 9 and 10 begin PART WAY DOWN a page that already carries
%  the end of the previous chapter; Caps 1, 2, 3, 8 and the Indledning begin
%  at the head of their page.
%
%  ---------------------------------------------------------------
%  2. WHAT IS TRANSCRIBED, AND HOW
%  ---------------------------------------------------------------
%  Diplomatic. 19th-century orthography exactly as printed: Philosophie,
%  Consequentser, Substantsbegrebet, høiere, Øieblikket, aa, capitalised
%  nouns. Printer's errors are transcribed AS PRINTED, logged in a % comment
%  at the site, and never silently corrected.
%
%  TYPEFACE. The whole book is ANTIQUA -- title page, Indhold, heads, body
%  and notes alike. There is no Fraktur anywhere in it. This makes it the
%  odd one out among the Brøchner volumes in this repository, and several
%  conventions below are consequently the REVERSE of the sibling volumes'.
%
%  THE BOOK HAS NO å AND NO Å. Danish å is spelled aa throughout. Any Å in
%  a draft is an uncorrected OCR reading of a capital A.
%
%  ---------------------------------------------------------------
%  3. EMPHASIS, LATIN, GREEK
%  ---------------------------------------------------------------
%  EMPHASIS: the book's only emphatic device inside a paragraph is
%  letterspacing (Sperrsatz). All of it is rendered \emph{}. Bold occurs
%  ONLY in display heads, on the title page and in the Indhold; it never
%  appears inside a paragraph. Small caps are not used at all.
%    Letterspacing is applied to Danish, to Latin and to proper names
%    indifferently -- there is no separate device for names.
%    As in the sibling volumes the compositor is inconsistent about the
%    extent of a spaced run: p. 71 prints  n a t u r a  naturans  with only
%    the first word spaced. Transcribe AS PRINTED and flag it in a comment;
%    do not regularise.
%
%  LATIN IS SET IN PLAIN ROMAN, indistinguishable from the Danish, and so
%  are the Ethics citations (Eth. II, 11. Coroll.), the French and the
%  German. IT IS THEREFORE NOT ITALICISED HERE. This reverses the practice
%  of `problemet-tro-viden' and `det-religioese', where the Latin stands in
%  antiqua against Fraktur and is rendered \textit{}; in an all-antiqua book
%  there is no such distinction to record, and inventing one would be an
%  editorial claim the page does not support. \textit{} therefore appears
%  in the body of this file exactly ONCE, for the Greek (below); check.py
%  flags every other occurrence.
%
%  GREEK: ONE passage only, printed p. 63 -- the Aristotelian tag
%  »τοῦ μὴ ὄντος …« -- and the book sets it in an accented ITALIC Greek
%  fount. That italic is the print's own, not an editorial italicisation of
%  a foreign phrase, so it is rendered \textit{} and is the single sanctioned
%  exception to the rule above. The ABBYY layer reads the passage;
%  tesseract -l dan does not.
%
%  THE OLD ID-EST MARK ɔ: (a reversed c followed by a colon) is used
%  constantly -- about 75 instances, from p. 37 to p. 170, concentrated in
%  Caps 2--8. It is typed as U+0254 + colon and rendered \reflectbox{c}.
%
%  ---------------------------------------------------------------
%  4. QUOTATION MARKS -- TWO PAIRS, BOTH GENUINE
%  ---------------------------------------------------------------
%  This book uses BOTH »…« and „…“, in the body and in the footnotes alike,
%  and it mixes them: the pair »…“ in note 1 on p. 148 is what the page
%  prints. Neither pair is »the book's« mark, and neither is normalised to
%  the other. There is no space inside the marks; the spaces the ABBYY layer
%  puts there are an artefact of letterspacing (§6).
%    Consequence for verification: NEITHER quote balance is expected to come
%  out at zero. The mixed pair alone leaves the guillemets at +1 and the
%  low/high pair at -1 over the whole book. check.py reports the two counts
%  separately; watch the delta between splices, and log every deliberate
%  imbalance in RESUME-NOTES.md so that nobody later »fixes« it.
%
%  Em-dash: ---. Measured on the 600 dpi renders of pp. 28, 71 and 108, the
%  book's long dash is 78--86 px against an 83 px em, so it is an em dash;
%  the hyphen is 29--33 px. The Indhold's page ranges use an en rule (--).
%
%  ---------------------------------------------------------------
%  5. FOOTNOTES
%  ---------------------------------------------------------------
%  134 of the 188 pages carry at least one. The mark is a SUPERSCRIPT ARABIC
%  NUMERAL FOLLOWED BY A CLOSE PARENTHESIS -- ¹) ²) ³) ⁴) ⁵) -- and THE
%  NUMBERING RESTARTS AT 1 ON EVERY PRINTED PAGE (up to five on p. 156).
%  There is no *) anywhere in the book.
%    Hence the preamble sets \thefootnote to \arabic{footnote}) and every
%  page marker that opens a page carrying a note is followed by
%  \setcounter{footnote}{0}. footmisc[perpage] would be wrong: it resets on
%  the TYPESET page, not the printed one.
%    Notes are set in smaller roman, the mark at the left margin, turnovers
%  indented, under a short rule flush left about one fifth of the measure --
%  omitted where the page is tight (p. 96 has notes and no rule).
%    NOTES RUN OVER THE PAGE BOUNDARY: p. 130's note 2 continues at the head
%  of p. 131's note block, and p. 95's notes 3 and 4 at the head of p. 96's.
%  The whole note belongs to its mark, not to the page it finishes on.
%    TRAP: 1) 2) 3) 4) also occurs as an in-text enumeration in the body, on
%  pp. 95--96 (Spinoza's four modes of perception) and p. 108 (the four uses
%  of the doctrine). Those are not footnotes.
%
%  ---------------------------------------------------------------
%  6. WHAT THIS SCAN CANNOT SETTLE
%  ---------------------------------------------------------------
%  The file is a Google Books JBIG2 scan with symbol-matching compression:
%  199 /JBIG2Globals streams, and on a sample page 57--59% of glyph
%  components are pixel-identical to another instance of the same letter
%  (largest class: 27 identical `e'). The scan therefore RECONSTRUCTS its
%  letters from a dictionary rather than reproducing each impression.
%    NO SORT-LEVEL CLAIM MAY BE MADE FROM IT: no wrong-sort inventory, no
%  broken-sort log, no argument about the compositor's case. This is the
%  hard version of TRANSCRIPTION-PLAYBOOK §2, and it is not a matter of
%  rendering at a higher resolution -- the evidence is not in the file.
%    Structural defects -- dropped or repeated words, turned lines, wrong
%  punctuation, a missing quotation mark -- survive compression and remain
%  fair game, logged as printed.
%    TESTED, NOT ASSUMED. The clustering was calibrated on the three pages
%  where batches had claimed turned sorts (printed 9, 36, 39) plus a
%  neighbour of each: 664 glyphs read as `n' and 161 read as `u' were cropped
%  from the NATIVE 1-bit page bitmaps and hashed. 28% of the n's and 20% of
%  the u's are pixel-identical to an earlier instance -- impossible under
%  lossless coding, so the dictionary is confirmed. And TWO BITMAPS OCCUR IN
%  BOTH ROLES: one stands in `sur', in `under Titelen' and in `tractatus',
%  all on printed p. 36; another in `Middelpunct' (p. 9), `conseqvente'
%  (p. 9) and `Absolute' (p. 10). The words are unambiguous, so a single
%  symbol is being placed where the text requires n and where it requires u.
%  The n/u substitution mechanism is therefore ACTIVE IN THIS FILE, not
%  merely possible in principle. Full method, numbers and the one measurement
%  that failed are in JBIG2-TEST.md.
%
%    STANDING RULE FOR THE BOOK, from that result: WHERE A READING IS NOT A
%  WORD IN DANISH OR LATIN AND DIFFERS FROM A REAL WORD ONLY BY n<->u, THE
%  SENSE-READING IS SET and the printed appearance is recorded in a comment
%  at the site. Three sites are so treated -- `Element' (p. 9), `under'
%  (p. 36), `Foruden' (p. 39), where the scan shows `Elemeut', `nnder',
%  `Fornden'. Earlier drafts logged all three as turned sorts, with confident
%  arguments from the glyph at 600 dpi and from both OCR witnesses agreeing;
%  both arguments are void here, the second because the two witnesses read
%  the same reconstructed image and so are not independent.
%
%    Two further apparent misprints NOT logged as defects, same reason:
%  `Stcd' for Sted (p. 71) and `Vanskcligheden' (p. 50), where both OCR
%  witnesses and the image show a weak crossbar on an `e'. The sense-reading
%  is transcribed and the doubt is recorded; the printing house is not
%  accused.
%
%  ---------------------------------------------------------------
%  7. NORMALISATION
%  ---------------------------------------------------------------
%  Almost none is needed, the book being antiqua: there is no I/J single
%  sort to resolve, as there was in the Fraktur volumes.
%    ø / ö: the body, the heads and the Indhold all print the slashed ø/Ø.
%  ONLY the title page's display capitals are ö -- BRÖCHNER, KJÖBENHAVN --
%  and that is the printer's own setting, verified at 600 dpi. It is
%  transcribed as printed, and not carried into the body.
%    Spaces before commas and semicolons in a draft are an ABBYY artefact of
%  letterspaced lines, not the print. Set ordinary punctuation.
%
% =====================================================================

\documentclass[12pt,a4paper]{book}
\usepackage[utf8]{inputenc}
\usepackage[T1]{fontenc}
\usepackage{libertinus}
\usepackage{libertinust1math}
\usepackage{textalpha}   % Greek (the p. 63 passage), typed directly
\usepackage{graphicx}    % for \reflectbox (the old Danish »ɔ:« id-est mark)
\DeclareUnicodeCharacter{0254}{\reflectbox{c}}  % ɔ (U+0254) = reversed c
\usepackage[danish]{babel}
\usepackage[protrusion=true,expansion=true]{microtype}
\usepackage{setspace}
\usepackage[margin=1.2in]{geometry}
\usepackage{enumitem}
\usepackage{fancyhdr}
\setlength{\headheight}{28pt}
\addtolength{\topmargin}{-13.5pt}
\usepackage{hyperref}

\hypersetup{
  pdftitle={Benedict Spinoza. En Monographie},
  pdfauthor={H. Brøchner},
  hidelinks
}

% The book prints ¹) ²) ³) ⁴) ⁵), restarting on every printed page; hence
% \setcounter{footnote}{0} at the marker of every page that carries a note.
\renewcommand{\thefootnote}{\arabic{footnote})}

\pagestyle{fancy}
\fancyhf{}
\fancyhead[LE,RO]{\thepage}
\fancyhead[RE]{\textit{Benedict Spinoza. En Monographie}}
\fancyhead[LO]{\textit{\leftmark}}

\setstretch{1.3}

\title{\textbf{Benedict Spinoza}\\[10pt]
  \large En Monographie}
\author{Dr.\ H.\ Brøchner}
\date{Kjøbenhavn: P.\ G.\ Philipsens Forlag, 1857}

\begin{document}
\maketitle
\thispagestyle{empty}
\newpage

% =====================================================================
%  TITLE PAGE  (PDF 7), as printed. Everything on it is antiqua; the
%  display lines are letterspaced, hence \emph{}. Note ö, not ø, in
%  BRÖCHNER and KJÖBENHAVN -- the printer's own setting on this leaf only
%  (header §7).
% =====================================================================
\chapter*{Titelblad}
\addcontentsline{toc}{chapter}{Titelblad}
\markboth{Titelblad}{Titelblad}

\begin{center}
{\Large\textbf{\emph{BENEDICT SPINOZA.}}}\\[8pt]
\rule{2cm}{0.4pt}\\[8pt]
{\large\emph{EN MONOGRAPHIE}}\\[6pt]
AF\\[6pt]
\textbf{Dr.\ H.\ \emph{BRÖCHNER}.}\\[10pt]
\rule{1.5cm}{0.4pt}\\[10pt]
KJÖBENHAVN.\\[4pt]
\textbf{\emph{P.\ G.\ PHILIPSENS FORLAG.}}\\[4pt]
{\small THIELES BOGTRYKKERI.}\\[4pt]
1857.
\end{center}

\newpage

% =====================================================================
%  INDHOLD  (PDF 9)
%  Read from the image at 300 and 600 dpi, not from the OCR layer
%  (playbook §6). `Indhold.' is letterspaced bold and centred, with a short
%  centred rule below; `Side' is a small bold column head at the right; the
%  Fortale entry is divided from the rest by a second short rule; the
%  `Cap. N.' labels and the chapter arguments are bold; leaders are spaced
%  full stops; ranges use an en rule; NO ENTRY ENDS IN A FULL STOP (the
%  printed heads do -- a setting convention, not a variant).
%
%  TWO THINGS RECORDED RATHER THAN NORMALISED:
%   (a) The Indhold's `Indledning' has NO SUBTITLE, but the printed head on
%       p. 17 reads `Indledning.' followed by `Spinozismens historiske
%       Genesis.' Both readings stand.
%   (b) The ranges are inconsistent about sharing an endpoint -- 27--43 /
%       44--64 do not overlap, while 65--71 / 71--82 do. That is not an
%       error: it tracks exactly which chapters begin part way down a page.
%
%  `naturæ pars' in the Cap. 6 entry is in the same bold roman as the rest
%  of the entry, not italic and not letterspaced; likewise in the head on
%  p. 108. The second `1' of Cap. 9's `171' has lost its foot serif and both
%  OCR witnesses read `17!'; compared at 600 dpi with the sound `1' two
%  characters earlier it is a 1.
% =====================================================================
\chapter*{Indhold}
\addcontentsline{toc}{chapter}{Indhold}
\markboth{Indhold}{Indhold}

\begingroup
\setstretch{1.15}
\begin{list}{}{%
  \setlength{\leftmargin}{3.6em}\setlength{\itemindent}{-3.6em}%
  \setlength{\itemsep}{2pt}\setlength{\parsep}{0pt}}

\item[] \hfill Side

\item \textbf{Fortale} \dotfill 1--14

\item[] \vspace{2pt}

\item \textbf{Benedict Spinoza} \dotfill 15--188
\item \textbf{Indledning} \dotfill 17--26
\item \textbf{Cap.\ 1.\ Spinozas Liv og Skrifter} \dotfill 27--43
\item \textbf{Cap.\ 2.\ Spinozismens Grundbegreber} \dotfill 44--64
\item \textbf{Cap.\ 3.\ De nærmeste Consequentser af Spinozas Opfattelse af
  Substantsbegrebet} \dotfill 65--71
\item \textbf{Cap.\ 4.\ Spinozas Naturbegreb} \dotfill 71--82
\item \textbf{Cap.\ 5.\ Den menneskelige Aand} \dotfill 82--108
\item \textbf{Cap.\ 6.\ Mennesket som naturæ pars. Affecterne. Menneskets
  Trældom} \dotfill 108--132
\item \textbf{Cap.\ 7.\ Naturretten. Religionen} \dotfill 132--161
\item \textbf{Cap.\ 8.\ Menneskets Frihed og Evighed} \dotfill 162--171
\item \textbf{Cap.\ 9.\ Charakteristik og Kritik af Spinozismen} \dotfill 171--181
\item \textbf{Cap.\ 10.\ Spinozas Forhold til den sildigere Philosophie} \dotfill 181--188

\end{list}
\endgroup

\newpage

% =====================================================================
%  FORTALE  (printed pp. 1--14 = PDF 11--24)
%  Head `Fortale.' is letterspaced bold, centred, with a short centred rule
%  below and a two-line raised initial D on the first word. The initial is
%  transcribed as an ordinary letter and recorded here.
%  No footnotes before p. 13.
% =====================================================================
\chapter*{Fortale}
\addcontentsline{toc}{chapter}{Fortale}
\markboth{Fortale}{Fortale}

% --- p. 1 ---
% The BSB ink stamp stands at the head of this page; it is not text.
% `Fortale.' is letterspaced bold, centred, over a short centred rule (both in
% the skeleton). The first word carries a two-line raised initial D,
% transcribed here as an ordinary letter.
Den Fremstilling af den spinozistiske Philosophie, som jeg i
dette Skrift har givet, er bestemt til at begynde en Række af
Monographier, der, medens de hver for sig danne et Hele, tillige i
Forening skulle give en Udsigt over hele Philosophiens Historie. Jeg
har ikke villet forsøge paa med lige Udførlighed at behandle denne i
dens hele Sammenhæng; jeg troer, at en Enkelts Kræfter ere
utilstrækkelige til et saadant Arbeide, naar det skal være grundet paa
et virkeligt Studium af Kilderne, paa en paalidelig Autopsie;
idetmindste føler jeg, at mine Kræfter vilde være det. Ogsaa antager
jeg, at det for Øieblikket, efter de mange Forsøg, der med større og
mindre Dygtighed og Held i de sidste 30 Aar ere gjorte i denne
Retning, vilde være mindre fornødent, og at det nu mere kommer an
paa, ved en ny omhyggelig Behandling af de enkelte Hovedpunkter at
skaffe et paalideligere Materiale tilveie for en gjentagen
sammenfattende Behandling af hele Stoffet. En fragmentarisk Behandling,
der dog, ligesom fra de fremragende Høider, giver et Overblik over det
hele Omraade, vil maaskee ogsaa, idet den samler Blikket paa det
Væsentlige, paa Tænkningens historisk udfoldede Grundformer, der ikke
er noget blot Forbigangent, som tilhører Erindringen og den lærde
Forskning, men et indenfor sin Grændse bestandig Berettiget og
Blivende, der beholder sin Betydning til alle Tider, og idet den
fremstiller de enkelte Tænkeres Systemer i en saadan Ud\-%
% --- p. 2 ---
førlighed, at der kan aabnes et dybere Indblik i disses indre
Drivefjedre, saa at de sees som Udfoldelsen af en \emph{levende}
Tanke, mere end den sammenhængende Fremstilling være skikket til at
bane Veien, ikke blot for en grundigere Forstaaen, men ogsaa, gjennem
den, hos Enkelte for nogen Interesse for en Gjenstand, som Øieblikkets
Stemning og Stræben ikke er gunstig. Ingen, der har gjort Philosophien
til sit Livs Syssel, vil vel have kunnet undgaae Tvivl og Mismod ved
at see den Videnskab, der for faa Decennier siden var Gjenstand for et
enthusiastisk, forhaabningsfuldt Studium, ved Skuffelsens Reaction at
blive Gjenstand for Ligegyldighed, og næsten at skjules under
fremmede, tidt ideeløse Interesser. Men som den menneskelige Aands
Selvbesindelse er Philosophien dens umistelige Eiendom, der, engang
vunden, aldrig vil blive opgivet; den er som de Floder, der paa en
Strækning af deres Løb forsvinde under Jorden, men kun for atter at
komme tilsyne, mægtigere og rigere efter Foreningen med usynlige
Kilder.

Idet jeg gjør det til min Opgave, ved en fragmentarisk Fremstilling
af Philosophiens Historie at give en Udsigt over den i dens hele
Omfang, har mit Valg maattet falde paa de Systemer, der betegne
Høidepunkter, fra hvilke der paa een Gang aabner sig Udsigt tilbage
over en længere, afsluttet Udvikling, og fremad over en ny, fra det
fremstillede System udgaaende eller i væsentlig Grad ved det
paavirket, Række. Med dette Hensyn falder \emph{det} sammen, at
udvælge saadanne Systemer, der ved deres Princip og Retning
repræsentere Hovedformer af den menneskelige Tænkning; thi kun
saadanne træde i hint Forhold til Fortid og Eftertid. Den græske
Philosophie vil blive repræsenteret ved \emph{Platon},
\emph{Aristoteles} og \emph{Plotinos}, Renaissancetiden ved det
% `og' between the spaced names is NOT letterspaced (600 dpi), here and
% throughout the Fortale; nor is `Renaissancetiden', nor this `det'
% (600 dpi) -- the parallel phrase on p. 12 DOES space its `det'.
\emph{platoniske Akademie i Florents} og \emph{Giordano Bruno}; den
nyere Philosophie i dens ene Udviklingsrække ved \emph{Bacon} og
\emph{Système de la nature}, i dens anden ved \emph{Spinoza} og
\emph{Leibnitz}, og paa dens sidste Udviklingsstadium ved \emph{Kant}
og \emph{Hegel}. Ved Indledninger og Slutnings\-%
% --- p. 3 ---
afhandlinger til de enkelte Fremstillinger vil Sammenhængen mellem det
Fragmentariske blive tilveiebragt.

Den fuldstændige Begrundelse af Valget af de repræsenterende Systemer
maa selve Fremstillingen af dem give. Dog maa jeg her i Korthed, og
mere antydende end udførende, fremsætte den Opfattelse af Gangen i
Philosophiens Historie, der ligger til Grund for Valget, og det
saameget mere, som jeg begynder med et af Systemerne midt inde i
Rækken, og der altsaa under Fremstillingen af det ikke vil være nogen
naturlig Plads for en saadan Udvikling.

Efter min Opfattelse deler hele Philosophiens Historie fra dens
første Begyndelse og indtil vore Dage sig kun i to store,
\emph{afsluttede} Afsnit: den græske Philosophie og den nyere Tids. De
ere adskilte ved et stort Mellemrum, maaskee vil et lignende lægge sig
mellem vor Tids Philosophie og Fremtidens, der vil støtte sig paa en
rigere Udvikling af de enkelte Videnskaber, navnlig
Naturvidenskaberne. Den \emph{græske} Philosophie betragter jeg som
% only `græske' is spaced here; `Philosophie' beside it is not (600 dpi).
Philosophiens Begyndelse; thi Orienten kjendte kun Phantasiens, ikke
Tankens Uendelighed. Grækenlands philosophiske Gjerning var det, med
Orientens Naturliv til Forudsætning, at opdage \emph{Ideen},
anskuende at begribe Universet som det uendelige, levende, harmoniske
\emph{Kunstværk}, der fik sin sande Virkelighed ved Formen, Ideen, den
evigt værende plastiske Tanke. \emph{Skjønhedens} Tanke var
Grækenlands Grundtanke, den beherskede den græske Tænkning som den
beherskede det græske Liv, Livet i Staten og Individets sædelige Liv.
Den græske Tankes Opløsning er begrundet i den indre Modstrid i
Skjønhedens \emph{umiddelbare} Forening af sine Elementer: Formen og
Stoffet, i Umuligheden af at bringe Stoffet, Materien, til at gaae op
i Formen, at bringe »det Ikke-Værende« til at forsvinde ligeoverfor
det Værende, der har al Virkelighed i sig, men bestandig trænger til
Stoffet forat kunne realisere denne Virkelighed. Denne indre
Modsigelse er det, som den græske Tanke efter Aristoteles bliver sig
bevidst og forgjæves søger at løse indenfor sit Omraade, indtil den i
Nyplatonismen løfter sig op over sin
% --- p. 4 ---
Grundvold, Naturen, for i en høiere Region at finde en Forsoning, den,
efter sin Eiendommelighed, ikke kunde finde der, men bringer
Modsigelsen med sig over i det nye Element, og, i en ny Aands
Tjeneste, udtømmer sin sidste Kraft i et frugtesløst Forsøg. For denne
Udvikling blive de naturlige Repræsentanter at søge i \emph{Platon},
der i sig sammenfatter hele den ældre græske Philosophie, først
udtaler Grækenlands Grundtanke, og først drager Universet ind under
Ideens Herredømme; --- i \emph{Aristoteles}, der bestemmer den sidste
ægte græske Formel for Enheden af Form og Stof, der ved den supplerer
Platon, der, medens Alexander undertvinger en Verden for Grækenland,
vinder en Verdens Rigdom for den græske Tanke, og med den giver et
Aartusinde Næring; --- og endelig i \emph{Plotinos}, den største
græske Tænker i den Tid da Grækenland var bleven fremmed for sig selv,
da et nyt Livsprincip allerede bevægede Verden, og den græske Tankes
sidste energiske Forsøg paa at hævde sig selv kun viste dens Afmagt og
bebudede den nye Tids Seier.

Det hele Tidsrum fra Grækenlands Undergang til Videnskabernes
Gjenfødelse indtager ingen Plads i Philosophiens Historie. Jeg troer,
at jo mere man, uhildet af uklare Tendentser, vender tilbage til
Studiet af denne Tid, desto bestemtere vil man komme til den Indsigt,
at alt det \emph{Tanke-}Indhold, som den christne Middelalder tærede
% PARTIAL LETTERSPACING, verified at 600 dpi: `T a n k e -' is spaced,
% `Indhold' on the next line is not. Transcribed as printed, not regularised.
% The hyphen is the compound's own (the second element is capitalised);
% the book sets uncompounded `Tankeindhold' on pp. 6 and 12.
paa, skyldes den græske Philosophie, og at man endog kan udvide denne
Sætning saaledes, at den i alt Væsentligt ogsaa gjælder for
Renaissancetiden. Christendommen kunde og vilde ikke ud af sig udvikle
en Philosophie; den havde fornægtet de Kræfter, af hvilke en saadan
skulde udvikle sig, forladt den Grund, hvoraf den skulde spire frem.
Idet Christendommen vil fuldbyrde hvad Grækenland ikke havde kunnet
fordi det hvilede i Skjønhedens Tanke: at løsrive Aanden som det ene
Sande og Virkelige fra Naturen, at sætte denne som det Væsenløse, det
Onde, at fordre Død som Betingelse for Livet, at flytte Virkeligheden
over i et Rige, der »ikke er af denne Verden«, sætter den sig i et
blot negativt Forhold
% --- p. 5 ---
til det menneskelige og naturlige Livs Former og Kræfter, dens
Fordring bliver bestandig kun den: at fornægte dem. Der kan derfor
ikke være Tale i \emph{stræng Forstand} om en christelig Philosophie,
saalidet som om en christelig Kunst, en christelig Stat, en christelig
Ethik. Det er en Misbrug af Ordene naar man taler saaledes; man gjør
det kun idet man, forsætlig eller uforsætlig, overseer, at
Christendommen ikke er kommen ind i en \emph{ny} Verden, men at den
kom ind i en Verden, i hvilken en stor Culturperiode netop
afsluttedes, at den forefandt hele Oldtidens rige Arv, og at den, da
den første christne Tids consequente Haab om en nærforestaaende
Tilintetgjørelse af den naturlige Virkelighed skuffedes, maatte finde
sig tilrette i den \emph{fremmede} Virkelighed, tale dennes Sprog og
bruge dens Rigdom, dog uden nogensinde at kunne eller ville indgaae en
Enhed med den, saaledes at Alt, hvad der af Tænkning, af Kunst, af
Former for et ethisk Menneskeliv har udviklet sig \emph{samtidig med}
Christendommen, ikke er udgaaet af \emph{denne}, men kun
% `af' before `denne' is not spaced; the `af' inside the next run is
% (600 dpi).
\emph{under Paavirkning af den}, hyppigt under Kamp imod den, af
Menneskelivets Naturgrund, befrugtet ved Oldtidens Tanke.

Det var først i det femtende Aarhundrede, at Tænkningen, der i sin
Trældom havde vundet Styrke, vovede sig til, beaandet ved Oldtidens
nye Aabenbaring, og baaren af et rigtudfoldet verdsligt Liv, og af en
Natur, der ikke længere var Aanden fremmed, uafhængig af Kirkens
Autoritet at udvikle en Philosophie. Det var endnu en
Overgangsperiode; endnu maatte Tanken kæmpe med Phantasien om
Herredømmet, og havde ikke Kraft til, af sig selv alene at udvikle et
nyt Indhold; men i Oldtidens Philosophie fandt den et Indhold, i
hvilket Menneskeaanden gjenkjendte sig selv i sin Virkelighed, og med
Begeistring greb den dette, og formede det paany. Det var Tænkningens
bacchantiske Periode: med en uhyre Tillid, med en urokkelig \emph{Tro
paa Mennesket} svælgede den i den nye Rigdom, og formede, i Tilslutning
til et \emph{selvvalgt} Forbillede, nye Systemer. Tankeindholdet i
disse er Oldtidens, saaledes som det var
% --- p. 6 ---
tilberedt gjennem de sidste græske Philosopher: det var en storartet
pantheistisk Anskuelse af det guddommelige Universum, i hvilket
Mennesket indtog en ophøiet Plads som Foreningspunctet mellem
Sphærerne, som en Mikrokosmos, der afspeilede Universet i sig.
Systemernes Characteer var endnu overveiende theologisk, men
forskjellig fra Scholastikens. Naturen udøvede en stærkere Tiltrækning
paa Aanden, og Anskuelsen af den skjød sig, bevidst eller ubevidst,
mere og mere ind under Anskuelsen af det Guddommelige. Universet blev
en harmonisk Udfoldelse af Guddommens Liv, og Troen paa Gudmennesket
omformede sig til Troen paa den guddommelige, mikrokosmiske
Menneskeaand.

Men idet Tanken saaledes droges henimod Naturen, der gjennem hele
Middelalderen havde været banlyst, og idethøieste kun havde været en
Gjenstand for Phantasien, maatte paa den ene Side Blikket aabnes for
de Naturbetingelser, under hvilke Menneskeaanden staaer; paa den anden
Side maatte Striden mellem den nye Retning og den overleverede, der
endnu bevarede en hemmelig Magt, føre til en dualistisk Adskillelse
mellem det Naturlige og det Overnaturlige, mellem Legeme og Aand,
mellem en sensualistisk og en idealistisk Erkjendelsestheorie, mellem
en materialistisk og en spiritualistisk Opfattelse af Tilværelsens
Principer. Gjennem denne Dualisme, der maatte medføre
Utilfredsstillethed og Skepsis, lededes Opmærksomheden hen paa
Nødvendigheden af et fast Udgangspunkt for Tanken, en sikker Methode
for dens Udvikling, og paa dette Punkt begynder \emph{den nyere Tids
Philosophie}, der afbryder Oldtidens Tradition, og selvstændig og
selvbevidst vil udfolde et nyt Tankeindhold.

Som Arv fra Middelalderen havde den nye Philosophie modtaget
Forudsætningen af en Dualisme mellem Aanden og Naturen, en Dualisme,
der i Principet var givet ved Christendommen. For Tanken, der
selvstændig søgte efter, Sandhed, laae der to adskilte Sphærer,
% PRINTER'S DEFECT, as printed: a comma stands between `efter' and its
% object `Sandhed', at the line end. Verified at 600 dpi; both OCR
% witnesses read it. Not corrected.
ligesom to Verdener; men i Selvbevidsthedens Enhed ligger en Protest
mod enhver Dualisme, en Fordring paa Erkjendelsens Enhed og
% `paa' (last word but two of the page): the third sort is an ink-filled or
% damaged `a' and reads as e/v at 600 dpi; both witnesses give `pae'. The
% sense-reading `paa' is transcribed. NOT logged as a defect -- this JBIG2
% scan cannot settle a single sort (RECON §14).
% --- p. 7 ---
Sammenhæng. Men den dualistiske Forudsætning maa bestemme denne Enhed
saaledes, at den ene af Modsætningerne opsluger den anden, at i
Foreningen den enes Eiendommelighed opgives. Det vil vise sig at være
Tilfældet i de Systemer, der betegne denne philosophiske Epoches
Slutningspunkter. Og den Consequents, der i dem er dragen, er de
tidligere, paa samme Grundforudsætning hvilende, Systemers bevidste
eller ubevidste Drivefjeder; ved det sidste Maal forstaaer man deres
Stræben.

Den første Form, hvorunder Dualismen fremtræder, bliver den, at den
philosophiske Erkjendelse tager et dobbelt Udgangspunkt, fra Naturen
eller fra Aanden, og dette dobbelte Udgangspunkt, hvorved
Erkjendelsens Characteer og væsentlige Gjenstand bestemmes væsenlig
forskjelligt, constituerer for den nyere Philosophie i dens
fremskridende Udvikling en dobbelt Række, der imidlertid, efter Sagens
Natur, strax maatte faae Berøringspunkter, idet begge gik ud fra
Forudsætningen af Aandens Magt til at erkjende, idet Naturen for begge
beholdt sin mægtige Tiltrækningskraft, og begge tillige begyndte med
at laane deres Methode fra Naturvidenskaberne og Mathematiken.

\emph{Bacon} og \emph{Descartes} betegne det dobbelte Udgangspunkt for
Tænkningen. \emph{Bacon} sætter den mod Naturen og dens reale Fylde
rettede \emph{Erfaring} som Erkjendelsens Princip, han gjør Naturen
til Tænkningens væsentlige Gjenstand, til Sandhedens Kilde,
Naturphilosophien til Videnskabernes Moder. Men han vil ikke lade
Erfaringen tabe sig i det Enkelte, han vil finde de skaberiske Former,
Naturens Enheder, det Almindelige i det Concrete, og han anerkjender
Tænkningens Activitet og Idealismens Fordring. Men i
Erfaringsprincipet ligger Faren for den Consequents, at Jeget
forholder sig passivt, at det, som blot optagende gjennem Sandserne,
kun optager det Endelige, at dette, der opfattes strængt som Endeligt,
opfattes mechanisk, og at Mechanismen ogsaa overføres paa Aanden.
Præmisserne til alle disse Consequentser ere allerede givne hos
\emph{Hobbes}, tildels ere de ogsaa dragne og denne
Erkjendelsesretnings
% --- p. 8 ---
sidste Maal antydet; men en Rest af et rationalistisk Element hos ham,
og det, at hans Interesse væsentligen vendte sig mod det Practiske,
har forhindret, at han blev for Philosophiens Udvikling hvad senere
\emph{Locke} blev.

Medens Empirismen saaledes hurtigt udviklede sine Consequentser, var
den modsatte Retning, der ud fra det tænkende Jegs Indre vil udvikle
Tankens Verden, udpræget i betydelige Systemer. \emph{Descartes} søgte
gjennem Tvivlen om Alt, gjennem Udskilningen af alt Fremmedartet,
tilbage til Tvivlens urokkelige Forudsætning, Jeget, der bestemtes som
Eet med Tænkningen og som sondret fra Materie og Udstrækning. Men
forat faae et Indhold maa han forlade dette Enhedspunkt, han maa tye
tilbage til det, Jeget har skilt sig fra som sin Modsætning, og som
det end ikke kan \emph{tænke} som sin Modsætning, uden at ophæve det
abstracte Modsætningsforhold. Gjennem Gud kommer Tanken tilbage til
det Materielle; Sjælen anerkjender, ligesom nødtvungen, sin
Forbindelse med Legemet, om end denne Forbindelse reduceres til et
Minimum af Rum og til Paavirkningens abstracteste Form. Men
Forbindelsen af Aand og Legeme bliver atter opfattet som en
\emph{Sammensætning} af to Uensartede, den tænkende og den udstrakte
Substants vedblive at staae adskilte, og \emph{Geulincx} og
\emph{Malebranche} drage Occasionalismens Consequentser. Den
guddommelige Substants bliver hos Descartes staaende udenfor
Virkeligheden, adskilt fra de to Substantser, der udfylde denne, og
kun beslægtet med den tænkende Substants. Tankens Stræben efter
Erkjendelsens Enhed er saaledes paa to Punkter hæmmet. Men
\emph{Spinoza} hæver Tvetydigheden i Substantsbegrebet (jfr.
Indledningen): den guddommelige Substantses uendelige Enhed udfylder
Alt, og derved blive de bestemte Substantser satte som dens
Attributer. Tænkning og Udstrækning sættes i Parallelismens Enhed som
ligeberettigede, og Overvægten er endog ifærd med at gaae fra Tanken
over til Materien, Gud er ifærd med at falde sammen med Naturen; men
Tankens Attribut griber saa igjen ud over Udstrækningen, Naturen
opfattes igjen gjennem en Tanke-%
% The hyphen falls at the page break, but it is the compound's own: the
% second element is capitalised, exactly as in `Tanke-Indhold' (p. 4).
% --- p. 9 ---
Abstraction, og Tanken hævder den Betydning, den havde som
Udgangspunkt. --- I Modsætning til Spinozas pantheistiske Absorberen
af Endeligheden fastholder \emph{Leibnitz}, indenfor samme
Grundretning, Tidsalderens Interesse: Individualiteten, og hævder ved
sin Monadelære det Individuelles Bestaaen, medens han mod Empirismen
hævder Idealismens Ret, det Aprioriskes Gyldighed, og hans Systems
theologiske Baggrund bevarer Uendelighedens Element og tilstræber
% DOUBTFUL READING, not a defect. The scan shows `Elemeut'; `Element' is
% set here. An earlier draft logged this as a turned n on the strength of the
% glyph at 600 dpi and of both OCR witnesses reading u. Neither argument can
% work on this file: it is a symbol-matching JBIG2 scan that demonstrably
% places one bitmap both where the text requires n and where it requires u
% (JBIG2-TEST.md, Result 2), and both OCR witnesses read the same
% reconstructed image, so their agreement is not independent. The burden of
% proof is on the claim of a defect (playbook §2 rule 3) and this scan cannot
% discharge it, so the sense-reading stands and the printing house is not
% accused.
%   NB: `Middelpunct' four lines below is one of the two words whose bitmap
%   the test caught being shared with a `u' (JBIG2-TEST.md, bitmap B).
Enheden. Ved \emph{Wolf} reduceredes hans Philosophie fra speculativ
Dybde til forstandig Klarhed, blev systematiseret og udtværet, og
derved egnet til at gaae over i den almindelige Bevidsthed i den Form,
i hvilken vi gjenfinde den i den \emph{tydske Populærphilosophie},
Forstandens Apotheose, hvor det endelige, aandløse Jeg, Parodien paa
Leibnitz's Monade, blev Universets Middelpunct, indirecte og ubevidst
gjort til et Absolut, medens det Absolute forsvandt i det Relative,
det Uendelige i Abstractioner, saa at Endeligheden udfylder
Philosophiens hele Gebet.

Paa dette Udviklingstrin mødes den idealistiske Retning ved Grændsen
af det Uphilosophiske med den franske Materialisme, der var den
consequente Udvikling af den Lockeske Philosophie.
% `Lockeske' here without apostrophe, `Locke'ske' six lines below: both as
% printed, verified on the image. Not normalised.

I \emph{Locke} var Philosophien vendt tilbage til Selvbevidstheden;
men ikke, som hos Descartes, for ud af Jeget at udvikle et Indhold,
der oprindelig laa deri, men forat paavise Jegets Passivitet i Forhold
til Objectet, forat tilnærme sig Enheden mellem Tanke og Væren ved at
reducere den første til den sidstes passive Afbillede. De
Inconsequentser, der endnu bleve tilbage i den Locke'ske Philosophie,
fjernedes idet den gik over til Frankrig. \emph{Condillac} drog
Sensualismens Consequents ud deraf, og hos de franske Materialister
forsvandt Aandens Begreb, og Forskjellen mellem Tanken og Tankens
materielle Gjenstand blev kun en Gradforskjel indenfor den mechanisk
bevægede Materie. Ved denne Udvikling vare Empirismens sidste
Conseqventser dragne; Modsætningerne vare reducerede til Enhed, men
% `Conseqventser' here, `Consequentser' on pp. 7--8: both as printed.
saaledes, at den ene var tilintetgjort. En tilsvarende Ud\-%
% --- p. 10 ---
vikling havde den samtidige Idealisme ikke naaet. Den
\emph{Berkeley}'ske Idealisme var kun formel; den beholdt under et
% PARTIAL LETTERSPACING: `B e r k e l e y' is spaced, the suffix `'ske'
% after the apostrophe is not (600 dpi). As printed.
andet Navn hele Realismens Indhold. Den \emph{skotske} Philosophie
hvilede paa en postuleret \emph{Erfaring}; den \emph{tydske} var kun
in abstracto Idealisme, og sigtede ubevidst mod den modsatte Retnings
Maal. \emph{Rousseau} endelig henviste vel paa et bestemt
Udgangspunkt for en kraftigere Idealisme; men uden Klarhed og uden
Gjennemførelse.

Ved det attende Aarhundredes Slutning var Alt forberedt til et
afgjørende Omsving i Philosophien. Materialismen havde draget sin
yderste Conseqvents, Idealismen sygnede hen, men var ikke udtømt; i
dens abstracte Forudsætning laa der et Incitament til en ny Form, og
Utilfredsstilletheden ved dens Indhold maatte hos den dertil Mægtige
kalde Bevidstheden om dette til Grund Liggende frem og fra det af
fremdrage en Udvikling. Den Udvei, der endnu stod Tanken aaben, var at
gjøre al virkelig Væren til \emph{tænkt Væren}, at sætte Tanken som al
Virkeligheds Kilde, og ud af Jeget, der er Eet med det Absolute og med
Naturens sande Væsen, at udvikle et System, der paa een Gang
fremtræder som Tankens og Virkelighedens. Det er denne Retning, den
tydske Philosophie tog ved Slutningen af forrige Aarhundrede, og som
culminerede i \emph{Hegel}. Det er, trods al Forskjel fra, ja
Modsætning til Christendommen, den sidste philosophiske Conseqvents af
dennes Opfattelse af Forholdet mellem det Aandelige og det Naturlige.
\emph{Ikke Berkeley, men Fichte og Hegel ere den abstracte
Materialismes Modsætninger og Paralleler.}

Hos \emph{Kant}, i hvis Lære om \emph{den transcendentale
Apperception} denne Retning er givet som i en Spire, fremtræder den
endnu begrændset og bunden ved sin Modsætning. Kant støtter sig ikke
paa en enkelt af de samtidige Retninger, men paa dem begge; han gaaer,
i Overensstemmelse med sin Tids hele Tendents, ud fra en Kritik af
Bevidstheden, men i Bevidstheden anerkjender han baade det active og
det passive Moment, og hævder saaledes Idealismens Ret medens han
anerkjender Empirismens Betyd\-%
% --- p. 11 ---
ning. Han sætter endnu ikke alt Objectivt ind i Selvbevidstheden,
lærer endnu ingen Tankens Monisme, tværtimod bliver netop ved ham
Modsætningen mellem Tænkning og Væren gjort til Philosophiens
Gjenstand; men Objectiviteten fortyndes for ham til en abstract
\emph{Ding an sich}, der hvert Øieblik vil falde sammen med Jegets
abstracteste Form. Denne Conseqvents drog \emph{Fichte} idet Jeget som
practisk sattes som producerende Ikke-Jeget, og \emph{Schelling}
fjernede den Tvetydighed, der laae i Bestemmelsen: Jeget, idet han
opfattede det som det absolute, med den under Selvbevidsthedens Former
anskuede Natur identiske Jeg, hvis Identitet i den ideale og reale
Potents-Række det blev Philosophiens Opgave at fatte. Methoden, der
skulde realisere denne Fordring, fremstillede \emph{Hegel}. For ham
bliver den rene logiske Tanke al Virkelighed, dens evige Bevægelse i
sig selv al Virkeligheds Kilde, Naturen dens Udtræden af sig selv,
Aanden dens Tilbagevenden til sig selv, Methoden selve Tankens indre
Bevægelse, derfor Eet med Philosophiens System og dette ligt med
Virkeligheden.

Med Hegel staaer en heel Udviklingsrække som afsluttet; hvad der er
sildigere end han, er kun Bestræbelser, der endnu ikke tilhøre
Historien. \emph{Hans System paa den ene Side, den franske
Materialisme paa den anden, repræsentere Endepunkterne af den dobbelt
udviklede Philosophie, hvis Forudsætning er Middelalderens Dualisme.}
% The whole sentence is letterspaced, four and a half lines of it; the
% first `Med Hegel' of the paragraph above is NOT.

Valget af de Systemer, der skulle repræsentere den nyere Tids
Philosophie, vil ved det Udviklede have fundet sin Begrundelse.
\emph{Bacon} og \emph{Système de la nature} betegne Begyndelsen og
Slutningen af en conseqvent fremskridende Række, der gaaer fra
Empirisme til Materialisme. \emph{Spinoza} repræsenterer den modsatte
Rækkes Begyndelse fordi han først conseqvent har udtalt hvad Descartes
% both `conseqvent' spellings on this page verified at 600 dpi (q + v);
% ABBYY reads the second as `consequent'.
kun antydede. \emph{Leibnitz}, som den, der philosophisk integrerer
Spinoza, og i sig indeholder den frugtbare Spire til en følgende
Udvikling, repræsenterer Afslutningen af Idealismens
% --- p. 12 ---
første Hovedstadium (som Dogmatisme). \emph{Kant} og \emph{Hegel}
omslutte den nyeste Philosophie.

For Overgangsperioden har Valget været vanskeligere, og Isolertheden i
dens Systemer, det Mangfoldige og Usammenhængende i deres Tendentser,
vil altid gjøre et enkelt Systems Overvægt betinget. Men \emph{det
platoniske Akademie i Florents} repræsenterer som Collectivindividuum
% here `det' IS inside the spaced run, unlike the parallel phrase on p. 2.
den første begeistrede Tilslutning til den gjenfundne Oldtid, medens
\emph{Giordano Bruno} kraftigst udvikler og fyldigst sammenfatter sin
Tids Tankeindhold, i en Form, der bærer Præget af en ved Tankens
Lidenskab bevæget Aand, i hvilken Poesien og Philosophien, hin Tids
store Magter, kæmpe om Herredømmet. ---

Ved Fremstillingen af de enkelte Systemer vil min Stilling væsentlig
være historisk characteriserende og kun forsaavidt criticerende, som
der under den reproducerende Udvikling vise sig Inconseqventser i
Gjennemførelsen, navnlig ved mislykkede Forsøg paa at optage hvad der
er udelukket ved Principernes Ensidighed eller Mangelfuldhed. Det
bliver saaledes væsentlig \emph{Grundtankens Udførelse}, ikke
Grundtanken selv, der bliver Gjenstand for Kritiken. Selve
Grundtankens Kritik er for alle de tidligere Systemer givet i den
historiske Udvikling; Opgaven ved det enkelte System bliver kun at
bringe den frem i dens hele Eiendommelighed og Bestemthed, og derved
at paavise dens historiske Berettigelse som et Led i Tankens
Udvikling. En uberettiget Kritik vilde den være, der traadte hen til
Systemerne med Fordring paa Løsningen af bestemte Opgaver, som
Systemerne ikke have stillet sig selv, Kunde det lykkes et System at
% PRINTER'S DEFECT, as printed: a comma where the sense requires a full
% stop -- `stillet sig selv, Kunde det lykkes'. Verified at 600 dpi: a
% comma, followed by the wide space of a sentence break, then a capital K.
% Not corrected.
afslutte sig til en virkelig Heelhed, uden at enkelte Problemer
traadte frem, da vilde dette være et Tegn paa, at Problemerne vare
falsk stillede, at Opgaven vilde være, ikke at løse dem, men at afvise
dem.

For fuldstændig at bringe Grundtanken frem, vil det være nødvendigt at
medtage alt det af den detaillerede Udførelse, der giver et Indblik i
Systemets indre Motiver, eller kaster et bestemtere Lys tilbage paa
Principet. Ogsaa hvad
% --- p. 13 ---
\setcounter{footnote}{0}%
der kun ved løsere Traade hænger sammen med dette, men har faaet
eiendommelig historisk Betydning, tør ikke savnes i Fremstillingen.

Jeg har begyndt Rækken af de monographiske Fremstillinger med Spinoza,
dels af den ydre Grund, at Materialet her laae færdigt i størst
Fuldstændighed, dels i den indre, der er givet i hans Systems
Betydelighed og i den store Indflydelse, det har faaet paa
Philosophiens Udvikling i dette Aarhundrede. Jeg har saameget mindre
taget i Betænkning at begynde midt inde i Rækken som hver enkelt
Monographie udgjør et Hele for sig, der ogsaa uden de andre vil have
sin Fuldstændighed.

Under Fremstillingen af Spinoza har jeg sjelden omtalt den tidligere
Literatur. De, som ere fortrolige med denne, ville snart see, at der
er taget Hensyn til alt blot nogenlunde Betydeligt i den, ogsaa til
hvad der tilhører den seneste Tid\footnote{Foruden de bekjendte Værker
om Philosophiens Historie, dels i dens hele Omfang, dels i kortere
Perioder, af \emph{Brucker}, \emph{Tennemann}, \emph{Tiedemann},
\emph{Buhle}, \emph{Hegel}, \emph{Schleiermacher}, \emph{Ritter},
\emph{Feuerbach}, \emph{Erdmann}, K.\ \emph{Fischer} (Gesch.\ der
% `K.' is NOT letterspaced; `Fischer' is (600 dpi). As printed.
neueren Philosophie Bd.\ 1.), \emph{Cousin} (fragments de philosophie
% `Bd. 1.': a digit one, not a capital I -- the glyph carries the flag of
% the `1' in `1854' two lines below, not the symmetrical serifs of an I.
Cartésienne), og \emph{Bouillier} (histoire de la philos.\ Cartésienne.
Paris 1854.\ 2.\ Vol.), samt Monographierne af \emph{Sigwart} (der
Spinozismus historisch und philosophisch erläutert mit Beziehung auf
ältere und neuere Ansichten), \emph{Schaarschmidt} (Descartes und
Spinoza), \emph{Helfferich} (Spinoza und Leibnitz),
\emph{Montbeillard} (de l'Ethique de Spinoza), \emph{Erdmann}
(Vermischte Aufsätze), og \emph{Trendelenburg} (Ueber Spinozas
Grundgedanken und dessen Erfolg), har jeg benyttet hvad der findes af
critiske og polemiske Udviklinger af Spinozismen hos \emph{Leibnitz}
(navnlig i den af \emph{Foucher de Careil} udgivne: réfutation inédite
de Spinoza par Leibnitz), \emph{Wolf}, \emph{Jacobi}, \emph{Herder},
\emph{Mendelssohn}, \emph{Maimon}, \emph{Schelling}, \emph{Hegel},
\emph{Baader}, \emph{Herbart} og \emph{Trendelenburg}. Spinozas
samtidige theologiske og philosophiske Modstanderes Skrifter have kun
tildels været mig tilgængelige. De give kun Udbytte gjennem deres
Misforstaaelser.}. Jeg har ment, at den rette og for den historiske
Fremstillings Enhed tjenligste Maade at benytte den paa vilde være,
efterat have brugt den som Hjælpemiddel til at
% --- p. 14 ---
undgaae Misforstaaelser og en ufuldstændig Opfattelse, da at lade
selve den sammenhængende Reproduction af Spinozas Tanker ved sin Enhed
godtgjøre sig som dèn rigtige og med det Samme afvise de falske
% `dèn': the accent slopes down to the right at 600 dpi, i.e. a grave, not
% the acute of `Cartésienne' / `inédite' on p. 13. ABBYY reads plain `den'.
Opfattelser paa en indirecte Maade. Lykkes det fra et Systems levende
Midtpunkt at gjengive det i dets Heelhed, da vil selve denne organiske
Heelhed godtgjøre Opfattelsens Rigtighed i det Hele og det Enkelte, og
af sig selv fjerne den falske eller ufuldstændige.

Med dette Forord, hvis Udførlighed maa finde sin Undskyldning deri, at
Meget har maattet søge sin Plads i det, der i en sammenhængende
% `Meget', not `Méget': the mark above the e is a speck, not an accent
% (600 dpi); ABBYY reads `Meget'.
Philosophiens Historie vilde have fundet en Plads i selve Bogen,
overgiver jeg da dette Skrift til det ikke talrige Publicum, som Bøger
af dette Indhold kunne gjøre Regning paa at finde i Danmark. De øvrige
Monographier ville blive udgivne saa hurtigt som Forholdene tillade
mig det. Nærmest efter denne ville følge de to Afhandlinger om
\emph{Frants Bacon af Verulam} og \emph{Système de la nature}, der,
paa Grund af deres Forhold til hinanden, ville udkomme samtidig.

% The date line is set in smaller type, indented from the left margin,
% with `K j ø b e n h a v n' letterspaced; the signature is bold, smaller,
% and ranged short of the right margin. A short centred rule closes the
% Fortale at the foot of the page.
\bigskip
\begin{flushleft}
\hspace{2em}{\small\emph{Kjøbenhavn}, den 6te April 1856.}
\end{flushleft}

\begin{flushright}
{\small\textbf{H.\ Brøchner.}}\hspace{4em}
\end{flushright}

\bigskip
\begin{center}\rule{2.5cm}{0.4pt}\end{center}

\newpage

% =====================================================================
%  DIVISIONAL HALF-TITLE  (printed p. 15 = PDF 25), verso p. 16 blank.
%  `B e n e d i c t  S p i n o z a .' in large letterspaced bold, sunk about
%  a third down the page, over three Ethics epigraphs in small roman,
%  indented to about the middle of the measure. Æ is a real ligature sort.
%  The letterspacing INSIDE the epigraphs is the print's own: `D e o',
%  `n e c e s s i t a t e', `D e i  c o g n i t i o'. The em dash before
%  each citation is the body's em dash.
% =====================================================================
% --- p. 15 ---
\chapter*{Benedict Spinoza}
\addcontentsline{toc}{chapter}{Benedict Spinoza}
\markboth{Benedict Spinoza}{Benedict Spinoza}

\begin{center}
{\Large\textbf{\emph{Benedict Spinoza.}}}
\end{center}

\vspace{2\baselineskip}

\begingroup
\setstretch{1.15}
\begin{quote}
\small
Quicquid est, in \emph{Deo} est; et nihil sine Deo esse neque concipi
potest. --- Eth.\ I.\ prop.\ 15.

\medskip

Æternum illud et infinitum ens, quod Deum seu naturam appellamus, eadem,
qua existit, \emph{necessitate} agit. --- Eth.\ IV.\ præf.

\medskip

Summum mentis bonum est \emph{Dei cognitio}, et summa mentis virtus Deum
cognoscere. --- Eth.\ IV.\ prop.\ 28.
\end{quote}
\endgroup

% --- p. 16 ---
% blank verso of the divisional half-title.

\newpage

% =====================================================================
%  INDLEDNING  (printed pp. 17--26 = PDF 27--36)
%  Printed head: `Indledning.' letterspaced bold, then the subtitle
%  `Spinozismens historiske Genesis.' -- which the Indhold omits (header
%  §Indhold (a)). Two-line raised initial N. Opens at the head of p. 17.
% =====================================================================
\chapter*{Indledning}
\addcontentsline{toc}{chapter}{Indledning}
\markboth{Indledning}{Indledning}

\begin{center}
\textbf{Spinozismens historiske Genesis.}
\end{center}

% --- p. 17 ---
% The head (`Indledning.' letterspaced bold, the subtitle `Spinozismens
% historiske Genesis.' and the short centred rule below it) is in the
% skeleton. The first word carries a two-line raised initial N, transcribed
% here as an ordinary letter. No footnote on this page; the `2' at the foot
% is the signature mark of the third gathering and is not transcribed.
% No letterspacing anywhere on p. 17 -- `Spinozas', `Fichte', `Kant' and
% `Descartes' are all set normally (600 dpi).
Naar det, forat forstaae et philosophisk System i dets historiske
Betydning i videste Forstand, forat opfatte de bevægende Kræfter, der,
hyppigt kun anede af Systemets Ophavsmand, virke igjennem det, forat see
dets Berettigelse og Grændsen for denne, er nødvendigt at kaste et Blik
ud over den menneskelige Tankes hele Udviklingsgang; saa er det, forat
opfatte et System i dets bestemte Form, forat forfølge dets Grundtanke i
dens Consequentser, og forat opdage dets historiske og psychologiske
Genesis, tilstrækkeligt at gaae tilbage til det nærmeste store Afsnit i
Historien, hvor den nye Begyndelse danner en Kløft mellem sig og Fortiden
og optræder med Eiendommelighedens berettigede Fordring paa
Selvstændighed. Vi finde ogsaa, at de Systemer, der følge nærmest efter
saadanne Epocher, i disse see Grændserne for deres historiske
Sammenhæng, og, hyppigt med miskjendende Ringeagt for Fortiden, kun tage
dem som Udgangspunkt for en conseqvent Tankeudvikling. Saaledes gaaer
% `conseqvent' as printed (600 dpi), with qv; the book sets `Consequents',
% `Consequentser' and `consequent' with qu elsewhere. Not corrected.
Spinozas historiske Parallel: Fichte, i sin første Optræden kun tilbage
til Kant, Spinoza selv kun tilbage til Descartes.

Med Descartes er (jfr. Fortalen, S. 8) den ene af den nyere Philosophies
Udviklingsrækker begyndt med Fordringen paa, ud fra Selvbevidstheden, fra
det tænkende Jeg, efter en fast Methode at udvikle Philosophiens Indhold.
Antyd\-%
% --- p. 18 ---
\setcounter{footnote}{0}%
ninger til denne Retning vare allerede givne hos tidligere Tænkere. De
\emph{franske Skeptikere} havde viist tilbage til Selvbevidstheden,
% `De' before the run is NOT letterspaced (600 dpi); the run begins at
% `franske'. Same on the next line: only `Campanella' is spaced.
\emph{Campanella} havde atter erindret om Augustinus's Slutning fra
Jegets Tænken til dets Væren; men disse Antydninger havde staaet som
løsrevne, de vare ikke forfulgte i deres Consequentser. Det blive de
først idet Descartes, afkastende Traditionen, og polemisk vendende sig,
ikke saameget mod den allerede tilbagetrængte scholastiske Philosophie,
som mod Oldtidens gjenoplivede Tænkning, gjennem Tvivlen om Alt fører
tilbage til Jeget som det sikkre Tilknytningspunkt, og, idet Tvivlen
bliver en Udskillen af Alt, hvad der ikke hører til Jegets Væsen,
bestemmer dette som Tænkning, og nu i en umiddelbar Anskuelse (simplici
mentis intuitu), ikke gjennem en Slutning --- thi cogito ergo sum er ikke
nogen Slutning --- lader det tænkende Jeg bekræfte sig som værende (sum
cogitans). I Udgangspunktet er saaledes givet Tankens umiddelbare Vished
om sig selv, dens umiddelbare Selvbekræftelse, dens active Adskillelse
fra og Modsætning til det, hvis Væsen ikke bestemmes som Tanke, og som
consequent maatte have været sat som et Ikke-Værende. Men i denne
Adskillelse og Modsætning taber Tanken sit
Indhold\footnote{Denne Consequents gjælder for Descartes, ikke for
Spinoza. Hos hin er Jeget sondret fra Naturen og denne udelukket af Guds
Væsen; hos Spinoza er Jeget som tænkende i Gud i Enhed med Naturens
Væsen.}, og medens Fordringen stilles: at udvikle Erkjendelsen af sig,
bliver der ingen Vei aaben til, ud fra sig selv alene atter at naae
dertil. \emph{Guds}begrebet maa her bane Veien. Til det naaer Descartes,
% PARTIAL LETTERSPACING, verified at 600 dpi: `G u d s' is spaced,
% `begrebet' on the same line is not. Transcribed as printed, not
% regularised (cf. `n a t u r a naturans', p. 71).
% `Descartes' here carries a round speck over the second e -- both OCR
% witnesses read `Descartés'. It is the same size and shape as the specks
% scattered over this leaf and the i-dots beside it; the sense-reading is
% transcribed and no defect is logged (JBIG2 scan, RECON §9).
idet Tvivlen, der, dybere opfattet, var Tankens Fastholden af sig selv
under Udsondringen af Alt, hvad der var den fremmed, opfattes som
Endelighedens Tegn, og det tvivlende, endelige Jeg nu i sig finder Ideen
om det uendelige Væsen, en Idee, der garanterer dette Væsens Tilværelse,
da den kun kan have sin Oprindelse fra et saadant Væsen, ikke fra det
endelige Jeg, idet ellers Virkningen
% --- p. 19 ---
\setcounter{footnote}{0}%
vilde indeholde en større Realitet end Aarsagen. I Forestillingen om Gud
som fuldkommen findes nu Forestillingen om hans Sanddruhed; den giver os
Tillid til vore klare og tydelige Ideers
Sandhed\footnote{Consequent opfattet er selve den klare og tydelige Idee
Sandhedens Kriterium, idet Tanken i den er i Enhed med sig selv og
bekræfter sig selv. Cogito ergo sum er det tænkende Jegs
Selvbekræftelses første Form, Ideen om Gud som Tankens absolute Væsen
dens høieste Form, og i Begrebet dens første. Men idet Tvivlen opfattes
som Endelighed, bliver den klare og tydelige Idee usikker, og maa støttes
ved Gudsforestillingen.}; thi det uendelig fuldkomne Væsen kan ikke ville
bedrage, dette vilde forudsætte en Ufuldkommenhed. Men til de klare og
tydelige Ideer regner Descartes, ved at give Begrebet Tænkning et
Omfang, der staaer i Modsætning til hans dybere Opfattelse deraf (idet
ikke blot intelligere og velle, men ogsaa \emph{imaginari} og
\emph{sentire} regnes dertil), ogsaa Ideen om et Legeme, der er forenet
% `intelligere og velle' is NOT spaced; only the second pair of Latin
% infinitives is (600 dpi).
med Aanden, og gjennem dette Legeme aabnes Forbindelsen med Legemernes
Verden, med Naturen, der, trods Modsætningen til Tankens Væsen,
bestandig udøver en magisk Magt over Descartes's Tanke, og bliver dennes
væsentlige og kjæreste Gjenstand, medens han ligesom skyer Aandens
concretere Gebeter, navnlig det Ethiske.

I Modsætningen til Aanden er Naturens Uafhængighed af denne givet; idet
Aanden opfattes som Substants, maa Legemet ogsaa opfattes som saadan; den
tænkende Substants faaer ligeoverfor sig en \emph{udstrakt} Substants;
thi for Naturen, opfattet \emph{gjennem} Tanken, abstract som dennes
% PARTIAL LETTERSPACING, both verified at 600 dpi: `u d s t r a k t' is
% spaced and `Substants' beside it is not; `g j e n n e m' is spaced and
% `Tanken' is not.
Modsætning, bliver Quantiteten, Udstrækningen, den væsentlige
Bestemmelse. --- Adskilt fra disse endelige Substantser kommer den
\emph{uendelige Substants, Gud}, til at indtage sin Plads. Idet
% the comma inside this run is itself letterspaced from `Substants' and
% from `Gud'; `den' before the run is not spaced (600 dpi).
\emph{Tanken} er Udgangspunktet, bestemmes Gud saaledes, at Tanken
gjenfinder sig i ham, Udstrækningen udelukkes fra hans Væsen. Men idet
\emph{Jeget} er Udgangspunktet, og fæstner sig fra reen Activitet til et
tænkende \emph{Væsen}, forsvinder det ikke i Gud, skjøndt hans Væsen
% PARTIAL LETTERSPACING: `tænkende' is NOT spaced, `V æ s e n' at the head
% of the next line is (600 dpi).
% --- p. 20 ---
\setcounter{footnote}{0}%
bliver Tankens absolute Væsen; Consequentserne af dette Begreb drages
ikke, de antydes blot\footnote{Idet der tillægges Tanken om det Uendelige
en Prioritet for Tanken om det Endelige (jfr. Ep. I, 119), og den sættes
som oprindeligt værende i os (pr. phil. I, 20-22 o. fl. St.).}, og
% `20-22' is a hyphen, not the en rule the Indhold uses for page ranges
% (measured on the 600 dpi render against the hyphen of `Aimé-Martin' in
% note 2, which is set in the same size).
Substantsbegrebet holdes i en Tvetydighed, der gjør det muligt, at bevare
en tvivlsom Selvstændighed for de endelige Jeg'er. Vel bestemmer
Descartes Substantsen som: une chose, qui existe en telle façon, qu'elle
n'a besoin que de soi-même pour exister (principes de la philos. I,
51)\footnote{Jeg citerer efter: \emph{Oeuvres philosophiques de
Descartes} publiées d'après les textes originaux par
L. \emph{Aimé-Martin}. Paris 1839.} og ved denne Definition begrændses
% in note 2 the initial `L.' is NOT spaced, exactly as `K.' before
% `F i s c h e r' on p. 13; the run resumes at `A i m é'.
den til Gud; men den udvides strax igjen til dem af de skabte Ting, der
ikke behøve \emph{andre skabte} Ting forat existere, men dertil kun
% PARTIAL LETTERSPACING: `a n d r e  s k a b t e' ends the line spaced and
% `Ting' at the head of the next line is not; the first `skabte Ting' in
% the same sentence is not spaced at all (600 dpi).
behøve Guds concours ordinaire, medens det øvrige Endelige kaldes disse
Substantsers \emph{Qualiteter} eller \emph{Attributer}. Fra den
»\emph{uskabte Substants}, der tænker og er uafhængig« (I, 54), adskilles
% no space inside the guillemets, as printed; the spaced run stops at
% `Substants' and `der tænker og er uafhængig' is set normally.
da de to Classer af \emph{skabte Substantser}, den \emph{legemlige} og
den \emph{aandelige}, der hver have sit særegne Attribut, som
% `den' and `og den' between the runs are NOT spaced (600 dpi). The
% compositor broke `legem-lige' at the line end inside the spaced run.
constituerer deres Natur og Væsen, og af hvilket deres øvrige Attributer
ere afhængige. Saaledes constituerer \emph{Udstrækningen} i Længde,
Bredde og Høide den udstrakte Substantses Væsen, \emph{Tænkningen} den
tænkende Substantses (I, 53). De to Classer af endelige Substantser have
Intet tilfælles, men Bevidstheden om deres Forening i Mennesket, der
opfattes som en \emph{Sammensætning} af begge, tvinger Descartes til
Hypothesen om et Berøringspunct i glandula pinealis (conarion), Sjælens
umiddelbare Sæde i Hjernen, gjennem hvilket Legemets Bevægelser meddele
sig som Fornemmelser til Sjælen, Sjælens Tanker som Bevægelser til
Legemet\footnote{Denne Hypothese fastholder Descartes dog ikke consequent,
da der allerede hos ham findes Antydninger til den af hans Skole udviklede
Theorie om Causæ occasionales, ved hvilken den guddommelige Causalitet
bliver det baade i Aand og Legeme Virkende, der mægler mellem begges
Forandringer.}; men Muligheden af denne For\-%
% --- p. 21 ---
vandling under Meddelelsen paaviser han ikke, og den kan ikke paavises ud
fra hans Principer.

Istedenfor med en sammenhængende og i sig selv begrundet Erkjendelsens
Enhed ender derfor Descartes med en dobbelt, uforsonet Modsætning:
Modsætningen mellem den tænkende og den udstrakte Substants, og mellem de
endelige Substantser og den uendelige. For de Tænkere, der sluttede sig
til ham, var der altsaa stillet de to for hans hele Tænkning afgjørende
Problemer, at finde Enheden mellem Aand og Legeme, og Enheden mellem Gud
og det Endelige, og derved at forsone den dobbelte Splid.

Consequentserne af Descartes's Opfattelse af de to endelige Substantsers
Væsen droges af \emph{de la Forge}, \emph{Geu\-%
lincx} og \emph{Malebranche}, idet Væsensforskjelligheden fastholdtes, og
% `og' between the spaced names is NOT letterspaced, here and in the
% parallel list below (600 dpi) -- the same usage as in the Fortale.
Foreningen søgtes, ikke, gjennem en phantastisk Hypothese, i Mennesket,
men gjennem et Postulat i Gud, der i begge de to Rækker af Substantser
frembringer tilsvarende Forandringer, medens de endelige Substantser kun
ere Anledninger (causæ occasionales) til hinandens Forandringer. Og baade
\emph{Geulincx}, \emph{Malebranche} og \emph{Clauberg} førte
Consequentserne af det af Descartes opstillede Forhold mellem den
uendelige og de endelige tænkende Substantser saavidt, at kun den
religiøse Forestillings Forudsætninger forhindrede den endelige tænkende
Aands Sammensmeltning med den uendelige. Det var ogsaa kun disse
Forudsætninger, der afholdt \emph{Malebranche} fra, med Bestemthed at
sætte den \emph{reale} Udstrækning som Guds Attribut.

Kun idet hine Forestillinger forlodes, kunde den strænge Consequents af
det Descartes'ske Substantsbegreb drages, og det andet Problem føres
videre; men idet Consequentsen droges med Hensyn til dette, indvirker den
forandrede Opfattelse umiddelbart paa Opfattelsen af det første Problem
og en ny Formel opstilles for dettes Løsning. Denne Udvikling finder Sted
ved \emph{Spinoza}, og idet han har tænkt til Ende hvad hans Forgænger
havde ladet staae tvetydigt, bliver han den egentlige Repræsentant for
den ved hin angivne Retning.
% --- p. 22 ---
\setcounter{footnote}{0}%

For Descartes havde \emph{Tanken i Jegets Form} været Centrum, fra Jeget
havde han taget sit Udgangspunkt, og selv ligeoverfor det Uendelige, der
opfattedes under samme Grundbegreb, havde han fastholdt det, og afvist de
Problemer, der truede dets Existents\footnote{Saaledes navnlig det ved
Descartes's Opfattelse af Opholdelsen skærpede Problem om den
guddommelige Almagts Forhold til de endelige Substantsers
Selvvirksomhed, specielt til den menneskelige Frihed.}. Men idet Jeget
bestemtes som \emph{tænkende}, med Udelukkelse af Individualitetens
Betingelser, bestemtes dets Væsen kun ved Tankens \emph{Almindelighed},
og den Consequents maatte drages, der lader den endelige Forms
selvstændige Bestaaen forsvinde, og sætter Væsenet, Tanken i dens
Uendelighed, som det ene Værende. Det Uendeliges Tanke taaler ikke
Begrændsning ved et for sig bestaaende Endeligt, og det saameget mindre
som den ikke er en secundær, men den oprindelige Tanke; thi det
Begrændsede --- dette havde allerede Descartes udtalt --- begribes kun
ved det Ubegrændsede; dette er det Positive, det til Grund Liggende,
Begrændsningen er en Bestemmelse indenfor det og forudsætter
det\footnote{Hvorledes Spinozismens eiendommelige Bestemmelse af den
uendelige Substantses Væsen fører til at bestemme den almindelige
Sætning, at det Endelige ikke kan existere \emph{ved Siden} af det
Uendelige, saaledes, at det, der gjør det til Endeligt, kun er et
Negativt (at Begrændsning er Negation), det Endelige altsaa ikke blot er
i det Uendelige, men er i det \emph{kun som Modification}, vil senere
blive paavist (jfr. Cap. 3.).}. Den tænkende Substantses \emph{første
Tanke} er Uendelighedens Tanke, Tanken om Gud; men i denne Tanke er den
ikke skilt fra det tænkende Uendelige, den er i Enhed med det: en
selvstændig Stilling for den endelige tænkende Substants bliver en
Umulighed. Hvad Descartes benævnede endelige tænkende Substantser, maa
derfor forandre Navn og blive til \emph{Modi\-%
ficationer af den ene uendelige Substants}, de maae kun have deres
% the whole phrase, `af den ene' included, is spaced through the line
% break; `de maae' after the comma is not (600 dpi).
Realitet i denne. Og det Samme gjælder om de \emph{udstrakte}
Substantser. De sættes af Descartes
% --- p. 23 ---
som skabte af Gud, existerende, ligesom de tænkende, alene ved hans
fortsatte Opholdelse, men Udstrækningens Væsen sættes ikke i Gud, han
bestemmes kun som tænkende, skjøndt denne Bestemmelse, hos Descartes, kun
har Betydning som Modsætning til Udstrækningen. Men Aarsagen kan kun være
Aarsag til hvad der er begrundet i dens Natur, den kan kun meddele hvad
den selv har i sig; sættes derfor Gud som Aarsag til de udstrakte
Substantser, maa Udstrækningen som uendelig sættes i Gud, og ligesom de
tænkende Substantser kun havde deres Realitet som Modificationer i Gud,
maa det Samme gjælde om de udstrakte. Først derved hæves ogsaa den
Modsigelse hos Descartes, at den tænkende og den udstrakte Substants
begribes alene gjennem det Attribut, der constituerer deres Væsen, uden
at der i deres Begreb ligger en Relation til Gud, af hvem de altsaa med
Hensyn til Væsenet blive uafhængige, medens de med Hensyn til Existentsen
ere afhængige af ham. Denne Modsigelse mellem Afhængighed og
Uafhængighed hæves kun idet det, som Descartes opfattede som
selvstændige Substantser, opfattes som Attributer ved den ene uendelige
Substants, der er udelt i begge, saa at begge udtrykke dens Væsen helt,
og medens de ere uafhængige af hinanden idet de hver for sig udtrykke
Substantsens \emph{hele} Væsen, tillige ere i Enhed ved at udtrykke den
\emph{ene} Substants (jfr. Cap. 2.)
% AS PRINTED: the paragraph ends `(jfr. Cap. 2.)' with no full stop after
% the closing parenthesis (600 dpi, both witnesses). The parallel
% reference in note 2 on p. 22 prints `(jfr. Cap. 3.).'

Naar nu, som Fremstillingen af Spinozas Philosophie vil vise, denne helt
er concentreret i hans Substantsbegreb, og dette umiddelbart udvikler sig
som Consequents af det hos Descartes Givne, saa ligger heri, at det ikke
er nødvendigt at forfølge Spinozismens Oprindelse længere tilbage.

Man har, skuffet med en overfladisk Lighed, og misforstaaende
Spinozismens Eiendommelighed, villet føre denne tilbage til
\emph{Kabbalah's} Emanationslære. Udviklingen i de følgende Capitler vil
% the possessive `'s' is inside the spaced run; `Emanationslære' is not
% spaced (600 dpi).
vise, hvor uforenelig Spinozas Opfattelse af den evige Væren, af det
Endelige som kun værende i det Uendelige, og, efter sin Sandhed, med
ophævet Begrændsning, er med Emanationslærens Opfattelse af den endelige
Tilværelse som opstaaet ved en ufri Udfoldelse i
% --- p. 24 ---
\setcounter{footnote}{0}%
\emph{Tiden}, som selvstændig og som frafalden. Den lidet anerkjendende
Maade, paa hvilken Spinoza udtaler sig om den jødiske Tradition (jfr.
tract. theol.-pol. Cap. 7 og 9), og de faa og ubestemte Henvisninger til
jødiske Tænkeres Anskuelser (Ep. 21, 29. Eth. II, 7, schol.) ere ikke
heller skikkede til at bevise et nærmere Afhængighedsforhold.

Men idet et saadant Afhængighedsforhold med Hensyn til det Væsentlige,
til Systemets Grundtanke, afvises, afvises derved ikke, hverken en
Indflydelse af andre Philosopher paa Enkeltheder i Gjennemførelsen, eller
en almindelig Paavirkning af bestemte Retninger i Tidsalderen og af
Spinozas jødiske Nationalitet. Paa Spinozas Lære om Staten har
\emph{Hobbes} indvirket, maaskee ogsaa \emph{Hugo Grotius}, skjøndt
Spinoza ikke nævner ham; for Bestemmelsen af Religionens Forhold til
Philosophien kan \emph{Herbert af Cherbury} have havt Betydning, for den
% `af' inside this run IS letterspaced (600 dpi), unlike the `og' between
% the spaced names on p. 21.
skarpere Opfattelse af de ved Descartes forberedte Problemer
\emph{Arnold Geulincx}\footnote{Navnlig med Hensyn til Opfattelsen af de
enkelte Aander som modi i Gud, og af Udstrækningen som res singularis og
uendelig.}. Paa den pantheistiske Grundanskuelse hos Spinoza kunne de
pantheistiske Elementer, der fra Overgangsperioden vare blevne tilbage i
Tidsaanden, have havt Indflydelse; paa den strænge Opfattelse af
Substantsen, Gud, som det ene Værende, den jødiske Gudsforestilling. Men
enhver ydre Paavirkning, gjennem bestemte Theorier eller gjennem
Tidsalderens hele aandelige Atmosphære, er hos Spinoza forarbeidet til
indre Consequents; hans System danner en Helhed, der forstaaes ud fra sig
selv, fra eet Princip.

Ville vi, forud for den egentlige Fremstilling af Spinozas Philosophie,
til foreløbig Orientering angive dette Princip og dets nærmeste
Consequentser for hans Tænknings Form og hans Methode, saa ligge
Momenterne dertil i det, der ovenfor blev udviklet angaaende hans Forhold
til Descartes. Idet Substantsbegrebet opfattes saaledes, at det i sin
Enhed udfylder hele Tilværelsen, behersker det Tanken som dens ikke blot
væsentlige, men, strængt taget, eneste
% --- p. 25 ---
\setcounter{footnote}{0}%
Gjenstand. I Substantsen finder Tanken al Virkelighed, i den finder den
sig selv. Men det, at al Virkelighed er i Substantsen, begrunder den
\emph{pantheistiske} Opfattelse. Det, at Tanken umiddelbart er i Enhed
med Substantsen, der er dens eget uendelige, objectiverede Væsen, og
tillige Udstrækningens Væsen, giver Tanken dens umiddelbare Sikkerhed og
Systemet \emph{Dogmatismens} Characteer. Det, at al \emph{Virkelighed}
indesluttes i Substantsens Begreb, medfører, at dette Begreb for Tanken
% `at al' before the run is NOT spaced (600 dpi); nor is the parallel
% `at al Virkelighed er i Substantsen' three sentences earlier.
maa udfolde sig, og Characteren af denne Udfoldelse, der er i Evighedens,
ikke i Tidens Element, maa bestemmes ved Naturlovens Form,
Nødvendigheden. Thi idet Tanken først har sondret Naturen ud fra sig som
sin Modsætning, og bestemt den ved Udstrækningens abstracte Begreb,
virker denne Opfattelse tilbage paa den selv, der maa søge sit Indhold i
Naturen, og idet Tanke og Udstrækning forenes i Substantsens Enhed,
bliver den mechaniske Nødvendighed, den abstract opfattede Natnrs Lov,
% AS PRINTED: `Natnrs' for `Naturs'. At 600 dpi the fourth sort is plainly
% an n -- arched and closed at the top, with two footed stems -- unlike the
% u's beside it, and both OCR witnesses read it so. Recorded here, not
% corrected; no claim is made about the compositor's case (JBIG2 scan).
Tankens og Substantsens Lov. Og om end i Substantsen som det ene
Existerende Nødvendigheden ikke er en ydre Nødvendighed, ikke Tvang, men
en indre, ved dens eget Væsen bestemt, Nødvendighed, en libera
necessitas, fortrænger \emph{Causalitetens} Begreb dog Øiemedets, og det
baade paa Grund af Naturbegrebets Indvirkning paa Substantsbegrebet, og
paa Grund af den strænge Fastholden af Uendelighedsbegrebet i Modsætning
til anthropomorphistiske Forestillinger.

Systemets \emph{Methode} bestemmes umiddelbart ved Grundbegrebet. Tanken
er ikke adskilt fra sin Gjenstand; denne er Virkeligheden i hin, Tankens
egen Realitet; Gjenstandens Udfoldelse maa derfor give Tænkningens Form
dens Præg; dens af Causalitetsbegrebet betingede Nødvendighed bestemmer
% `Præg': the P has lost the lower half of its bowl in the impression and
% reads as I/F at 600 dpi (tesseract `Iræg', ABBYY `Fræg'). The
% sense-reading is transcribed; no defect is logged (JBIG2 scan).
\emph{Tankeudviklingens mathematiske} Form\footnote{Ikke just den
synthetiske Methode, som Spinoza kun har anvendt i Ethiken og i Principia
phil. Cartes., og hvis Anvendelse istedenfor den analytiske han navnlig
tillægger et didactisk Øiemed, men i Almindelighed den med mathematisk
Stringents bevisende, der til det ved sig selv Klare og Enkelte med
mathematisk Nødvendighed knytter Consequentsernes fast sammenkjædede
Række.},
% this note begins on p. 25 and its last two lines stand in the note block
% at the head of p. 26, below a rule and with no repeated mark. The whole
% note is attached here to its mark, so p. 26 carries no \setcounter.
% `Form' after the spaced run is not itself spaced (600 dpi).
% --- p. 26 ---
medens Tankens Udgangs- og Slutningspunkt bliver \emph{den umiddelbare
Anskuelse}, svarende til Substantsens Væren i sig selv.

En \emph{Inddeling} af Philosophien ud fra sit Princip har Spinoza ikke
givet, men kun antydet. Det vil i det følgende Capitel blive fremhævet,
hvilke eiendommelige Vanskeligheder der maatte opstaae for ham paa dette
Punkt, og disse Vanskeligheder ville blive endnu tydeligere ved den
nærmere Udvikling af hans Lære.

% A short centred rule, about one fifth of the measure, stands well below
% the last line of text and closes the Indledning; the rest of p. 26 is
% blank. Cap. 1 opens at the head of p. 27.
\bigskip
\begin{center}\rule{2.5cm}{0.4pt}\end{center}

\newpage

% =====================================================================
%  FØRSTE CAPITEL  (printed pp. 27--43 = PDF 37--53)
%  Opens at the head of p. 27, two-line raised initial E.
% =====================================================================
\chapter*{Første Capitel.}
\addcontentsline{toc}{chapter}{Første Capitel. Spinozas Liv og Skrifter}
\markboth{Spinozas Liv og Skrifter}{Spinozas Liv og Skrifter}

\begin{center}
\textbf{Spinozas Liv og Skrifter.}
\end{center}

% --- p. 27 ---
\setcounter{footnote}{0}%
% The head (`Første Capitel.' bold, `Spinozas Liv og Skrifter.' bold, and the
% short centred rule below) is in the skeleton. The first word carries a
% two-line raised initial E, transcribed here as an ordinary letter.
% `C o l e r u s' is the only letterspaced run on the page (600 dpi);
% spacing.py recovered it as `olerus'.
En Tænkers Livshistorie er hans Tankes Historie, thi hans Livs Indhold er
hans Tanke. Hans ydre Liv faaer ikke umiddelbar Interesse som hos en
udadtil handlende Personlighed, hvis Liv afpræger sig i hans Samtid som
det selv bærer Afpræget af denne. Det bliver enten fuldkommen ligegyldigt
eller kun forsaavidt af Betydning, som dets Overensstemmelse med hans
Tænkning afgiver Borgen for dennes Magt til at gjennemtrænge et
Menneskeliv. Interessen bliver en afledet, væsentlig kan den ikke blive,
thi i Tankeindholdets Tilegnelse forsvinder det personlige Forhold til
Tænkeren, og maa forsvinde, hvis Tilegnelsen skal blive fuldstændig og
fri. Man kan derfor hos en Tænker, selv om man fordrer mere end hans
Tankes Genesis, nøies med et let Omrids af hans Liv, naar kun dette Omrids
er bestemt nok til, at hans Tankes Præg bliver synligt derigjennem; man
kan give Afkald paa alle de selv i deres tilsyneladende Tilfældighed
betydningsfulde Enkeltheder, der gjøre Omridset til et udført Billede.
Hvad Spinoza angaaer, saa er en saadan Nøisomhed nødvendig; han har kun
funden een noget udførligere Biograph, den lutherske Præst
\emph{Colerus}\footnote{La vie de Spinoza, tirée des écrits de ce fameux
philosophe et du témoignage de plusieurs personnes dignes de foi, qui
l'ont connu particulièrement. Haag 1706. 8. En tidligere hollandsk Udgave
er udkommen i Utrecht.},
% --- p. 28 ---
og denne enfoldige, redelige, trods den theologiske Antipathie med den
uforstaaede Philosoph anerkjendende, Fortæller har i Spinozas simple,
udadtil ensformige Liv kun funden ringe Stof til at meddele, og hans
Meddelelser fuldstændiggjøres kun lidet ved de Enkeltheder, der findes i
Fortalen til Spinozas efterladte Skrifter, hos \emph{Bayle} (i hans
Dictionnaire historique et critique), i »La vie de Spinoza par un de ses
disciples« (sandsynligviis Lægen \emph{Lucas}), Amsterdam 1719, samt hos
\emph{Nicéron}, Kortholt og Leibnitz.
% `Kortholt og Leibnitz' on the same line as the spaced `N i c é r o n' is
% NOT spaced (600 dpi). No footnote on this page.

Af Charakteren i Spinozas Tænkning, saaledes som denne allerede viser sig
i den Maade, paa hvilken han opfattede sin Forgængers Princip og førte det
videre, kunde man vente, at hans Liv ikke vilde være som mange af
Tænkernes i Overgangsperioden, navnlig Italienernes, hvem Tanken greb som
en rastløs, aldrig tilfredsstillet, fortærende Lidenskab, og, idet den
fyldte dem som med prophetisk Raseri, tumlede dem ustadigt om i Verden som
sine Forkyndere, lod dem, stolte af deres Kald, trodsigt støde sammen med
det Bestaaende, forat ende som Ideens Martyrer. Hos Spinoza maae vi vente,
at Grundanskuelsens Ro vil kaste sit Skjær over hans Liv, at det Endeliges
Forsvinden i det Uendelige vil afpræge sig i Resignation, i en Opgiven af
alle Livets concretere Forhold for kun at leve for Erkjendelsen, for
Tankens Enhed med det Guddommelige i intellectuel Kjærlighed.

Dette finde vi ogsaa bekræftet naar vi læse de faa og tarvelige
Efterretninger, der ere givne os om Spinoza. De ere utilstrækkelige til at
danne os et helt og fyldigt Billede af et menneskeligt Liv, tilstrækkelige
til at lade os forfølge et Brud med Alt, hvad der ikke er det Uendelige,
med Nationalitet, med den fædrene Tro, med Kjærlighed til en Qvinde, med
ydre Anseelse, for alene at leve i og for det Høieste. ---

\emph{Baruch Spinoza} blev født i Amsterdam den 24de November 1632. Hans
Fader var en jødisk Kjøbmand, der paa Grund af Religionsforfølgelser havde
forladt sit Fædreland Portugal. Af de faa Træk, der findes af ham hos
% --- p. 29 ---
Sønnens Biographer, faaer man Indtrykket af en klar Forstand og Uhildethed
i religiøse Anskuelser. Da Familiens Kaar kun vare middelmaadige og Sønnen
tidlig røbede udmærkede Evner, blev han bestemt til den lærde Stand og
betroet til Rabbineres Undervisning, blandt hvilke \emph{Moses Morteira}
har et Navn i den ny-hebraiske Literatur. Denne Undervisning vedblev efter
Spinozas Forældres tidlige Død, og til den knyttedes senere, efter
Talmudens Forskrift, Oplærelsen i et Haandværk, Forfærdigelsen af optiske
Glas, hvori han efterhaanden opnaaede en stor Fuldkommenhed. Det varede
ikke længe, inden hans klare Forstand, skærpet ved den talmudiske
Spidsfindighed, men drevet af en Tørst efter Sandhed, der ikke lod den
finde Ro i Rabbinernes kunstige Bygning, bragte hans Lærere i Forlegenhed
og paa een Gang vakte hans Religionsfællers Stolthed og deres hemmelige
Uro. Allerede før han var 20 Aar gammel, var Troen paa Rabbinernes
Autoritet tilintetgjort hos ham, og hans Tanke, der ikke fandt
Tilfredsstillelse i deres Lære, søgte Sandheden i sig selv. Desværre ere
netop for de Aar, der ligge imellem det første Spor af en vaagnende
Selvstændighed hos Spinoza og hans afgjørende Brud med sit Folk og sin
Religion, Biographerne fattige og uoverensstemmende. Af indre Grunde synes
Forfatteren til »la vie de Spinoza« paa dette Punkt at fortjene Tiltro
fremfor Colerus, der tænker sig Gangen i Spinozas Opdragelse efter de
christelige lærde Skolers Mønster, og derfor paa et tidligere Punkt
bringer humanistiske Udviklingselementer ind, med hvilke Spinoza først
senere kom i Berøring. Efter hin Biographs Fremstilling synes det, som om
Spinoza i Taushed har ladet sine Tanker modnes, men efterhaanden som de
antoge en fastere Form, og hans Tvivl fra den jødiske Videnskab naaede til
den jødiske Religion, er hans Læreres og Troesfællers Mistillid bleven
vakt, og den er bleven bestyrket derved, at han mere og mere fjernede sig
fra Synagogen og søgte Omgang med Christne. Hans Lærere, der i ham havde
seet en vordende Støtte for Synagogen, bleve ham nu fjendske, og søgte ved
Magtsprog at bringe ham
% --- p. 30 ---
tilbage, Menigheden søgte at kjøbe ham tilbage ved et aarligt Stipendium
paa 1000 Gylden, og, da det ikke hjalp, ved Snigmord at skaffe sig af med
ham. Ogsaa dette slog feil, og nu greb man til det sidste Middel, at
excommunicere den Frafaldne, og det lykkedes snart efter Rabbinerne,
understøttede af de reformerte Geistlige, at udvirke hos Magistraten i
Amsterdam, at han for en Tid blev forvist fra Byen. Fra denne Tid af
staaer Spinoza som fremmed for sit Folk og dets Religion. Sit jødiske Navn
Baruch ombytter han med dets latinske Form Benedictus, og Philosophien
bliver hans Livs Syssel.

Allerede paa den Tid, da han begyndte at føle sig fremmed i sit Folk,
havde han søgt, udenfor den jødiske Videnskabs snævre Kreds at skaffe sig
Midlerne til en rigere videnskabelig Udvikling og til humanistisk
Dannelse. Hos Fritænkeren, Lægen \emph{Frants van den Ende}, lærte han de
classiske Sprog, navnlig det latinske, og ved dettes Hjælp skaffede han
sig Adgang til Studiet af Naturvidenskaberne og Philosophien, i hvilken
han snart valgte Descartes til Veileder. Da han var bleven
excommuniceret, optog van den Ende ham en Tidlang i sit Huus, og over
denne Periode af hans Liv falder der et romantisk Skjær ved hans
Kjærlighed til sin Lærers aandrige Datter. Hun foretrak en af hans
Meddisciple, og hans Kjærlighed søgte nu sin Gjenstand i det
Guddommelige.

Spinozas Excommunication af Synagogen falder omtrent i Aaret 1660, og i
samme Aar vistnok hans Fjernelse fra Amsterdam. En kort Tid tilbragte han
hos en Ven paa Landet i Nærheden af Byen, det næste Aar bosatte han sig i
Landsbyen \emph{Rynsburg} ved Leyden, hvor han i 1663 udgav sit første
Skrift: Principia philosophiæ Cartesianæ, more geometrico demonstrata, til
hvilket Cogitata metaphysica danne et Anhang. I 1664 flyttede han til
\emph{Voorburg} ved Haag, og derfra, efter et femaarigt Ophold, paa sine
Venners Opfordring til selve \emph{Haag}, hvor han levede Resten af sit
Liv.

% p. 31 opens a new paragraph (the first line is indented on the image).
% --- p. 31 ---
Om denne Deel af Spinozas Liv er Colerus udførligere, idetmindste
udførlig nok til, at hans Fortælling giver os et Indtryk af dets strænge,
alvorlige, ensomme, energiske Charakteer. Vi see Spinoza anstrængt og
uafbrudt arbeidende henimod sit Maal: Erkjendelsen af Gud, af Naturen, og
af den Enhed, i hvilken vi ere med denne, og imod den rolige Ophøiethed
over de menneskelige Lidenskaber, der erhverves gjennem Erkjendelsens
Klarhed. Engang, fortæller Colerus, havde han i 3 Maaneder ikke forladt
det simple Kammer, han beboede hos Maleren van der Spyck i Haag, saaledes
fordybede han sig i sine Studier. Ved sin overordentlige Nøisomhed gjorde
han det muligt for sig, skjøndt han maatte arbeide for sit Livs Ophold, at
kunne anvende sin meste Tid paa disse. Hans Biograph fortæller os
omstændelig og med naive Details, hvor tarveligt han levede, hvor Lidt han
brugte, skjøndt han altid i sit Ydre vedligeholdt en vis beskeden
Elegance, hvor faa og simple Efterladenskaber der fandtes efter ham, og
hvor uegennyttig han var i Forholdet til sine Søstre og sine rigere Venner
og Velyndere.

Foruden de Mange, hvem Nysgjerrigheden lokkede til ham, var der Mange, der
paavirkedes af Spinozas Personlighed. Hans dybe Viden imponerede dem,
medens hans livlige og vittige Samtale, hans Humanitet, hans Aabenhed og
Trofasthed, og hans Sæders Adel (jfr. Oldenburgs Yttringer i Ep. 1)
tiltrak, men kun faa af dem forstode ham, og hans Forhold til dem blev
mere et humant Velvilliesforhold, deres til ham et beundrende. Af hans
Breve, der alle ere Svarskrivelser, seer man, at kun een eneste af hans
Correspondenter, den Anonym, der har skrevet de med forskjellige Mærker
betegnede Breve 61, 63, 65, 67, 69 og 71, har været istand til at trænge
dybere ind i hans Philosophie, og kun en Enhed i Tænkning kunde for
Spinoza, hvis Liv var hans Tænkning, begrunde et virkeligt
Venskabsforhold. Hvor lidet han blev forstaaet af dem, der nærmede sig til
ham og følte sig tiltrukne af ham, kan man bedst slutte sig til af hans
% `Brevvexling med' is NOT letterspaced (600 dpi) -- the line is set very
% tight and only `O l d e n b u r g' is spaced.
Brevvexling med \emph{Oldenburg} (jfr. S. 40), der med en enkelt
Afbrydelse fortsattes fra 1661 til Spinozas Død. Selv af de
% --- p. 32 ---
\setcounter{footnote}{0}%
sidste af disse Breve viser det sig, at Oldenburg, der dog tilbragte sit
Liv i Omgang med Philosopher og Videnskabsmænd, endnu ikke havde formaaet
at fatte een eneste af Spinozas Tanker, og det Samme maatte naturligviis
være Tilfældet med alle dem, der vilde følge Spinoza uden at blive sig
bevidst, at de maatte opgive hele den Grundvold, paa hvilken deres
Forestillingskreds hvilede.

Allerede før Spinozas Skrifter havde henledt en almindeligere
Opmærksomhed paa ham, see vi af hans Breve, at han er bleven bemærket af
Lærde og draget ind i den Correspondance, der for hin Tids videnskabelige
Liv var et saa vigtigt og saa ivrigt benyttet
Hjælpemiddel\footnote{Spinozas Correspondance synes, efter Yttringer i L.
Meyer's Fortale til hans Opera posthuma, at have været temmelig vidtløftig,
og den Samling af Breve, der findes i hans Skrifter (74) har vistnok kun
udgjort en lille Deel deraf. Paa det store kongelige Bibliothek her i Byen
findes en utrykt Billet fra Spinoza til \emph{Grævius}, men dens Indhold
er uden Betydning.}. Hans første Skrift skaffede ham strax et Navn, men
fremkaldte ogsaa Angreb fra de strænge Cartesianere\footnote{Om Spinozas
Forhold til disse jfr. Ep. 19.}. Endnu mere steg hans Ry ved den 1670
udgivne tractatus theologico-politicus\footnote{Ikke blot Lærde bleve
opmærksomme paa ham, men ogsaa politiske Personligheder, saaledes, foruden
Jan de Witt, Prindsen af Condé, der indbød Spinoza til en Sammenkomst med
sig i Utrecht 1673, og Fr. af Luxembourg.}, men i et endnu stærkere
Forhold end hans Ry voxede Antallet af hans Modstandere, især blandt
Geistligheden, og disses fanatiske Angreb vare Aarsagen til, at Spinozas
vigtigste Skrifter først bleve udgivne efter hans Død. Kun ved denne
Tilbageholdenhed kunde han skaffe sig den Ro, han satte saa høi Pris paa,
og som var ham saa nødvendig, og for dens Skyld afslog han ogsaa det
hæderlige Tilbud om at overtage et philosophisk Professorat i Heidelberg,
der i 1673 blev gjort ham fra den frisindede Churfyrste \emph{Carl
Ludvig} gjennem dennes, over dette Hverv øiensynlig kun
% --- p. 33 ---
lidet tilfredse Raad, den heidelbergske Professor \emph{Ludvig
Fabritius}.

\emph{I de politiske Begivenheder} i sit Fædreland tog Spinoza ikke activ
Deel, men han fulgte dem som Iagttager med den Interesse, som hans
Anskuelse af Statsforholdet maatte give dem. Hans Opfattelse af Frihedens
Uundværlighed til Opnaaelsen af Statens Øiemed (jfr. Cap. 7) gjorde ham
til en Tilhænger af den republicanske Statsform og knyttede ham til dem,
der forsvarede denne mod monarchiske Tendentser, navnlig til \emph{Jan de
Witt}, der blev hans Beskytter og Velgjører, og hvis beklagelige Død
opfyldte ham med bitter Smerte.

Hans Forhold til \emph{Religionen} var givet i hans Opfattelse af dennes
Væsen. For ham selv var hans Philosophie hans Religion: Erkjendelsen af
Gud er Kjærligheden til ham og Enheden med ham. Ligeoverfor Andre, for
hvem de almindelige religiøse Forestillinger havde bevaret deres
Betydning, søgte han at fastholde det praktiske og humane Element i
Religionen, og at opfatte den saaledes, at den kunde blive en Tilnærmelse
til den sande Erkjendelse og til den sande Overensstemmelse mellem
Menneskene. Han fremdrog, med Omgaaen af Dogmerne, det, der kunde virke
tilskyndende til Kjærlighed, Samdrægtighed og Redelighed, og fremhævede
Lydighedens Betydning og den fælles Gudsdyrkelses Værd. Netop dette
dannede den hyppigste Gjenstand for de Samtaler med hans Huusfæller, i
hvilke han søgte en Hvile fra mere anstrængende Beskæftigelse, og i hvilke
han forstod at blande Alvor med Venlighed og tidt med muntert Lune og
godmodig Ironie.

Spinozas Helbred havde allerede fra hans Ungdom været svagt; anstrængt
aandeligt Arbeide og Nattevaagen nedbrød det endnu mere, og der udviklede
sig en Svindsot, som den 21de Februar 1677 gjorde Ende paa hans Liv i hans
45de Aar. Hans christelige Biograph fortæller omstændelig hvor taalmodig
og standhaftig han havde baaret sine Lidelser, hvor rolig og blid hans Død
var, og han afviser de usande Rygter, der vare i Omløb angaaende hans
Dødsmaade. En
% the `3' at the foot of p. 33 is the signature mark of the gathering and is
% not transcribed. No footnote on this page.
% --- p. 34 ---
\setcounter{footnote}{0}%
anden Theolog (\emph{Kortholt}) undrer sig over, at en Atheist kunde faae
et saa smukt Endeligt, og de, der ikke kunde begribe dette, og ikke havde
saa gode Oplysninger som Colerus, udbredte Rygter om, at Spinoza skulde
være død af Skræk for at komme til at dele Skjæbne med Jan de Witt, eller
ved Selvmord.

Et Portræt af Spinoza, der endnu er bevaret, viser os et Ansigt, i hvilket
den mørke Hudfarve, den bøiede Næse, den stærkt udprægede Mund og Hage
tyde paa hans østerlandske Nationalitet, medens den høie, kraftigt formede
Pande, Øienbrynenes faste og fine Linie, og det levende, klare Øie røbe
den skarpe Tænker. Ansigtets Udtryk bærer Præget af Resignationens Alvor,
men ogsaa af den Ro i det Uendelige, der erhverves gjennem den.

Af Spinozas \emph{Skrifter} ere de to, der udkom medens han levede,
allerede ovenfor blevne nævnede. Det første af disse: \emph{Renati Des
Cartes Principiorum philosophiæ pars I \& II}, more geometrico
demonstratæ par \emph{Benedictum de Spinoza}, Amstelodamensem, --- til
hvilket \emph{Cogitata metaphysica}, in quibus difficiliores, quæ tam in
% Extent of the spaced runs verified at 600 dpi: `more geometrico
% demonstratæ par' and `Amstelodamensem' are NOT spaced, though they stand
% inside the same title; the compositor spaces only the four blocks marked.
parte metaphysices generali quam speciali occurrunt quæstiones breviter
explicantur, slutte sig som Anhang, er oprindelig blevet til ved en ydre
Foranledning, idet en Deel deraf er udarbeidet til Brug for en Discipel,
hvem Spinoza underviste i den cartesiske Philosophie, men ikke ansaae for
moden til at indvie i sine egne, allerede dengang selvstændig udviklede,
Anskuelser (jfr. Ep. 9). Idet Spinoza, paa sine Venners Opfordring,
bearbeidede det til Udgivelse, er, under Fremstillingen af Descartes's
Lære, de Punkter antydede, hvor Vanskeligheder, som ere uopløselige fra
dennes Standpunkt, gjøre Consequentser nødvendige, der lede hen til
Spinozas\footnote{Det er navnlig ved Problemet om \emph{Guds Forhold til
Udstrækningen} (jfr. pr. ph. C. I, 9. schol. 16. 21. II, 2. schol., Cog.
met. I, c. 2.), ved Bestemmelsen af Begrebet \emph{creatio} (jfr. pr. ph.
C. I, 12. Cog. met. II, c. 10), og af Begrebet \emph{endelige Ting},
existerende \emph{extra Deum} og \emph{in se}, ved Spørgsmaalet om dette
sidste Begrebs Forenelighed med Guds Egenskaber, som Alvidenhed,
Allestedsnærværelse, Uforanderlighed (jfr. pr. ph. C. I, 11 dem. 12,
coroll. 2. Cog. met. I, c. 2. II, c. 1. 3. 7. 10 f.), samt endelig ved
Problemet om \emph{den menneskelige Frihed} (jfr. Cog. met. I, c. 3. II,
c. 11 f.), at der vise sig Vanskeligheder, til hvis Løsning fra Spinozas
eget Standpunkt der ere givne Antydninger (saaledes f. Ex. med Hensyn til
Begrebet Skabelse, der ophæves idet det indskrænkes til at gjælde
existentia, ikke essentia, og idet dets Anvendelse paa modi benægtes). Om
sit Forhold til Descartes har Spinoza ogsaa udtalt sig i Ep. 2 (fra
1661).},
% This note begins on p. 34 (`... existe-') and its remaining nine lines
% stand at the head of p. 35's note block, above p. 35's own note 1, with no
% repeated mark. The whole note is attached here to its mark, so p. 35's
% \setcounter still leaves its own note as number 1.
% `til' in the spaced run `G u d s  F o r h o l d  t i l' shows a dotless
% second sort at 600 dpi and could be read `tll'; the sense-reading is
% transcribed and no sort-level claim is made (JBIG2 scan, header §6).
og i Fremstillingens hele Anlæg er der banet
% --- p. 35 ---
\setcounter{footnote}{0}%
Vei for hans egen Opfattelse. Denne Skriftets Charakteer antyder Hensigten
med dets Offentliggjørelse.

Medens dette Skrift saaledes var bestemt til at forberede Modtagelsen af
Spinozas Philosophie, skulde det næste skaffe denne Plads ved at afstikke
Grændsen mellem Theologien og Philosophien og derved sikkre denne et
neutralt Gebeet, samt ved at paavise den frie Tænknings og den frie Tales
Forenelighed med Statens Øiemed og med dens Ro. Dette er den nærmeste
Opgave for den 1670 udkomne \emph{Tractatus theologico-politicus},
continens dissertationes aliquot, quibus ostenditur, libertatem
philosophandi non tantum salva pietate et reipublicæ pace posse concedi,
sed eandem nisi cum pace reipublicæ ipsaque pietate tolli non posse. Denne
Afhandlings væsentlige Indhold vil blive udviklet i Cap. 7, hvor Spinozas
Opfattelse af Religionen og Staten fremstilles. Den udkom uden Spinozas
Navn og blev paa Geistlighedens Anstiftelse hurtigt undertrykt, men faa
Aar efter (1673) atter udgivet under forskjellige fingerede
Titler\footnote{Som; \emph{Danielis Heinsii} operum historicorum Collectio
prima. Editio secunda priori editione multo emendatior et auctior.
Accesserunt quædam hactenus inedita. Lugd. Batav. 1673. 8. --- og:
\emph{Francisci Henriquez de Villacorta}, doctoris medicinæ, a cubiculo
regali Philippi IV et Caroli II archiatri, opera chirurgica omnia, sub
auspiciis potentissimi Hispanorum regis Caroli II. Amstelod. 1673. 8.}.
% PRINTER'S DEFECT, as printed: `Som;' for `Som:' at the head of the note.
% Verified at 600 dpi (a dot over a comma-tail, not two dots); both OCR
% witnesses read a semicolon. A punctuation fault is structural and survives
% the JBIG2 compression, so it stays logged.
Sildigere var det Spinozas Hensigt
% the `3 *' at the foot of p. 35 is the signature mark and is not transcribed.
% --- p. 36 ---
\setcounter{footnote}{0}%
at udgive den paany, ledsaget med Anmærkninger\footnote{Jfr. Ep. 75
(tilføiet i Bruders Udgave). Afhandlingen omtales desuden i Ep. 8. 14
(„tractatum de scriptura“), 17. 19 f. 47 og 49.}, men han blev forhindret
deri ved sin tidlige Død. De Anmærkninger, han vilde have tilføiet, ere
optagne i den franske Oversættelse af \emph{St. Glain}, der udkom 1678 i
Leyden under Titlen: La clef du sanctuaire par un savant homme de notre
% DOUBTFUL READING, not a defect. The scan shows `nnder'; `under' is set
% here. This page is the one on which the JBIG2 test caught a single
% dictionary bitmap standing in `sur', in `under Titelen' and in
% `tractatus' -- i.e. serving as both u and n within one page
% (JBIG2-TEST.md, bitmap A). Whatever this glyph looks like at 600 dpi is
% the encoder's reconstruction, so the sense-reading stands; the compositor
% is not accused of a turned sort.
siècle. Foruden denne franske Oversættelse, til hvilken der ogsaa findes
to andre Titler: Traité des cérémonies superstitieuses des Juifs tant
anciens que modernes, og: Reflexions curieuses d'un esprit désinteressé
sur les matières les plus importantes au salut tant public que
particulier; --- gives der en hollandsk af \emph{Joh. Henrik Glasemaker}
fra 1693, og flere nyere franske og tydske.

Først efter Spinozas Død\footnote{Kort før sin Død tilintetgjorde Spinoza
en Afhandling: de Iride, og en paabegyndt hollandsk Oversættelse af det
gamle Testamente med Noter, af hvilken Pentateuchen var færdig. Tabte ere
Størstedelen af hans Breve, en Afhandling om Djævelen, og et paa Spansk
skrevet Forsvarsskrift for hans Udtrædelse af Synagogen.} bleve de
Skrifter udgivne, der ere den egentlige Kilde til Kundskaben om hans
Philosophie. Hans Ven \emph{Ludvig Meyer} udgav i hans Dødsaar (1677)
denne Samling under Titelen: B. d. S. Opera posthuma (det var efter
Spinozas Ønske, at hans Navn ikke nævnedes). Den indeholder, foruden
Spinozas Hovedværk: \emph{Ethica}, ordine geometrico
demonstrata\footnote{Den er delt i 5 Dele: de Deo; de natura et origine
mentis; de origine et natura affectuum; de servitute humana seu de
affectuum viribus; de potentia intellectus seu de libertate humana.}, de
to ufuldendte Afhandlinger: \emph{tractatus politicus} og: \emph{tractatus
de emendatione intellectus}, samt 74\footnote{Det 75de Brev, der er
tilføiet i \emph{Bruders} Udgave, er for faa Aar siden fundet i Holland og
udgivet af \emph{H. W. Tydemann}. Det er til Lambert van Velthuysen.}
Breve fra og til ham, og, som Anhang, et ufuldendt \emph{Compendium
grammatices linguæ hebrææ}. Samlingen blev i samme Aar oversat paa
Hollandsk af \emph{Jarrig Jellis} under Titelen: \emph{De nagelate
Schriften van B. d. S.}
% `B. d. S.' IS inside the spaced run at the end of the paragraph (600 dpi),
% but the earlier `B. d. S. Opera posthuma' is set normally.

% p. 37 opens a new paragraph (the first line is indented on the image).
% --- p. 37 ---
\setcounter{footnote}{0}%
Til \emph{Ethiken} havde Spinoza allerede tidlig gjort et Udkast. I det
ældste af hans Breve til \emph{Oldenburg}, fra 1661, angives de første
Definitioner, om end i en fra Ethikens nuværende Redaction noget
forskjellig Form (jfr. Cap. 2), samt Hovedsætningerne af Ethikens første
Deel. Dette er af Vigtighed for Spinozas Udvikling, idet det afgiver et
Beviis for, at han allerede da han udgav Principia phil. Cartes. var paa
det Rene med sit fra Descartes afvigende Grundbegreb\footnote{For Spinozas
Forhold til Descartes bliver det ogsaa af Betydning, at han først efter
lang selvstændig Tænkning har begyndt Studiet af dennes Skrifter.}. De
væsentlige Mangler, han finder ved sin Forgængers Philosophie, angiver han
i det samme Brev (Ep. 2). I et Brev fra \emph{Simon de Vries}, der er
skrevet i Februar 1663 (Ep. 26) citeres bestemte Sætninger af Ethikens
første Bog. I et Brev til \emph{Blyenbergh} fra 1665 (Ep. 36, jfr. Ep. 38)
henviser Spinoza med Hensyn til Attraaens Begreb til Udviklingen i hans
»Ethica necdum edita«, hvor den findes i tredie Bog. I Brevvexlingen med
L. \emph{Meyer} (Ep. 63---71) fra Spinozas sidste Leveaar (1675---76)
citeres hyppigt Sætninger af de forskjellige Dele af Ethiken, der da
vistnok var bragt i dens endelige Form. Endelig opfordrer
\emph{Oldenburg} i et Brev fra 1675 (Ep. 18) Spinoza til at udgive
»tractatum illum --- --- quinque-partitum«. Af disse Citater, og af mange
Udviklinger i Brevene, der saagodtsom ordret stemme overeens med Ethiken,
sees det saaledes, at dette Spinozas Hovedværk længe havde været færdigt i
Udkast, i hvilket det var meddelt Adskillige af dem, der stode ham
nærmest\footnote{Dette siges ogsaa udtrykkeligt i Udgiverens Fortale: „ea
(ɔ: Ethiken) namque ante multos annos a diversis fuit descripta et iis
communicata.“}; men af Varianterne i Citaterne, der, hvor de ere af
Vigtighed for Begrebsbestemmelserne, senere ville blive berørte, seer man
tillige, at han bestandig paany beskæftigede sig dermed forat finde det
exacteste Udtryk for sin Tanke.
% The second footnote mark on p. 37 (after `nærmest') is printed as a
% superscript numeral WITHOUT the base serif of the `1' used four lines from
% the top of the page, and it is certainly not a 2; the note it points to is
% numbered 2) in the note block. Read on the image at 600 dpi. Recorded, not
% adjudicated: no sort-level claim is made from this scan (header §6).
% First `ɔ:' of the book, in note 2 above (600 dpi).

% p. 38 opens a new paragraph (the first line is indented on the image).
% --- p. 38 ---
Afhandlingen \emph{de intellectus emendatione} er, efter Udgiverens
Udsagn, et af Spinozas første Arbeider, med hvilket han idelig
beskæftigede sig, men uden nogensinde at bringe det til Afslutning. Den
omtales allerede 1663 i et Brev fra Oldenburg (Ep. 8) som: opusculum, in
quo de rerum primordio earumque dependentia a prima causa, ut et de
intellectus nostri emendatione tractas. Senere omtales den i flere Breve
(Ep. 42. 63 f.). Det Væsentlige af dens Indhold (jfr. Cap. 5) er optaget i
Ethikens 2den Deel, navnlig i Scholierne til Sætningerne 17 og 40. Den
udvikler, efterat have bestemt Erkjendelsen af det Evige og vor Enhed med
det som det høieste og sande Gode i Modsætning til de usande,
forgængelige Goder, og efterat have stillet det som Opgaven, at naae denne
Erkjendelse, de forskjellige Erkjendelsesformer, af hvilke
Fornufterkjendelsen viser sig at være den eneste, der fører til Maalet.
Dennes Redskaber ere de sande Ideer, der indeslutte Visheden i sig, og den
sande Methode er den, der viser, hvorledes Aanden sammenkjæder Ideerne
efter dens høieste Idees Norm. I Afhandlingens \emph{første} Deel
adskilles nu de sande Ideer fra ideæ fictæ, falsæ og dubiæ, der tilhøre
Imaginationen, og dennes Forskjel fra intellectus, de sande Ideers Kilde,
bestemmes. I \emph{anden} Deel gives Reglerne for, hvorledes det
Ubekjendte med Sikkerhed udledes af det Bekjendte, i en saadan
Sammenknytning, at vor Aand bliver den subjective Gjengivelse af Naturens
objective Væsen som Hele og i dens Dele. Her føres nu, gjennem
Adskillelsen mellem det, der begribes ved den nærmeste Aarsag, og det, der
begribes ved sit Væsen alene (sola essentia) til at afhandle Lovene for
Definitionerne af de »skabte og uskabte« Ting, idet Erkjendelsen kun
skrider frem gjennem Definitionerne, der udtrykke Tingenes eiendommelige
Væsen. Idet der nu skulde gaaes over til Angivelsen af de Midler, ved
hvilke res particulares æternæ, Fornuftens væsentlige Gjenstand, kunne
udledes af det høieste Begreb, afbryder Afhandlingen paa det afgjørende
Punkt, efterat der endnu kun i Korthed er handlet om Forstandens Kræfter
og dens Egenskaber.
% --- p. 39 ---
Efter Planen for Afhandlingen (S. 371) stod det endnu tilbage, først at
udvikle Ordenen i Erkjendelsens Sammenknytning, og dernæst det
fuldkomneste Væsens Idee. Det vil ved Udviklingen af Spinozismens
Grundbegreber (Cap. 2) vise sig, at der netop paa det Punkt, hvor
Afhandlingen afbrødes, opstaaer en uovervindelig Vanskelighed for Spinoza,
saa at det ikke er tilfældigt, at den har maattet blive ufuldendt.

\emph{Tractatus politicus}, der ligeledes kun er et Fragment, er skrevet
kort før Spinozas Død. I de 5 første Capitler udvikler han med en
overordentlig Skarphed og Klarhed sin Opfattelse af Naturretten, af
Begrebet Stat, af Statsmagtens Ret, dens Myndigheds Omfang, Staternes
Forhold til hverandre og Statssamfundets Øiemed. I det 6te og 7de Capitel
omhandler han Monarchiet, i det 8de---10de Aristokratiet, baade det, hvor
een By er hele Statens Hoved, og hvor flere Byer danne en Forbundsstat. I
det 11te Capitel er Udviklingen af Demokratiet begyndt, men her er
Afhandlingen afbrudt efterat nogle faa almindelige Bestemmelser ere givne.
Foruden den fuldstændige Fremstilling af denne Statsform vilde Spinoza
endnu (jfr. Fortalen) have omhandlet »Lovene og andre specielle politiske
Spørgsmaal.«
% DOUBTFUL READING, not a defect. The scan shows `Fornden'; `Foruden' is
% set here. As at `Elemeut' (p. 9) and `nnder' (p. 36): the appearance of
% the glyph at 600 dpi and the agreement of the two OCR witnesses are both
% worthless on an n/u question in this file, which demonstrably places one
% bitmap in both roles (JBIG2-TEST.md). Sense-reading; no defect claimed.
% `6te', `7de', `8de---10de', `11te Capitel' are Brøchner's references to
% Spinoza's own chapters, not heads in this book.

Den Samling af \emph{Breve}, der findes i Opera posthuma, udgjør vistnok
kun en ringe Deel af den Brevvexling, Spinoza har ført. Den indeholder
foruden eet Brev (Ep. 48), der er skrevet af en Fremmed til en af Spinozas
Bekjendte angaaende tractatus theologico-politicus, 32 Breve til og 41 fra
Spinoza. De fleste af Brevene ere oprindelig skrevne paa Latin, 24 (nemlig
30---41, 43---47, 50 og 55---60) paa Hollandsk og oversatte paa Latin af
L. Meyer. Foruden adskillige Oplysninger om Spinozas personlige Forhold og
om hans Skrifter, indeholde Brevene talrige Udviklinger, dels af Punkter,
der ere berørte i hans Værker, men som her i Regelen komme frem i en
modificeret Fremstilling, fremkaldt ved Spørgsmaal og Misforstaaelser;
dels af saadanne Punkter, som enten slet ikke eller kun i Forbigaaende
% --- p. 40 ---
ere omtalte andensteds (saaledes af forskjellige christelige Dogmer, af
physiske, mathematiske og optiske Problemer, af logiske Bestemmelser o.
desl.). En stor Deel af Brevene, og det netop af de i philosophisk
Henseende vigtigste, ere til og fra unavngivne Forfattere. Under flere af
de forskjellige Mærker synes een og samme Forfatter, Udgiveren af hans
efterladte Skrifter, at skjule sig. Til ham er nemlig ikke blot Ep. 29,
men udentvivl ogsaa Ep. 62, 64, 66, 68, 70 og 72, skjøndt de Breve, disse
ere Svar paa, ere daterede fra forskjellige fremmede Byer, og tildels
udgive sig for at være skrevne af en Anden end den, der har sendt dem til
Spinoza. Blandt de navngivne Correspondenter er, næstefter
\emph{Leibnitz}, med hvem Spinoza har vexlet nogle faa Ord i Anledning af
et tilsendt Skrift, \emph{Heinrich Oldenburg} den bekjendteste, og den,
til hvem det største Antal Breve er rettet. Han var under Cromwell den
nedersachsiske Kredses Gesandt i London og beklædte siden den samme Post
under Carl II, der havde Godhed for ham, og, da det engelske
Videnskabernes Selskab i 1662 blev kongeligt, gjorde ham til dets
Secretær. Han var en ivrig Mellemmand mellem flere af hin Tids Philosopher
og Lærde, og, trods sin ringe Evne til at opfatte en dybere Tænkning, ikke
uden Blik for Talenternes Eiendommelighed og Anerkjendelse af denne, og om
end ikke fri for et Forfængelighedens Motiv i sin Tilnærmelse til overlegne
Personligheder, dog besjælet af en virkelig Interesse for Videnskabernes
Udvikling. De andre navngivne Correspondenter ere: Spinozas Ven
\emph{Simon de Vries}, den dortrechtske Kjøbmand \emph{Willem van
Blyenbergh}, den heidelbergske Professor \emph{Joh. Lud. Fabritius},
\emph{Albert Burgh}, der gjorde et uheldigt Forsøg paa at omvende Spinoza
til Katholicismen, og \emph{Peter Balling}. Med Forbogstaverne alene er
bl. A. betegnede \emph{Ludvig Meyer} (Ep. 29) og den til Jødedommen
overtraadte Spanier \emph{Isaac Orobio de Castro}, samt maaskee \emph{Joh.
Bredenburg} (der 1675 skrev mod tractatus theol.-pol.) og Mennoniten
\emph{Jarrig Jellis}, Oversætteren af Spinozas Opera posthuma.

% p. 41 opens a new paragraph (the first line is indented on the image).
% --- p. 41 ---
Spinozas \emph{Compendium grammatices linguæ hebrææ}, der danner et
Anhang til hans efterladte Skrifter, er ikke af nogen Betydning for hans
Philosophie; idethøieste kunne Enkeltheder af Terminologien (saaledes
Anvendelsen af Betegnelsen: \emph{Attribut}) og Opfattelsen af
Ordclassernes Forhold til hverandre, idet de alle, med Undtagelse af
Interjectioner og Conjunctioner og enkelte Partikler, bringes ind under
Begrebet: Nomen, have nogen Interesse.

Foruden de ved Spinozas Død ufuldendte Skrifter, havde han endnu, efter
Meyer's Udsagn, lagt Planen til en »\emph{integra philosophia}«, til
hvilken der hyppigt henvises i tractatus de emendatione intellectus (i
Fortalen og i mange af Noterne) og maaskee i Ep. 74; samt en »Algebra
breviori et magis intelligibili methodo«.

I hvilket Forhold denne \emph{integra philosophia} vilde staae til
Ethiken, og i hvilke enkelte Videnskaber den vilde sondre sig, har Spinoza
ikke angivet. I Indledningen til Afhandlingen de intellectus emendatione
opregner han, men, som han udtrykkelig tilføier, uden systematisk Orden,
de forskjellige Videnskaber, der alle sigte hen imod det samme Maal, den
høieste menneskelige Fuldkommenhed ɔ: den for alle Mennesker fælles
Erkjendelse af den unio, quam mens cum tota natura habet. Dertil, siger
han, hører først Erkjendelsen af Naturen, dernæst et Samfund, der saa let
og sikkert som muligt kan føre til Erkjendelsen. Desuden maa der lægges
Vind paa Moralphilosophien og paa Læren om Børnenes Opdragelse, og,
eftersom Helbredet ikke er noget ringe Middel til at naae til det
tilsigtede Maal, maa der opstilles en fuldstændig Lægevidenskab, og da
Meget, der er vanskeligt, gjøres let ved Kunsten, og den kan spare os
megen Tid og yde os megen Fordeel i Livet, maa Mechaniken paa ingen Maade
forsømmes. Men fremfor Alt maa der udfindes en Maade til at helbrede
Forstanden, og, saavidt det forud for Videnskaben (initio) kan skee, at
rense den, saa at den begriber Tingene med Lethed, fuldkomment og uden
Vildfarelse. --- Denne Opregning angiver i det Væsentlige de samme
Videnskaber som Descartes' Ind\-%
% --- p. 42 ---
\setcounter{footnote}{0}%
deling (jfr. Descartes's Brev til den franske Oversætter af principia
philosophiæ); men medens de hos denne danne en enkelt Række, i hvilke
Moralen bliver den afsluttende, men uopnaaelige Videnskab, deles Rækken
hos Spinoza, i Overensstemmelse med Parallelismen mellem de to Attributer,
Mennesket erkjender. Ville vi bringe de opregnede Videnskaber i en
systematisk Orden, da bliver Læren om Erkjendelsen en Indledning, men som
saadan ufuldstændig idet den først kan blive klar efterat Systemets
Grundtanker ere udviklede. Grundvidenskaben bliver Erkjendelsen af Naturen
ɔ: af Substantsen, Gud. Dette bliver Systemets metaphysiske Deel, i
hvilken baade Ethikens og Physikens almindelige Begreber blive udviklede.
Til den slutter sig saa paa den ene Side Ethiken\footnote{Ep. 38: magnam
Ethices partem, quæ, ut cuivis notum, \emph{metaphysica et physica}
fundari debet.} med Politik og Pædagogik, paa den anden Side Mechaniken og
Medicinen. Men efter Spinozas Grundopfattelse vil der ikke i stræng
Forstand kunne blive Tale om en virkelig Sondring af Videnskaberne, og den
af ham opstillede Parallelisme mellem Tænkning og Udstrækning, hvis
Gjennemførelses Vanskelighed netop vilde have viist sig ved det
speciellere Physiske, maatte forhindre at begge Rækker udførtes, da dette
enten vilde have ført til en for stor Sammensmeltning eller Adskillelse.
Det er derfor vistnok ikke tilfældigt, at Spinoza kun har behandlet den
ene Række, og han kan neppe have paatænkt at give Mere af den anden end
»generalia in physicis« (jfr. Ep. 63). Paa et sildigere Punkt ville vi
see, at den af Udgiveren af hans efterladte Skrifter udtalte Formening, at
Spinoza i sin integra philosophia vilde have udviklet »veram motus
natu\emph{ram atque qua ratione a priori tot varietates in materia
deducendæ essent}«, efter Sagens Natur ikke vilde have kunnet gaae i
Opfyldelse.
% PARTIAL LETTERSPACING, verified at 600 dpi: the quotation opens unspaced
% (`veram motus natu-' at the line end) and the spacing begins with the
% second half of the broken word, `r a m', running on to the closing mark.
% Transcribed as printed, not regularised. `integra philosophia' just above
% is NOT spaced here, though it is on p. 41.

I Spinozas Hovedværk, der er overleveret os under Titelen:
\emph{Ethiken}, findes det væsentlige Indhold af de opregnede
Videnskaber. Læren om \emph{Erkjendelsen}, der først
% --- p. 43 ---
kan udvikles fuldstændig efterat det metaphysiske Grundlag, der er
Gjenstand for en umiddelbar Vished, er gaaet forud, er givet i Ethikens
anden Deel, summarisk i prop. 40, schol. 2. \emph{Metaphysikens}
almindelige Deel findes i Ethikens første Bog; de almindelige
\emph{physiske} Sætninger, der ere Forudsætninger for det speciellere
Ethiske, ere givne i Ethikens anden Deel, prop. 13. \emph{Politiken} er i
sit Grundlag givet i Læren om Affecterne, i hvilken ogsaa Principerne for
en \emph{Pædagogik} ligge, og dens Grundbegreber ere bestemte i Eth. IV,
37. schol. 2. Til \emph{Mechaniken} og \emph{Medicinen} findes der, efter
Sagens Natur, kun faa Elementer, men i Udviklingen af Parallelismen mellem
% `og' between the two spaced words is NOT spaced (600 dpi), nor is the
% preceding `Til'. Transcribed as two runs, not regularised.
Aandens Magt til at tænke og Legemets Magt til at handle (jfr. især Eth.
III, 28. dem. V, præfat., Ep. 45) ligger Grundprincipet for Opfattelsen af
dem. --- Kundskaben til Spinozas System hviler altsaa væsentlig paa
Ethiken; fra den maa Fremstillingen gaae ud, og de andre Skrifter tjene
kun som fuldstændiggjørende paa enkelte Punkter.

Af \emph{Udgaver} af Spinozas Skrifter findes der, foruden de
oprindelige, særskilte, 3 samlede: af \emph{Paulus} (Jena 1802---3. 2 Bind
8.), \emph{Gfroerer} (Stuttgart 1830. 8. I denne mangler dog den hebraiske
Grammatik) og \emph{Bruder} (Leipzig 1843---46. 3 Bind. 12.
Stereotypudgave.) Ethiken, tractatus politicus og de emend. intellectus
ere optagne i: Renati des Cartes et Spinozæ dissertationes philosophicæ
ed. C. \emph{Riedel} (Leipzig 1843. 2 Vol. 12.) Af nyere Oversættelser
findes to fuldstændige: en tydsk af B. \emph{Auerbach} (Stuttgart 1841. 5
Bind) og en fransk af E. \emph{Saisset} (Paris 1843. 2 Vol.).
% The initials `C.', `B.', `E.' before the spaced surnames are NOT spaced
% (600 dpi), exactly as `K.' before `F i s c h e r' on p. 13.

Hvor der i denne Afhandling ved Citaterne er angivet Sidetal, referere
disse sig til Originaludgaverne naar ikke en anden Udgave udtrykkelig er
nævnet.

% A short centred rule closes the chapter well below the last line; on this
% page it is printed as a DOUBLE rule, unlike the single rules that close
% the Fortale (p. 14) and the Indledning (p. 26). The rest of p. 43 is
% blank; Cap. 2 opens at the head of p. 44.
\bigskip
\begin{center}\rule{2.5cm}{0.4pt}\end{center}

\newpage

% =====================================================================
%  ANDET CAPITEL  (printed pp. 44--64 = PDF 54--74)
%  Opens at the head of p. 44.
% =====================================================================
\chapter*{Andet Capitel.}
\addcontentsline{toc}{chapter}{Andet Capitel. Spinozismens Grundbegreber}
\markboth{Spinozismens Grundbegreber}{Spinozismens Grundbegreber}

\begin{center}
\textbf{Spinozismens Grundbegreber.}
\end{center}

% --- p. 44 ---
Det, der i Indledningen blev udviklet angaaende Spinozas Forhold til
Descartes, fører umiddelbart til det Punkt, fra hvilket hans System kan
overskues. Idet Substantsbegrebet opfattes i dets Strænghed, bliver det
det altomfattende Begreb, altsaa ogsaa Udgangspunktet, og et
Udgangspunkt, der ikke mere kan forlades. Medens Descartes maatte gaae
ud fra det endelige Jeg, gaaer Spinoza ud fra dette Jegs absolute
Væsen, fra det Uendelige, og kan kun begynde med dette fordi kun det
Uendelige er. Det er altsaa ikke gjennem en Deduction, at dette naaes;
det er Gjenstand for en umiddelbar Vished; i det tænker Tanken kun sig
selv i sin Uendelighed. Herimod strider det ikke, at der hos Spinoza
findes Beviser for Substantsens Tilværelse, Uendelighed og Enhed, og
det ikke blot i principia phil. Cartes., hvor i det Væsentlige den
Cartesiske Bevisførelse gjengives, og Spinozas eiendommelige Opfattelse
kun antydes, men ogsaa i Ethiken (I, 7. 8. 14). Thi denne Bevisførelse
har kun Betydningen af noget Formelt, af en Begrebsudvikling for
Subjectet, ikke af noget Begrundende. Dette vil paa sit Sted blive
klarere naar den menneskelige Tænknings Væsen i dens Forhold til
Substantsen er udviklet.

I. Den angivne Opfattelse af \emph{Substantsen} ligger nu til Grund for
dens Definition: Per \emph{substantiam} intelligo id, quod in se est,
et per se concipitur; hoc est id, cujus conceptus non indiget conceptu
alterius rei, a quo formari debeat. (Eth. I, def. 3, jfr. Ep. 27.).
Substantsen »begribes ved sig selv« (\emph{concipitur per se}) idet
dens Begreb ikke dannes ɔ: udledes af et høiere Begreb (saaledes som f.
Ex. Bevægelsens af Udstrækningens), men selv er det høieste Begreb. Den
er »i sig« (\emph{in se}), ikke i et Andet (saaledes som modus er i
Substantsen).
% `Andet' here shows a weak crossbar on the `e' (reads close to `Andct'
% on the image), the same impression defect as `Vanskeligheden' (p. 50)
% and `Stcd' (p. 71, RESUME-NOTES): sense-reading transcribed, doubt
% logged, not a printer's defect and no sort-level claim made (JBIG2,
% header §6).
Men i den Bestemmelse, at den er i sig selv og begribes ved sig selv,
% --- p. 45 ---
\setcounter{footnote}{0}%
ligger umiddelbart, at den ogsaa er \emph{per se}, at den er
\emph{Aarsag til sig selv}\footnote{De emend. intell. pag. 386
sammenstilles: esse in se, og: causa sui.}; thi en Virknings Begreb
indeslutter Aarsagens Begreb (»effectus cognitio a cognitione causæ
% The opening » before `effectus' has no matching closing « anywhere
% before the parenthesis closes --- printed as a missing quotation
% mark, a structural defect (BATCH-AGENT.md: What this scan cannot
% settle). »« balance in this batch: +1 from this site.
dependet, et eandem involvit. Eth. I, ax. 4); det altsaa, der kun
indeslutter sit eget Begreb, kan ikke være Virkning uden af sig selv.
Men idet den er sin egen Aarsag, er dens Existents ikke tilfældig, men
nødvendig, sat med dens Væsen: Per \emph{causam sui} intelligo id,
cuius essentia involvit existentiam, sive id, cuius natura non potest
concipi nisi existens. (Eth. I, def. 1.). Denne Enhed af Existents og
Væsen er \emph{Evighed}: per æternitatem intelligo ipsam existentiam,
quatenus ex sola rei æternæ definitione necessario sequi concipitur
(Eth. I, def. 8.)\footnote{Æternitas = infinitas et necessitas
existentiæ; Eth. I, 23. dem.}. Og idet Evigheden er den absolute
Affirmation af en Naturs Existents, følger deraf umiddelbart
\emph{Uendeligheden} medens alt Endeligt er ex parte negatio (Eth. I,
8. schol. 1.). Men i Uendelighedens Begreb ligger \emph{Enheden}, baade
som Udelukkelse af \emph{Flerhed}, (thi denne forudsætter ikke blot en
gjensidig Begrændsning, men ogsaa en Adskillelse mellem Væsen og
Existents, og tillige, da den ikke er givet ved Tingens Definition, en
\emph{ydre} Aarsag, jfr. Eth. I, 8 schol. 2), --- og som Udelukkelse af
\emph{Delelighed}\footnote{I disse Bestemmelser af Substantsbegrebet
gjenfinde vi Sporene af den Genesis, gjennem hvilken det er fremgaaet.
Idet Substantsen sættes som Enhed af Væsen og Existents, sættes den som
den fra det Sandselige sondrede Tankes objective Væsen; thi
Sandseligheden adskiller Væsen og Existents, begrændser Forstanden,
denne gjenfinder kun sig selv hvor hin Adskillelse er hævet, hvor den
bekræfter sig selv i Enheden af Tanke og Existents. I Bestemmelsen af
Substantsen som Causalitet gjenfinde vi Consequentsen af, at den fra
det Sanselige sondrede Tanke, forat komme til et Indhold, maa søge
dette i Naturen; thi fra denne er Causalitetsbestemmelsen taget, og
overført paa Tanken. Men det vil vise sig, at den atter berøves sine
eiendommelige Betingelser og eiendommelige Betydning, saa at Naturen
kun gjør sig gjældende som tænkt ɔ: som ophævet Natur.}.
% `Sandselige' / `Sanselige' both occur within this one footnote, as
% printed --- book's own inconsistency, not normalised (cf.
% Consequentser/Conseqventser, RESUME-NOTES).

% --- p. 46 ---
\setcounter{footnote}{0}%
Substantsen er saaledes bestemt som det ved sig selv værende og ved sig
selv begrebne Evige, Uendelige, Ene. Al Virkelighed er sammentrængt i
dens Begreb, i det har Tanken sit Væsens Uendelighed og Erkjendelsens
Enhed.

Udsagte om Substantsen betegne Enhedens og Uendelighedens Begreber det
Samme; thi som blot \emph{Talbestemmelse} har Enheden ikke Anvendelse
paa Substantsen\footnote{Jfr. Ep. 50. 39 ff., princ. ph. Cart. I, 11.
schol., Cog. met. I, 6. Eth. I, 8.}. Men Uendeligheden er ikke opfattet
som Negationen af et positivt Endeligt, men som det negative Endeliges
positive Væsen\footnote{Angaaende Forholdet mellem det Uendelige og det
Ubestemte jfr. Eth. I, 8. schol. 1. II, def. 5. Ep. 29. 41. 64. Cog.
met. I, 6. De int. em. p. 360.}. Som Negationer vilde begge
Bestemmelser føre til en abstract, forskjelsløs Tomhed og Livløshed,
men Tanken søger i dem netop Udtrykket for Fylde. Og i Bestemmelsen:
Aarsag til sig selv, anskuer den Activitet, og med Activiteten er
Forskjelsløsheden hævet. Idet Enheden bestemmes ved Uendelighed og
Activitet, føres vi til de to Begreber, der i det spinozistiske
System, i Forening med Substantsen, udgjøre Grundbegreberne: Begrebet
\emph{Attribut}, og Begrebet \emph{Modus}.

Uendelighed udelukker Begrændsning, thi Begrændsning negerer. Men
Begrændsning finder kun Sted ved hvad der er af samme Art, ikke ved det
Uensartede: Ea res dicitur \emph{in suo genere finita}, quæ alia
eiusdem naturæ terminari potest. Ex. gr. corpus dicitur finitum, quia
aliud semper maius concipimus. At corpus non terminatur cogitatione,
nec cogitatio corpore (Eth. I, def. 2.). Uendeligheden kan altsaa i
sig, uden at ophæves som Uendelighed, indeslutte ligesom en ny
Uendelighed, en Mangfoldighed, der, paa Grund af Artsforskjelligheden,
ikke medfører Begrændsning, og en saadan Mangfoldighed bliver, som
forenelig med Uendeligheden, forenelig med Enheden. Den fordres endog
af Uendeligheden, idet Alt, hvad der udtrykker
% --- p. 47 ---
\setcounter{footnote}{0}%
Væsen og ikke indeslutter Negation, henhører til den (Eth. I, def. 6.
explic.).

II. Ved denne Mangfoldighed, der uden at ophæve Uendeligheden, er sat
indenfor Substansens Enhed, er \emph{Attributets} Begreb givet: per
\emph{attributum} intelligo id, quod intellectus de substantia
percipit tanquam eiusdem essentiam constituens (Eth. I, def. 4). Denne
Definition maa ikke misforstaaes som om Ordene: id, quod intellectus
--- --- percipit, skulde gjøre Attributet til en blot \emph{subjectiv}
Opfattelse\footnote{Til Oplysning af Forholdet mellem Substants og
Attribut kan Brugen af Bestemmelsen \emph{essentia} i Princip. phil.
Cart. og Cog. metaph. tjene. Den sættes paa den ene Side som synonym
med \emph{definitio} (pr. ph. C. II, 15. schol.) og med
\emph{conceptus}, paa den anden Side som synonym med \emph{natura}
(jfr. pr. ph. C. I, 4. ax. 11 not. og def. 9, med II, 15. schol.), med
\emph{vis}, qua in meo esse me conservo (pr. ph. C. I, 7. schol., Cog.
met. II, 1), med \emph{potentia} (pr. ph. C. II, 2. schol.) og med
\emph{perfectio} (Cog. met. I, c. 6), og der tales om \emph{essentia
formalis} ɔ: objectiv (Cog. met. I, 2). Den subjective og objective
Opfattelse er saaledes, efter Spinozismens hele Charakteer, i Enhed.
Efter denne Charakteer maatte ogsaa Substantsen bestemmes baade som
det, der er i sig, og som det, der begribes ved sig.}. Misforstaaelsen
har imidlertid en stor Autoritet at støtte sig paa, jeg maa derfor
udførligere modbevise den. Allerede Definitionens Slutningsord: tanquam
eiusdem essentiam constituens, modsige den subjective Opfattelse,
endnu mere Spinozismens hele Charakteer (jfr. Cap. 5, Læren om
Erkjendelsen) og bestemte Udsagn af Spinoza. Sammenligner man først de
to Steder i Brevene, hvor der gives Definitioner paa Attributet, vil
man finde, at det paa det ene Sted (Ep. 2, skrevet 1661, jfr. Ep. 4)
bestemmes i det Væsentlige paa samme Maade som Ethiken bestemmer
Substantsen: notandum, me per attributum intelligere omne id, quod
concipimus per se et in se, adeo ut ipsius conceptus non involvat
conceptum alterius rei; --- ligeledes, at i dette Brev Definitionen af
Attributet umiddelbart er knyttet til Definitionen af Gud uden at
Substantsbegrebet
% --- p. 48 ---
\setcounter{footnote}{0}%
kommer imellem\footnote{Man seer af Ep. 2 ff. ogsaa, at de 4 første
Sætninger i Ethiken oprindelig have været opstillede som Axiomer, og at
Sætningernes Række da begyndte med den nuværende prop. 5. Ligesom der i
Brugen af Substantsbegrebet har funden en formel Afvigelse Sted fra
Ethiken i dens nuværende Form, saaledes er ogsaa Bestemmelsen accidens
brugt istedenfor Bestemmelsen modus, medens Spinoza sildigere adskilte
dem.}. Paa det andet Sted (Ep. 27, fra 1663) defineres først
Substantsen som det, quod in se est et per se concipitur, hoc est,
cujus conceptus non involvit conceptum alterius rei. Derpaa tilføies
der: \emph{Idem} per \emph{attributum} intelligo, nisi quod attributum
dicatur respectu intellectus, substantiæ certam talem naturam
tribuentis. Paa det Samme vise mange Steder i Ethiken hen. Saaledes
udsiger allerede Definitionen af Gud som den af uendelige mange
Attributer bestaaende Substants (\emph{constantem} infinitis
attributis) den reale Enhed. Ligeledes Udtrykkene i Ethiken I, 15.
schol.: hinc conclusimus, \emph{substantiam extensam unum ex infinitis
Dei attributis} esse. Beviserne for Sætningerne 10, 19 og 20 i Ethikens
første Bog grunde sig paa samme Forudsætning. I Beviset til prop. 10
overføres Substantsens Eiendommelighed, at kunne begribes ved sig selv,
til Attributet som \emph{constituerende} hins Væsen (tanquam eius
essentiam constituens). Substantsens \emph{Evighed} overføres i prop.
19. demonstr. paa Attributet som det, quod divinæ substantiæ essentiam
exprimit, hoc est, id quod ad substantiam
pertinet\footnote{jfr. Henvisningen til prop. 11 i Beviset til prop. 21
i Ethikens første Bog.}. De samme Udtryk bruges i Beviset til prop. 20,
og derfor kan der i Coroll. 2 til denne Sætning sluttes: Deum sive
omnia Dei attributa esse immutabilia\footnote{Hermed kan man
sammenligne princ. ph. C. I. ax. 4, hvor der siges, at Substantsen har
mere Realitet end et Accidens eller en modus, men derimod ikke, at den
har mere Realitet end Attributet.}. Til samme Resultat fører det, naar
der i Eth. II, 40. schol. tales om essentia
% --- p. 49 ---
\setcounter{footnote}{0}%
formalis\footnote{Formalis betyder hos Spinoza og den Tids Philosopher
det Samme som i vor philosophiske Terminologie: objectiv. Hvad vi kalde
subjectivt, kaldtes da objectivt f. Ex. esse objectivum ɔ: Tingenes
Væsen som objectiverede i Forestillingen. Hos nogle af Scholastikerne,
f. Ex Occam, betegner omvendt subjectivt (ɔ: existerende som Subject,
Substants) hvad vi kalde objectivt.} attributorum. Ligeledes viser den
Maade, paa hvilken Spinoza opfatter vor Viden om Attributerne (jfr. Ep.
28. Eth. II, 40 schol. 2), hen paa deres objective Realitet. Og endelig
udsiger Beviserne for Eth. I, 4 og 30\footnote{jfr. Cog. met. I, c. 2.
pr. ph. C. I, 7 schol.} udtrykkelig, at Attributet ikke blot i
Forstanden (intellectus) er ligt med Substantsen.

Attributet\footnote{Ubestemtere bruges Betegnelsen Attribut ikke blot
paa mange Steder i Spinozas tidligere Skrifter og i den hebraiske
Grammatik (f. Ex. C. 33), men ogsaa i Ethiken; saaledes naar res,
aliquid o. desl. kaldes Attributer (Eth. II, 40. Schol. 1), og naar der
(Eth. IV. 37. schol. 2) siges: apparet, iustum et iniustum, ---
notiones esse extrinsecas, non autem attributa, quæ mentis naturam
explicent.} er saaledes væsensligt med Substantsen, det begribes,
ligesom denne, ved sig selv, er evigt og uendeligt. Det er uendeligt
trods Attributernes Flerhed, idet Begrændsningen kun har Betydning
indenfor samme genus. Det er \emph{in suo genere infinitum} idet vi
kunne benægte uendelig mange Attributer om det; det er ikke
\emph{absolute infinitum}, idet der til det absolut Uendeliges Væsen
hører alt det, der udtrykker Væsen og ikke indeslutter Negation (Eth.
I, def. 6 explic.). Dette \emph{ens absolute infinitum}\footnote{ens
absolute indeterminatum. Ep. 41.} er Gud ɔ: substantia constans
infinitis attributis, quorum unumquodque æternam et infinitam
essentiam exprimit (def. 6). I Begrebet: Gud, har Substantsens Begreb
først naaet sit fuldstændige Udtryk. Substantsen er det Uendelige,
først ved det ubetinget Uendelige naaer den altsaa sit Begreb. Den
uendelige Substants, Gud, maa, som indeholdende uendelig Realitet,
bestaae af \emph{uendelig mange} Attributer (Eth. I, 9), der hvert
begribes for sig, men derfor ikke
% --- p. 50 ---
\setcounter{footnote}{0}%
constituere flere Væsener eller flere forskjellige Substantser: Id enim
est de natura substantiæ, ut unumquodque eius attributorum per se
concipiatur, quandoquidem omnia, quæ habet, attributa, simul in ipsa
semper fuerunt, nec unum ab alio produci potuit; sed unumquodque
realitatem sive esse substantiæ exprimit. Longe ergo abest, ut
absurdum sit, uni substantiæ plura attributa tribuere; quin nihil in
natura clarius, quam quod unumquodque ens sub aliquo attributo debeat
concipi, et quo plus realitatis aut esse habeat, eo plura attributa,
quæ et necessitatem sive æternitatem et infinitatem exprimunt, habeat,
et consequenter nihil etiam clarius, quam quod ens absolute infinitum
necessario definiendum (ut def. 6. tradidimus) ens, quod constat
infinitis attributis, quorum unumquodque æternam et infinitam certam
essentiam exprimit (Eth. I, 10, schol.).

Paa dette Punkt opstaaer der imidlertid en væsentlig Vanskelighed. De
forskjellige Attributer begribes ved sig selv, ere i sig selv, uden
Afhængighed af hverandre, uden gjensidig Begrændsning. I det enkelte
Attribut ligger saaledes Intet, der viser hen paa de andre; det er
afsluttet i sig, indeholder en i sin Art uendelig Realitet, og dette: i
sin Art, betegner ingen Sammenligning. Gaaer man ud fra
Attributet\footnote{Fra Spinozas Synspunkt maa man vel sige, at det er
uberettiget, at Tanken sondrer Attributet fra Substantsen, da de
umiddelbart maae sees i Enhed, og at Tanken, idet den opfatter
Substantsen adskilt fra Attributerne, opfatter den under et nyt,
selvmodsigende Attribut, Bestemmelsesløshedens. Men om man nok saa
meget seer Sagen fra hans Synspunkt, bliver Resultatet det samme, og
Vanskeligheden faaer kun en anden Form. I Læren om Erkjendelsen (jfr.
Cap. 5) ville vi see, at idea corporis og idea ideæ udtrykke Eet og
begge ere i Aanden. Med een idea ere saaledes de to Attributer tænkte i
deres Forskjellighed. Men Udstrækningens objective Væsen som
forskjellig fra Tænkningen tænkes ikke, den udtrykker sin actuosa
essentia gjennem en med Tanken parallel Virkning, der bliver
selvstændig og kun ideelt er optaget i Aanden. Her kommer altsaa
Dualismen tilbage.},
% `Vanskeligheden' above: the image shows a weak crossbar on the second
% `e' (reads close to `Vanskcligheden'), already logged in BATCH-AGENT.md
% as a known impression defect, not a printer's fault --- sense-reading
% transcribed, no sort-level claim made (JBIG2, header §6).
kommer man derfor ikke til en Enhed af de
% --- p. 51 ---
flere Attributer i Substantsen; man kommer da i ethvert Tilfælde til at
opfatte Attributet umiddelbart som Substants, men da selvfølgeligt som
Substants med \emph{eet} Attribut. Ikke heller vil man, ved at tage sit
Udgangspunkt fra Tanken om den absolute Substants, med Nødvendighed
komme til Attributernes Flerhed. Det absolut uendelige, absolut
fuldkomne Væsen maa, som indeholdende uendelig Realitet eller Væren,
have uendelige Attributer, slutter Spinoza. Men i Præmisserne ligger,
naar man seer bort fra Spinozas Forudsætning: Attributernes absolute
Forskjellighed, kun Udelukkelse af Negation, og da Negation, strængt
opfattet, bliver liig Begrændsning, og Begrændsning kun har Betydning
indenfor det enkelte Attribut, vilde Slutningen kun føre til Attributet
i dets Uendelighed uden nogen Begrændsning, ikke til uendelig mange
Attributer. Tænker man sig en anden historisk Forudsætning for
Spinozismen, vilde Slutningen kunne have ført til at sætte Tanken som
det eneste Attribut, saaledes, at de andre bleve Former af Tanken,
eller Udstrækningen som det eneste, saa at Tanken blev en Form af den.
Attributernes Selvstændighed i deres Mangfoldighed er en Antagelse, der
hviler paa Spinozismens historiske Forudsætning, paa Dualismen mellem
Tanke og Udstrækning, hvis Enhed fordredes. De to Attributer: Tænkning
og Udstrækning (jfr. Cap. 5) bekræfte sig umiddelbart for den
menneskelige Aand idet denne er corporis idea, og skjøndt der allerede
deri, at Legemet kan blive et ideatum, ligger en Ophævelse af den
abstracte Modsætning, sættes de, i Kraft af hiin Dualisme, som absolut
forskjellige, saaledes, at hvert begribes for sig, og de kun ere i
Enhed idet de udtrykke Substantsen, altsaa idet der abstraheres fra
deres Forskjel, ikke ved sig selv som disse bestemte, saaledes at deres
Forskjel er bevaret. Men idet der er sat en \emph{Flerhed} af
Attributer, maa Substantsens Uendelighed udtrykke
% --- p. 52 ---
\setcounter{footnote}{0}%
sig i Flerhedens, \emph{Antallets Uendelighed}, og Attributerne sættes
som \emph{uendelig mange}\footnote{Forholdet mellem Attributernes
uendelige Antal og den Tohed af Attributer, der udtrykkes i det
menneskelige Væsen, samt den Vanskelighed, der opstaaer paa dette
Punkt, vil blive berørt i Cap. 5.}.

Bestemmelsen af Attributernes Mangfoldighed er altsaa betinget ved
Spinozas historiske Forudsætninger, men Opgaven maatte ligefuldt blive,
at bringe denne Bestemmelse i Harmonie med Substantsens Enhed. I den
Afvisning af Indvendinger, der er givet i det ovenfor citerede Sted,
gaaes der ud fra Forudsætningen af de flere Attributer, og deraf
sluttes til deres Selvstændighed; men Muligheden af deres Enhed trods
Forskjellen paavises ikke. Dette skeer ikke heller ved Theorien om
Parallelismen mellem de Rækker af Modificationer, der henhøre under de
forskjellige Attributer (Eth. II, 7); thi denne Forskjellighed bliver
staaende som en Forudsætning, der ligesaalidt føres tilbage til en
Enhed, som den er fremgaaet af en saadan. Det er ikke tilstrækkeligt
(med Feuerbach) at sige, at netop derved, at Aand og Materie begribes
hver for sig og ved sig, derved at begge ere uafhængige i deres Begreb,
udtrykke de ikke sig selv, men Substantsen, Gud, og, idet de begge ere
Substants, udtrykke de kun dette, at de ere Substants, saa at det
Affirmative i dem kun er Substantsen, rent som saadan, der er
ligegyldig imod om den er Aand eller Materie\footnote{Efter Spinozas
Opfattelse affirmerer Attributets Eiendommelighed sig netop umiddelbart
for Tanken.}, og for hvilken de derfor begge, \emph{forsaavidt} de ere
et \emph{Bestemt}, fra hinanden \emph{Forskjelligt}, og
\emph{forsaavidt} der kun sees paa denne Bestemthed som Bestemthed,
ikke derpaa, at den er og udtrykker Substantsen, kun ere Attributer,
Egenskaber, Momenter, der constituere Guds Idee. --- Paa denne Maade
bliver nemlig Attributets Eiendommelighed opfattet som adskillelig fra
Substantsen, altsaa som et Forsvindende; Substantsen, der bliver
ligegyldig mod Differentsen, bliver enten en Abstraction, en notio
universalis,
% --- p. 53 ---
egentlig en subjectiv Imagination (jfr. Cap. 5), eller Differentsen
bliver noget Subjectivt og Substantsen bestemmelsesløs. Opfattes
Differentsen i denne Uvæsentlighed, kan man ikke engang sige, at
Bestemthederne som Momenter constituere Guds Idee; de blive nemlig kun
noget Subjectivt, denne et Realt.

En nærmere Begrundelse og Udvikling findes imidlertid ikke hos Spinoza.
Han har kun, idet han under Bestræbelsen for at give Substantsbegrebet
Fylde, ved sin Tids Forudsætninger førtes til at søge denne i
Attributernes uendelige Mangfoldighed, søgt en Bestemmelse, der gjorde
det muligt at opfatte de enkelte som uendelige, og denne har han
fundet i det eiendommeligt opfattede Begrændsningens Begreb, og i
Adskillelsen mellem det, der er uendeligt i sin Art og det ubetinget
Uendelige. Han har dernæst ført Beviset for, at \emph{under
Forudsætning} af de uendelig mange Attributer i den ene Substants,
disse maatte begribes i og ved sig selv, gjensidig uafhængige. Men
Flerhedens Optagelse i Enheden, de uafhængige Differentsers
Forenelighed i Substantsen, hvis Væsen kun udtrykkes i Attributerne,
har han ikke paavist, og, som det senere vil vise sig, ikke kunnet
paavise. Dette er en Ufuldstændighed i hans System, som den historiske
Fremstilling har at paapege og belyse, ikke at fjerne.

III. Substantsen, der i sin Enhed indeslutter Attributernes
Mangfoldighed, faaer nu sin videre Bestemmelse ved
\emph{Causalitetsbegrebet}. Substantsen er \emph{causa sui}; dette er
ikke blot et Udtryk for det Negative, at den ikke er bevirket ved et
Andet, men har \emph{positiv} Betydning, idet det bestemmer Guds Væsen
som Activitet. Dette viser sig idet essentia opfattes som Eet med
potentia (Eth. I, 34: potentia Dei, qua ipse et omnia sunt et agunt,
est ipsa ipsius essentia) og som indeholdende vis agendi (Eth. IV, 26
dem.); det fremgaaer ligeledes af Udtrykket: actuosa essentia (Eth. II,
3 schol. jfr. Cog. metaph. II, c. 11), af Opfattelsen af Attributerne:
Tænkning og Udstrækning, som active, af den vigtige Sætning 16 i
Ethikens første Bog
% --- p. 54 ---
\setcounter{footnote}{0}%
(»ex necessitate divinæ naturæ infinita infinitis modis sequi
debent«)\footnote{Actiones, quæ ex ipsius [Dei] naturæ necessitate
consequuntur. Eth. IV. App. Cap. 4.}, af Beviset for prop. 36 i samme
Bog, af Begrebet idea ideæ (jfr. Cap. 5), og af Læren om Affecterne,
der helt igjennem hviler paa Anskuelsen af Activiteten.

Idet Bestemmelsen \emph{causa sui} faaer positiv Betydning, sættes der
i Identiteten af det Virkende og det Bevirkede en bestandig
forsvindende, eller rettere \emph{evigt forsvunden} Differents, og
gjennem denne Differents som Differents naae vi til det tredie
spinozistiske Grundbegreb, Begrebet \emph{modus}\footnote{Som
Synonymer bruges: modi, modificationes, affectiones. Fra modus
adskiller Spinoza accidens (Cog. met. I, c. 1), der hos Descartes er
ensbetydende dermed og af Spinoza selv endnu bruges saaledes i Ep. 4.
Betegnelsen affectio bruges i Cog. met. I, c. 3 som Synonym med
Attribut.}. »Per \emph{modum} intelligo substantiæ affectiones, sive id
quod in alio est, per quod etiam concipitur«\footnote{Modificationes
er: Id, quod in alio est, et quarum conceptus a conceptu rei, in qua
sunt, formatur. Eth. I, 8. schol. 2.}. (Eth. I, def. 5). Idet modus
bestemmes som værende \emph{i et Andet og begreben ved et Andet, er
den bestemt som sat ved et Andet, ved det, i hvilket den er og ved
hvilket den begribes} (Eth. I, 16. Coroll. 1, 25). Bestemt som Virkning
er modus forskjellig fra Substantsen, dens Aarsag, og naar den nu
bestemmes ved et Pluralis (substantiæ \emph{affectiones}), og bringes
sammen med Bestemmelsen: \emph{enkelte Ting} (\emph{res particulares}
nihil sunt nisi Dei attributorum affectiones sive modi, quibus Dei
attributa certo et determinato modo exprimuntur. Eth. I, 25. Coroll.),
saa synes denne Forskjel, idet det af Substantsen Satte og i
Substantsen Værende bestemmes ved Flerhed og Begrændsethed, at blive en
absolut, og at faae tilbagevirkende Kraft til at omstøde Opfattelsen af
Substantsbegrebet.

Bestemmelsen: Enkelthed og Begrændsning fører os til \emph{Problemet om
det Uendeliges Forhold til det Endelige}, og Spørgsmaalet kunde
nærmest synes at maatte
% --- p. 55 ---
stilles saaledes: hvorledes det Uendelige hos Spinoza frembringer en
Endelighed, hvorledes \emph{Overgangen} finder Sted fra hint til
dette? Men i denne Form har Problemet allerede faaet en skjæv
Stilling, idet der, med Forglemmelse af Systemets Grundtanke,
tilstaaes det Endelige Realitet. Fastholdes Grundbegrebet, Systemets
Udgangspunkt, saa maa Opgaven tværtimod just blive den: at begribe det
Uendeliges Forhold til det Endelige saaledes, at det sees, at der
ikke kan være Tale om en Overgang fra hint til dette, at det
\emph{Endelige ikke er som Endeligt}, at Intet er uden Substantsen,
og at selve Erkjendelsen af det Endeliges Intethed ikke vorder ---
altsaa ikke er underkastet Endelighedens Form; --- \emph{men er
evigt}\footnote{Med Hensyn til Problemet om de endelige Ting jfr. Eth.
II, 8. Cor. 9. Cor. Ep. 58. Pr. phil. Cartes. I, 5. 7. Cog. metaph. I,
c. 2. II, c. 4. 7).}.
% The note's closing `7).' has no matching open paren -- as printed, both
% witnesses agree; a terse citation string, not flagged as a defect.

Det kunde synes, at Bestemmelsen af Substantsen som
\emph{Causalitet}, af Gud som causa omnium rerum (Eth. I, 16. Cor. 1)
maatte afgjøre Problemet om det Endeliges Væren paa en modsat Maade.
Den \emph{reale} Causalitet indeslutter et Tidsforhold, og med Tiden
er Endelighed, om end ikke selvstændig bestaaende endelige Ting,
givet. Men denne Consequents undviger Spinoza dels ved Begrebet causa
sui, i hvilket Causaliteten ligesom er bøiet tilbage i sig selv; dels
ved at statuere, at Gud er alle Tings Aarsag \emph{eo sensu} (i samme
Forstand) i hvilken han kaldes causa sui (Eth. I, 25. schol.), hvorfor
han maa bestemmes ikke som \emph{causa transiens} ɔ: som virkende over
paa et Andet, men som \emph{causa immanens}\footnote{„agentem sive
causam immanentem.“ Gr. hebr. c. XII. --- I Index til Opp. posth.
sammenstilles Bestemmelserne causa immanens og \emph{natura
naturans} (see nedenfor).}, som iboende Aarsag (Eth. I, 18, jfr. Ep.
21); dels endelig indirecte ved at udvikle den Forskjellighed, der
efter den almindelige Opfattelse af Causalitetsforholdet bliver
mellem Aarsag og Virkning, baade med Hensyn til Væsen og Existents
(Eth. I, 17. schol.)\footnote{At Argumentationen paa dette Sted er
polemisk mod den almindelige Opfattelse af Causalitetsbegrebet, seer
man af andre Steder f. Ex. „Gud har ogsaa formaliter (objectivt)
Noget tilfælles med de endelige Ting.“ Ep. 4. jfr. Eth. I, 25. Coroll.
(Prop. 10 i Ethikens 2den Bog, der angaaer den enkelte modus, faaer
ikke Anvendelse her.)} ---
% --- p. 56 ---
\setcounter{footnote}{0}%
hvilket, naar de endelige Ting opfattedes som selvstændig bestaaende,
vilde føre til Ophævelsen af Substantsbegrebet, og, anvendt paa
Begrebet: causa sui, til Absurditeter.

Betragte vi nærmere Causalitetsbegrebets Betydning hos Spinoza, see vi
ogsaa, at Eiendommeligheden i dets Opfattelse afskjærer hiin
Consequents. \emph{Causaliteten falder for Spinoza sammen med den
logiske Følge}\footnote{De logiske Følger kan Spinoza derfor ogsaa
betegne som actiones. Ogsaa ved denne Opfattelse bekræfter det sig, at
den spinozistiske Substants er Forstandens objectiverede Væsen.
Derfor sættes ogsaa i Læren om Erkjendelsen agere som liig med
Intelligere (Eth. IV, 24. dem.).}: »Alt udgaaer med Nødvendighed
eller følger bestandig med samme Nødvendighed af Guds uendelige
Natur, paa samme Maade, paa hvilken det fra Evighed til Evighed
følger af Triangelens Natur at dens tre Vinkler ere lige med to
rette« (Eth. I, 17 schol.); --- og i den ovenfor citerede 16de
Sætning i 1ste Bog føres Beviset deraf, at der af enhver Tings
Definition (ɔ: af Tingens Væsen) med Nødvendighed følge flere
Egenskaber; idet der altsaa af den guddommelige Naturs absolut
uendelige Attributer, af hvilke ethvert udtrykker en i sin Art
uendelig Væsenhed, maa følge Uendeligt i uendelige Modificationer.

Ved denne logiske Opfattelse af Causaliteten er Tidsbestemmelsen
fjernet\footnote{Guds Prioritet er Causalitetsprioritet; Pr. phil.
Cart. I, 12. Coroll 4.}; Forholdet bliver indenfor Evighedens Element,
der ikke bestemmes ved Tid, Maal eller Tal, Endelighedens og
Imaginationens Bestemmelser (Ep. 29).

Medens saaledes Causalitetsbegrebet, efter Spinozas Opfattelse, ikke
førte til Endelighed, udelukkes denne ligeledes naar vi i
Causaliteten urgere Substantsbegrebet som den uendelige causa sui.
Causaliteten er i Enhed med Substantsens Væsen, gjennem den
realiserer dette sig selv. Men
% --- p. 57 ---
\setcounter{footnote}{0}%
idet den causale Substantses Virkning er Substantsen selv, er den
uendelige Causalitet udtømt i den uendelige Virkning; der bliver ikke
Plads til en Bevirken af et Endeligt, og i Substantsens Væsen findes
der intet Tilknytningspunkt for dette. Hvad der følger af Substantsens
Væsen, følger med Nødvendighed; men Nødvendighed indeslutter
Uendelighed\footnote{jfr. Eth. I, 7, demonstr. med I, 8. schol. 1.}.
Og idet Substantsen bestemmes som det uendelige, absolut reelle
Væsen, der er Enhed af Væsen og Existents, bliver der ligesaalidt
Plads for en fra det adskilt virkelig Existents som for et fra den
adskilt Væsen. Om ogsaa Modificationens Existents, efter Spinozas
Sprogbrug, ikke bestemmes ved Evighed, men ved
duratio\footnote{\emph{Duratio est indefinita existendi continuatio}
(Eth. II. def. 5). Eth. II, 45. schol.: durationem, hoc est
existentiam, quatenus abstracte concipitur, et tanquam quædam
quantitatis species. Med Hensyn til Forholdet mellem æternitas og
duratio jfr. man. Eth. I, def. 8. explic., prop. 23. dem., Cog. met.
I, c. 4. II, c. 1. 2; med Hensyn til Forholdet mellem duratio og
essentia Eth. I, 24. dem., IV. præf.}, der er Udtrykket for den
Existents, der ikke følger af Tingens Væsen, men af Aarsagens Natur,
saa bliver dog duratio, paa Grund af Aarsagens Uendelighed, at tænke
som uendelig naar Virkningen sees i sin Sandhed ɔ: i sin Enhed med
Substantsen, og først naar den ved en Abstraction skilles fra denne,
komme Endelighedens Bestemmelser frem (Ep. 29).

Til det samme Resultat fører Spinozas Opfattelse af disse
Endeligheds-Bestemmelser: Begrændsning, Flerhed, Delelighed, Tid, Tal
og Maal. \emph{Begrændsningen} opfattes kun som Negation (Ep. 50), den
bliver altsaa utænkelig naar man gaaer ud fra Substantsen
(Attributet); thi da Begrændsning kun kan skee ved det Væsenslige,
maatte den skee ved Attributet selv, og dette som det Begrændsende
maatte opfattes som det Begrændsedes, altsaa som sin egen, Negation
(Ep. 50). Derfor er Begrændsning udelukket ved Substantsens Begreb; og
naar den sees fra et hypothetisk
% --- p. 58 ---
antaget Endeligts Synspunkt, bliver den, som noget blot Negativt, et
Forsvindende.

Bestemmelsen Flerhed ligger ikke i Tingens Definition, men
forudsætter en ydre Aarsag (jfr. Eth. I, 14 schol. Ep. 41). Forfølger
man denne Sætning tilbage, da vil, idet Aarsagernes Række føres
tilbage til den første Aarsag, Substantsens Enhed udelukke Flerhed,
altsaa Endelighed. --- Bestemmelserne Delelighed, Tid, Tal, Maal, der
opstaae, idet Varighed og Quantitet ved en Abstraction skilles fra
Substantsen og fra »den evige Orden« (jfr. Ep. 29), føre til
Modsigelser, der ikke kunne løses: en uendelig Række, der ingen Enhed
kan danne, en Uendelighed af Dele, der ikke kan naae til at udgjøre
en Helhed, en \emph{større og mindre} Uendelighed. Men disse
Modsigelser hæves idet Forestillingen om de endelige, enkelte Ting
hæves; kun den Realitet, der tillægges disse som endelige, fremkalder
hine.

Men medens saaledes Betragtningen af alle disse Begreber fører bort
fra et existerende Endeligt, kunde det synes, at Spinoza selv har
gjort et Forsøg paa at udlede det Endelige af det Uendelige ved
Sætningerne 21--23 i Ethikens første Bog. De lyde: »Alt hvad der
følger af et af Guds Attributers absolute Natur, har altid maattet
existere og existere som uendeligt, eller er evigt og uendeligt ved
dette Attribut.« --- »Hvad der følger af et af Guds Attributer,
forsaavidtsom det er modificeret ved en saadan Modification, der
existerer med Nødvendighed og som uendelig ved Attributet, maa ogsaa
existere med Nødvendighed og som uendeligt.« --- »Enhver modus, der
existerer med Nødvendighed og som uendelig, har med Nødvendighed
maattet følge enten af et af Guds Attributers absolute Natur, eller af
et Attribut, modificeret ved en Modification, der baade existerer med
Nødvendighed og som uendelig.« --- Der er her givet en Art Gradation:
\emph{Attributet; det, der er umiddelbart frembragt af Gud} (immediate
a Deo productum), \emph{og det, der er frembragt mediante infinita
quadam modificatione} (jfr. Ep. 65 og Spinozas Svar i
% --- p. 59 ---
Ep. 66, hvor der som Exempler paa den anden Art anføres, under
Tænkningens Attribut: \emph{intellectus absolute infinitus}, under
Udstrækningens: \emph{Bevægelse og Hvile}; som Exempel paa den tredie
Art: \emph{facies totius universi, quæ}, quamvis infinitis modis
variet, manet semper eadem). Men Gradationen bliver indenfor
Uendeligheden, og sammenligner man med disse Sætninger den 28de
Sætning i samme Bog: »Enhver enkelt Ting, eller enhver Ting, der er
endelig og har en begrændset Existents, kan ikke existere eller
bestemmes til Virksomhed uden ved en anden Aarsag, der ligeledes er
endelig og har en begrændset Existents; og denne Aarsag kan igjen
ikke existere eller bestemmes til Virksomhed uden at den bestemmes
til Existents og Virksomhed ved en anden, der ligeledes er endelig og
har en begrændset Existents, og saaledes i det Uendelige« --- saa
bliver ikke blot Udledningen ikke ført fra det Uendelige til det
Endelige; \emph{men en saadan Overgang sættes udtrykkeligt som
umulig}, idet med det Endelige den i det Uendelige gaaende aldrig
afsluttede Række er givet. Gaaer man altsaa ud fra det Uendelige,
kommer man kun til Uendeligt, aldrig til det Endelige; gaaer man ---
ved en Hypothese --- ud fra det Endelige, kommer man, i den uendelige
Række, kun til Endeligt, aldrig til det
Uendelige\footnote{Til den samme Consequents fører Spinozas
Opfattelse af de forskjellige Erkjendelsesarters Forhold til
hverandre (jfr. navnlig Eth. V, 28). Den klare og tydelige
Erkjendelse, der begriber under Evighedens og Nødvendighedens Form
(sub specie æterni), kan ikke opstaae af den uadæquate Erkjendelse,
hvis Gjenstand er det imaginerede Endelige, men kun af en anden klar
og tydelig Erkjendelse. Her er den samme udtrykkelige Benægtelse af
Overgangen, hvilket vil blive indlysende naar Erkjendelsens Natur og
Arter er udviklet (jfr. Cap. 5 og Cap. 8.).}.

Ved Bestemmelsen \emph{modus} ere vi saaledes endnu ikke traadte ud af
Uendeligheden. Allerede i Betegnelsen: modus eller modificatio, er
givet en Væren, ikke i sig, men i et Andet, hvis Affection den er.
Dette Andet er den uendelige immanente Aarsag, Substantsen. Og ved
nærmere
% --- p. 60 ---
\setcounter{footnote}{0}%
Bestemmelser sikkres det nu, at ikke modus, skjøndt værende i det
Uendelige, dog ved sin Uensartethed med dette skal blive udenfor det
Uendelige og altsaa ophæve dette som saadant. De Mellembestemmelser,
ved hvilke Spinoza søger at hæve Uforeneligheden indenfor det
Uendelige, ere navnlig tre: Begrebet \emph{uendelige modi},
Opfattelsen af modus under Evighedens og Nødvendighedens Form (sub
specie æternitatis et necessitatis), og endelig Opfattelsen af modi
som en Helhed, som \emph{universelt Individuum, som natura
naturata}.

Stiller man Spørgsmaalet saaledes: hvorledes kunne Modificationer ɔ:
Bestemmelser, Begrændsninger tænkes i det Uendelige? saa skyder
Begrebet; \emph{uendelige modi}, Besvarelsen et Skridt tilbage. Idet
der statueres uendelige modi faaer Spørgsmaalet den Form: hvorledes
kunne modi tænkes som uendelige? Svaret herpaa forberedes ved
Opfattelsen af modus under Evighedens Form, forsøges afgjørende givet
ved Opfattelsen af \emph{modi som Helhed}. Idet Fornuften
\emph{tænker} modus, opfatter den den ikke under Tidens og Rummets
Bestemmelser, --- der ere Imaginationens Bestemmelser, men under
Evighedens og Nødvendighedens Form, i Gud ɔ: efter dens
Sandhed.\footnote{For Spinoza ere Udtrykkene: res in se spectata, ---
vere spectata, --- ad Deum relata, ensbetydende.} Men idet den
opfattes saaledes, opfattes den ikke som enkelt og endelig, men i sin
umiddelbare Sammenhæng med sin evige uendelige Aarsag; dens Definition
bliver Substantsens Definition (Ep. 29), dens Existents væsentlig,
nødvendig, uendelig Existents (Eth. II, 45. schol.). Nødvendig
Existents strider mod Endelighedens Væsen; gjennem Opfattelsen af
modi som nødvendige og evige føres der derfor til Opfattelsen af
Modificationernes Totalitet som \emph{uendelig Enhed}. Begrændsning
er Negation, ikke noget Positivt; det, der, naar et Endeligt
forudsættes, i Forestillingen adskiller modi fra hverandre, er, seet
i Tanken, et Forsvindende; men idet dette Negative hæves, er den
afsluttede Enhed sat. Dette
% --- p. 61 ---
\setcounter{footnote}{0}%
giver Bestemmelsen af det \emph{universelle Individuum} (Eth. II, 13.
Lemma 7, schol.; jfr. Cap. 4). Dette \emph{universelle} Individuum er
\emph{natura naturata}\footnote{\emph{Natura naturata} (Udtrykket er
allerede brugt af Giordano Bruno og Hobbes) bestemmer Spinoza (Eth. I,
29 schol.) ved Modsætningen til \emph{natura naturans}. Denne er: id,
quod in se est et per se concipitur, sive talia substantiæ attributa,
quæ æternam et infinitam essentiam exprimunt; hoc est --- Deus,
quatenus ut causa libera consideratur. --- \emph{Natura naturata er:}
id omne, quod ex necessitate Dei naturæ sive uniuscuiusque Dei
attributorum sequitur, hoc est, omnes Dei attributorum modi, quatenus
considerantur ut res, quæ in Deo sunt et quæ sine Deo nec esse nec
concipi possunt.} som Helhed (Cog. metaph. II. c. 7. 9.). Man kan om
det ikke sige, at det er sammensat, man kan ikke tillægge det Dele;
thi Dele vilde forudsætte en bestaaende Begrændsning, Sammensætning
Begrændsningens Realitet; men Grændsen er her netop sat som det
Ikke-Reale.

Ved disse Mellembestemmelser er Uendeligheden hævdet, Endeligheden og
dens Former udelukkede fra Substantsen. Substantsen er som causa sui
\emph{Enheden af natura naturans og natura naturata}. Som Aarsag er
den naturans, som Virkning naturata; men i Substantsen, hvis Væsen er
Evighed, er denne Forskjel evig forsvunden: Begrændsning, Endelighed,
Tid, Tal, Maal, er ingen Virkelighed, kun et Skin, der end ikke kan
siges at forsvinde, men, seet i sin Sandhed, \emph{kun at være
forsvundet}.

Uendelighedens Interesse er altsaa hævdet; men den anden
% DOUBTFUL READING, not a defect. The scan shows `Interesre'; `Interesse'
% is set here. An earlier draft logged it as a probable broken or wrong
% sort on the ground that s/r is not an n/u pair and so falls outside the
% standing rule. That was too narrow a reading of the rule: the JBIG2
% argument (header §6, JBIG2-TEST.md) is not about n and u in particular
% but about the file placing dictionary symbols instead of impressions, so
% it voids EVERY single-sort claim, s/r included. Both OCR witnesses read
% the same reconstructed image and so do not corroborate each other.
Interesse, der bragte Bestemmelsen modus frem: den, at opfatte
Substantsen som det i sig uendeligt Fyldige, ikke i en hvilende Fylde,
men i en evig sig
selv frembringende, vel ikke saaledes, at der i dens Element skulde
kunne blive Tale om en Udvikling, men dog om en
Activitet,\footnote{Derfor kan Spinoza tale om en \emph{infinitus
amor}, quo Deus seipsum amat (Eth. V, 36). For dette Begreb er
Anskuelsen af Activiteten Forudsætningen. (Jfr. Cap. 6 og Cap. 8.)}
--- er ikke tilfredsstillet, og denne Interesse, der er
% --- p. 62 ---
\setcounter{footnote}{0}%
Tænkningens, Philosophiens egen Interesse, bringer Inconsequentser
frem hos Spinoza, som ikke kunne fjernes. I den bestemmelsesløse
Uendelighed bekræfter Aanden sit abstract-almindelige Væsen, benægter
sig i sin Enkelthed; men dens Maal er Selvbekræftelse. Det, at
\emph{philosophere}, har kun Realitet ved Enkeltheden, kun derved, at
den menneskelige Tanke ikke umiddelbart smelter sammen med den
uendelige. Fastholdes Uendeligheden i hiin abstracte Form, bliver
ogsaa det Endeliges Væren i Forestillingen, om den end nok saa
afgjørende sættes som Skin, som væsenløs,
ubegribelig.\footnote{Det er denne Modsigelse, paa hvilken allerede
Eleaterne stødte an. Idet al Ikke-Væren udelukkes, saaat kun Væren er,
og Væren nu som det Ene omfatter Tænkningen, bliver det Ikke-Værendes
Tilværelse i Tænkningen som Skin uforklarlig, netop derved at Tanken
er sat som Eet med Væren.}

Ved at sætte Activiteten, og ved af den at lade følge Uendeligt i
uendelige Modificationer søgte Spinoza at undflye Abstractionens
Tomhed, og at give Tanken et Indhold og Tænkningen Mulighed. Men
Fylden forsvandt i samme Øieblik, den sattes, den blev til et
flygtigt Skin, idet det, der skulde give Fylden, blev et Skin. De
endelige Ting kunde kun hypothetisk sættes, som et Skin, der kaster
sin Reflex i sig selv, thi de existere kun i den menneskelige
Forestilling, der selv, som endelig, er et Forsvindende. Intethed
afspeiler sig i Intethed. Kun idet Grændsen faaer en Realitet og
derved Modificationerne Bestaaen, kan Substantsen faae et Indhold, og
derfor maa en saadan Betydning af Grændsen komme frem hos Spinoza,
netop ved hiin indre Strid i Tænkningen, og om den end atter
bortvises, hævder den sig dog en Art Gyldighed. Der maa søges
Bestemmelser for det, der ikke umiddelbar gaaer op i Substantsen. Det
kan, om det end bestemmes negativt, eftersom Substantsen er bestemt
som uendelig Virkelighed, ikke rent bringes bort af Tanken, thi denne
vilde derved negere sig selv. Idet Substantsen tænkes, har Tanken
ligesom
% --- p. 63 ---
\setcounter{footnote}{0}%
Erindringen om Endeligheden; denne maa atter tvinge sig frem og
derved opstaae Inconsequentserne.

Disse har derfor Spinoza ikke heller kunnet undgaae. Det er allerede
uconsequent at beskjæftige Tanken med Endeligheden i dens
tilsyneladende Bestaaen; thi ved selve denne Beskjæftigelse med det
faaer det Endelige en Betydning, det efter Grundbegrebene ikke kunde
faae.\footnote{Dog kommer her det Dialektiske atter frem, at denne
Betydning, som støttende sig paa Tiden, ved Tidens Forsvinden som et
Ikke-Realt, atter ophører at være Betydning.} Men bestemtere træder
Inconsequentsen frem naar der tales om en Forskjellighed i det
Endelige; thi den gamle aristoteliske Sætning: \textit{τοῦ μὴ ὄντος
οὐκ ἐστιν εἴδη}, beholder bestandig sin Gyldighed. % Greek, printed in
% an accented italic Greek fount (RECON.md §4); the only \textit{} in
% the book -- not a Latin italicisation, so it does not conflict with
% the plain-roman-Latin rule. Read from the ABBYY witness at 600 dpi;
% tesseract -l dan does not recognise Greek at all.
Forskjelligheden sættes som en Gradforskjel i det Endeliges Væren; men
en saadan Gradforskjel er, om man end ved at fastholde Sætningen, at
al Begrændsning er Negation, tvinger den tilbage til at blive blot
quantitativ, alligevel tilstrækkelig til, ikke rent at lade det
Endelige forsvinde. Det er ligeledes en Inconsequents, naar der tales
om: Ordenen i Naturen, Aarsagernes Række, ordo et constitutio totius
naturæ, series et ordo causarum, der adskilles fra leges naturæ
æternæ; thi i denne Række hævde de enkelte Ting en Betydning som
enkelte. Det er en Inconsequents naar der af de virksomme evige Love
udledes som Følge et uendelig Mangfoldigt (Eth. I, 16). Det er en
Inconsequents, naar der om natura naturata, der opfattes som Enhed,
siges, at den er een »skjøndt dens Dele afvexle paa uendelige
Maader« (quamvis partes eius infinitis modis varient); thi om og
Begrændsning og Forandring sættes som bestandig hævede, sættes de dog
atter for at kunne hæves. Det er en Inconsequents, naar der tales om
res particulares æternæ, om æterni cogitandi modi, der bestemmes som
Dele af infinitus intellectus; thi ved Evigheden er den
Begrændsethed, der ligger i Begreberne: Enkelthed og Deel, hævet. Det
er en Inconsequents naar
% --- p. 64 ---
\setcounter{footnote}{0}%
Spinoza tillægger modi \emph{positiv} Betydning for Erkjendelsen af
Gud (Eth. V, 24); thi skjøndt dette ikke er at forstaae saaledes, at
Erfaringen skulde supplere en oprindelig abstract Gudsidee, men at den
os evig aabenbare Gudsidee kommer til Bevidsthed i sin Fylde og gjør
den forsvindende ufuldkomne Viden til et Skin, saa faaer dog, idet der
tillægges modi som modi positiv Betydning for dette fyldige Begreb,
det, der gjør dem til modi, Grændsen, positiv Realitet. Det er
endelig en Inconsequents at tale om en Gudserkjendelse og en deraf
fremgaaende amor intellectualis i den som en Modus, om end som en
evig Modus, bestemte menneskelige Aand; thi i denne Bestemmelse er det
Uensartede, Begrændsethed og Uendelighed, sammenknyttet (jfr. Cap. 9).

Disse Modsigelser blive saaledes tilbage i Spinozas Lære om
Modificationerne. De ere ikke tilfældige, men nødvendige Følger af,
at den menneskelige Tanke tænker Substantsbegrebet. Den
Inconsequents, at tillægge det Endelige Betydning \emph{om end kun
for atter at hæve den}, er Betingelsen for, at der kan philosopheres,
at ikke Tanken strax standser i Abstractionens
Ubevægelighed.\footnote{Hvormeget dette Problem har beskjæftiget
Spinoza, og hvor tydelig han var sig bevidst, i hvilken Grad hans
Opfattelse af det Endelige maatte være en Anstødssteen for hans
Samtid, seer man overalt i hans Skrifter, og navnlig i den
undvigende Maade, paa hvilken han endnu i sine sidste Leveaar, og
ligeoverfor den af hans Venner, der synes bedst at have forstaaet ham,
udtaler sig derom (jfr. Ep. 70--72). Ligesom for Spinoza Substantsen
er det reelle Udgangspunkt, der ikke mere forlades, saaledes er det
Endelige et hypothetisk Udgangspunkt, der bestandig kun forlades; men
netop heri ligger en Modsigelse.}

% A short centred rule closes the chapter well below the last line; on
% this page it is again printed as a DOUBLE rule, exactly like the one
% closing Cap. 1 on p. 43 -- and unlike the single rules on pp. 14 and
% 26. So the double rule recurs at chapter ends within this half of the
% book; the single rule marks only the Fortale and the Indledning. The
% rest of p. 64 is blank; Cap. 3 opens at the head of p. 65.
\bigskip
\begin{center}\rule{2.5cm}{0.4pt}\end{center}

\newpage

% =====================================================================
%  TREDIE CAPITEL  (printed pp. 65--70 = PDF 75--80)
%  Opens at the head of p. 65.
%  FJERDE CAPITEL  (printed pp. 71--81) opens PART WAY DOWN p. 71, after a
%  short centred ornamental rule closing Cap. 3. Its head, the rule and the
%  transition are therefore inside the pp. 65--81 fragment, not here.
% =====================================================================
\chapter*{Tredie Capitel.}
\addcontentsline{toc}{chapter}{Tredie Capitel. De nærmeste Consequentser af Spinozas Opfattelse af Substantsbegrebet}
\markboth{De nærmeste Consequentser}{De nærmeste Consequentser}

\begin{center}
\textbf{De nærmeste Consequentser af Spinozas Opfattelse af\\
Substantsbegrebet.}
\end{center}

% --- p. 65 ---
\setcounter{footnote}{0}%
Det Spinozistiske Substantsbegrebs constituerende Bestemmelser:
\emph{Uendeligheden og Causaliteten}, maae blive de afgjørende for
Systemets videre Udfoldelse. Substantsbegrebet var jo for Spinoza
det altomfattende Begreb, dets Grundpræg maa derfor gjenfindes i
enhver af hans Tanker. Dette vil i de følgende Capitler vise sig
ved de speciellere Udviklinger; paa dette Punkt vil det blive
paavist med Hensyn til de almindelige metaphysiske Forhold.

Vi saae, at \emph{Substantsbegrebet} i sin Fuldstændighed faldt sammen
med \emph{Gudsbegrebet}, og at \emph{Gud} bestemtes som det absolut uendelige
Væsen. Idet Uendeligheden nu strængt fastholdes, udelukkes først
og fremmest fra Gudsbegrebet enhver Forestilling, der paa
anthropomorphistisk Viis kunde endeliggjøre det. Derfor udelukkes
fra det Bestemmelserne: \emph{Forstand, Villie, Attraa, Kjærlighed}\footnote{Kjærligheden
som passio udelukkes naturligviis fra Gud. Derimod kan der i
Ethikens 5te Bog tales om en activ amor, quo Deus se ipsum amat,
og af hvilken Menneskets amor intellectualis Dei er „en Deel“.
Denne Kjærlighed bliver at parallelisere med infinitus intellectus
(jfr. Cap. 5), idet den er idea Dei seet fra en anden Side.} o.
s. v. fordi de ikke udtrykke det absolut uendelige Tænkningens
Attribut, men kun en bestemt Modification af det (jfr. Cap. 5).
Derfor udelukkes Christendommens Grundforestilling om Guds
Menneskevordelse, da det for Spinozas Opfattelse er ligesaa
selvmodsigende, at Gud, den uendelige Substants, skulde have
antaget Menneskets endelige Natur, som at Cirklen skulde have
antaget Quadratets Natur (Ep. 21).

% --- p. 66 ---
\setcounter{footnote}{0}%
Det med Nødvendighed existerende active Uendelige (causa sui) maa
virke efter sit Væsens Nødvendighed, idet der udenfor det Intet
gives, der kunde bestemme det. Denne Nødvendighed som indre, er
fri Nødvendighed (libera necessitas, Ep. 62),\footnote{Idet denne
Nødvendighed bliver en indre, selve Guds Væsen som udtrykkende sig
i Virksomhed, bliver den ikke Fatalisme, da denne skiller Fatum
fra Gud (Ep. 23. 49. 62.).} Frihed, Gud er causa libera (Eth. I,
def. 7). Som Eet mod Nødvendigheden er denne Frihed ikke et
liberum decretum, ikke Indifferents (jfr. Ep. 60); Guds Magt er
ikke som Kongernes Magt, Vilkaarlighed, den er bunden ved hans
eget Væsens Lov (tr. polit. II, 7). Det Forhold til Tiden, den
Adskillelse mellem Væsen, Forstand og Villie, som et liberum
arbitrium forudsætter, strider mod Evigheden, ved hvilken Guds
Væsen bestemmes. Væsenet er Magten, den nødvendige Existents
giver den nødvendige Activitet. Men Nødvendighedens Begreb
skjærpes ved det, der var Spinozas Tankes dybere Grund.
Substantsen er Forstandens objectiverede Væsen; men dette Væsen
maa, efter Spinozas historiske Forudsætning, bestemmes ved det
abstracte Udtryk for det, hvori Forstanden har sit Indhold ɔ: ved
Naturens abstracte Begreb, og derfor bliver den nødvendige
Causalitet bestemt saaledes, at \emph{Øiemedsbegrebet} udelukkes.
Benægtelsen af dette Begreb, der er afgjørende for Systemets
Charakteer, motiveres vel ud fra Begrebet om Guds Fuldkommenhed
(Uendelighedens Begreb), men hvad der er Polemikens dybere Grund,
og hvad der betinger den eensidige Opfattelse af Begrebet
Øiemed,\footnote{I den Forestilling, at Gud skulde handle sub
ratione boni, seer Spinoza kun en Underkastelse af Gud under et
Fatum, idet den „sætter Noget udenfor Gud, hvilket Gud i sin
Handlen har for sig som et Forbillede og stræber efter som et
Maal“ (Eth. I, 33. schol. 2).} er Naturbegrebets ubevidste
Indvirkning paa Spinozas Tanke.

I Anhanget til Ethikens 1ste Bog udvikler Spinoza dette Begrebs
psychologiske Oprindelse. Forestillingen om
% --- p. 67 ---
\setcounter{footnote}{0}%
Gud som handlende efter et Øiemed, siger han, er kun begrundet i
menneskelige Fordomme, der opstaae derved, at Menneskene fødes
uden Kundskab til Tingenes Aarsager, men med Attraa efter det for
dem Nyttige og med Bevidsthed om denne Attraa. Bevidstheden om
Attraaen, forbunden med Ubevidstheden om dens Grunde, fremkalder
Indbildningen om den menneskelige \emph{Frihed}, og gjennem denne
Indbildning forvandles Bevidstheden om Attraaen til en
Forestilling om en Handlen efter et Øiemed, og denne Forestilling
behersker nu, som afsluttende, Opfattelsen af Tingene, og idet
disse betragtes fra Menneskets Standpunkt, og tillige som
uafhængige af Mennesket, overføres den paa det Guddommelige,
trods alle de absurde Consequentser og trods alle de Modsigelser,
der opstaae ved denne Forestilling. Men Overførelsen af
Øiemedsbegrebet paa Gud vender op og ned paa Naturen. Thi det,
der i Virkeligheden er Aarsag, betragtes som Virkning og
omvendt; det, der efter sit Væsen er det Tidligere, gjøres til
det Sildigere, det, der er det Høieste og Fuldkomneste, til det
Ufuldkomneste,\footnote{Jo mere umiddelbart Noget fremgik af Gud,
desto fuldkomnere maatte det sættes at være, medens efter hiin
Forestilling omvendt det Sidste blev det Fuldkomneste.} og
endelig, --- hvad der er det Afgjørende, --- tillægges der derved
Gud en Ufuldkommenhed; thi naar han handler for et Øiemeds Skyld,
attraaer han nødvendigviis Noget, som han ikke er i Besiddelse
af. (Eth. I, 36. schol. IV, præf.).

Udelukkelsen af Øiemedsbegrebet mødes med Udelukkelsen af de
anthropomorphistiske Bestemmelser i at fjerne de Forestillinger,
der ligge til Grund for Anskuelsen af en \emph{personlig} Gud. I dennes
Sted træder \emph{Naturmagtens Nødvendighed}, der ogsaa bliver
Tænkningens Form, og dette Nødvendighedsbegrebs Betydning for
hele Spinozas System udtaler han selv naar han i Ep. 21 kalder
det præcipuum fundamentum eorum omnium, quæ in tractatu, quem
edere destinaveram (Ethiken), habentur, og i sin sammenfattende
% --- p. 68 ---
\setcounter{footnote}{0}%
Oversigt over Ethikens første Bog, (i Scholiet til prop. 36)
sætter den Bestemmelse: quod [Deus] necessario existat, der
suppleres ved: quod ex sola suæ naturæ necessitate sit et agat,
foran alle de øvrige.\footnote{Eth. IV, præf: æternum illud et
infinitum ens, quod Deum seu naturam appellamus, eadem qua
existit necessitate agit.}

Gaae vi fra natura naturans over til natura naturata, gjenfinde
vi den samme Grundtankes Præg. Causalitetens Nødvendighed, der
udtrykker Aarsagens Væsen, maa udtrykke sig i Virkningen, der kun
har sin Væren i hiin. Naar Substantsens \emph{Uendelighed} medførte, at
Alt, hvad der er, maatte være i Gud, og ikke kunde begribes uden
ham (Eth. I, 15); at de enkelte Ting maatte opfattes som
Affectioner eller modi ved Guds Attributer, ved hvilke disse
udtrykkes paa en bestemt, begrændset Maade (Eth. I, 25 Coroll.);
at Gud ikke blot er den bevirkende Aarsag til Tingenes Existents,
men ogsaa til deres Væsen (Eth. I, 25); saa ligger deri, at Gud
kaldtes alle Tings Aarsag i \emph{samme Betydning} som han kaldtes causa
sui, at enhver Bestemmelse til Virksomhed i Tingen er en
Bestemthed ved Gud (Eth. I, 26); at der, da Guds Virken er i
Enhed med hans Væsens Nødvendighed, i Naturen ikke kan gives
noget Tilfældigt (Eth. I, 29); at Tingene ikke have kunnet
frembringes af Gud paa anden Maade og i anden Orden end de ere
frembragte (Eth. I, 33), og at alt det, vi begribe som værende i
Guds Magt, nødvendigviis maa være (Eth. I, 35). Og i denne ved
Nødvendighedens Baand sammenholdte Tilværelse kan der intet
selvstændigt Udgangspunkt for en Virken danne sig, der kan ikke
blive Tale om en individuel \emph{Frihed}, om Valgfrihed (jfr. Ep. 62)
ligesaalidt som om Tilfældighed og Mulighed. Disse Forestillinger
have kun deres Oprindelse fra den menneskelige Erkjendelsesevnes
Ufuldkommenhed. Gyldighed for Tingene have Begreberne
\emph{Nødvendighed og Umulighed} (skjøndt det sidste, som blot Negation,
naturligviis ikke kan være en Affection ved et Værende --- entis
affectio. Cog. met. I. c. 3). En Ting kaldes
% --- p. 69 ---
\setcounter{footnote}{0}%
nemlig \emph{nødvendig},\footnote{Nødvendigheden indbefatter saaledes, i
sin dobbelte Betydning, Frihed og Tvang, og der kan tales om en
libera og en coacta necessitas (Ep. 62, jfr. 58. 60. Cog. met. I,
c. 3).} enten med Hensyn til dens Væsen eller med Hensyn til dens
Aarsag. Thi en Tings Existents følger med Nødvendighed enten af
dens eget Væsen eller Definition, eller af en given virkende
Aarsag (Eth. I, 33. schol. 1). Modsætningen til det Nødvendige
bliver det \emph{Umulige}, hvilken Bestemmelse vi bruge naar enten
Væsenet (Definitionen) indeholder en Modsigelse eller naar der
ingen ydre Aarsag gives, der er bestemt til at frembringe Tingen.
Derimod udsige Betegnelserne: \emph{det Tilfældige} og \emph{det Mulige}, kun en
Mangel ved vor Erkjendelse. Naar vi nemlig, idet vi blot betragte
en Tings Væsen, ikke finde Noget, der med Nødvendighed bekræfter
eller benægter dens Existents, saa kalde vi den \emph{tilfældig}; \emph{mulig}
kalde vi den naar vi, idet vi see hen til de Aarsager, af hvilke
den skulde frembringes, ikke vide om de ere bestemte til at
frembringe den (Eth. IV, def. 3 og 4).\footnote{Cog. metaph. I,
c. 3. Res possibilis --- dicitur, quum eius causam efficientem
quidem intelligimus, attamen an causa determinata sit ignoramus.
--- Si autem ad rei essentiam simpliciter, non vero ad eius
causam attendimus, illam \emph{contingentem} dicemus --- --- quia ex
parte essentiæ nullam in ipsa reperimus necessitatem existendi
--- --- neque etiam implicantiam sive impossibilitatem. I Ethiken
I, 33. schol. 1. bruges Udtrykkene som Synonymer.}

Ligesaalidt som Bestemmelserne: Tilfældighed, Mulighed og
Frihed, kunne de Bestemmelser finde Anvendelse paa natura
% p.69: `natura naturata' prints in plain roman here, unlike the
% genuinely spaced instances on pp. 68/71/72 -- verified at 600 dpi.
naturata, der ikke udtrykke den i dens nødvendige Sammenhæng med
sin Aarsag; men kun i dens tilfældige Forhold til den
menneskelige Indbildningskraft. Derfor udelukkes saadanne
Bestemmelser som Orden og Forvirring, Godt og Ondt,
Fuldkommenhed og Ufuldkommenhed, Skjønhed og Hæslighed; og
saadanne Bestemmelser, der ikke ere »notiones castæ, quæ
naturam, ut in se est, explicant«, men »notiones ex vulgi usu
factæ«, der blot udtrykke Naturen saaledes som den refereres til
de menneskelige Sandser,
% --- p. 70 ---
\setcounter{footnote}{0}%
f. Ex. Synligt og Usynligt, Koldt og Varmt, Flydende og Fast o.
desl. (Ep. 6). I Naturen fremgaaer Alt med evig Nødvendighed og
med samme Fuldkommenhed af Guds nødvendige Væsen; der tilkommer
ingen modus Andet end hvad der ligger i dens Begreb, og dette
Begreb bestemmes ikke ved det tilfældige, usande Forhold til os,
men ved det nødvendige, sande Forhold til Gud (jfr. Eth. II, 32.
Ep. 58). Den nærmere Udvikling heraf vil blive givet i Cap. 5.

Gaae vi til Slutning fra den i sit sande Forhold til det
Guddommelige opfattede natura naturata, ved hvilken det dog ikke
har kunnet undgaaes, paa Grund af de i forrige Capitel udviklede
Inconsequentser hos Spinoza, at bringe Bestemmelser ind, der
strængt taget ere udelukkede ved dens Begreb, over til den, i
Abstractionen fra sin Aarsag sondrede natura naturata, til det
hypothetisk fastholdte Endeliges uendelige Række, da gjenfinde vi
atter her Nødvendighedens Bestemmelse. Det løsrevne Endelige i
dets begrændsede Existents har ikke i sig en bestemmende Grund
til at være og virke, men maa bestemmes dertil ved et Andet. Men
dette Andet bliver her, paa det Endeliges Gebet, et andet
Endeligt, der atter fordrer en anden endelig Aarsag og saaledes i
det Uendelige (Eth. I, 28). Ved denne abstracte Opfattelse af det
Endelige bliver altsaa Tingenes Bestemmen ved Gud (jfr. Beviset
for den anførte Sætning) opfattet som den i en uendelig Række
fortløbende Bestemmen ved hverandre.\footnote{Bestemmetheden
bliver i de endelige Ting en endelig Bestemmethed, altsaa ikke
ved Gud, men ved et Endeligt.} Paa dette Punkt faaer Systemet
Charakteren af en stræng Determinisme; men det sees let, at
Determinismen hos Spinoza maa faae en væsentlig anden Betydning
end i de Systemer, hvor der tillægges det Enkelte en reel
Existents. Idet det Enkelte, opfattet i sin Sandhed, opfattes i
det Uendelige, bliver det ene Endeliges udenfra kommende
Bestemthed ved det andet forvandlet til dets Bestemmen ved det
Uendelige, der er dets sande Væsen. Determinismen bliver saaledes
en kun tilsyneladende, og
% --- p. 71 ---
% No footnote on this page (confirmed on image).
% Calibration page for this book's letterspacing: prints
% "n a t u r a  naturans" and "n a t u r a  naturata" with only
% the first word of each pair spaced -- transcribed as printed.
% "Stcd" for "Sted" (first line): weak crossbar on an e, already
% known from recon; sense-reading "Sted" transcribed, doubt
% logged, no printer's defect claimed (JBIG2 scan, see JBIG2-TEST.md).
det vil paa et senere Sted (i Cap. 6) vise sig, at det Samme
gjælder om den Egoisme, der bliver en Consequents af Opfattelsen
af Væsenets og Magtens Enhed ogsaa i det Enkelte; thi idet
Enkelthedens Grændse er et blot Negativt, forsvinder det, der
constituerer Egoismen.

Forestillingen om Nødvendigheden som \emph{ydre} Tvang hænger altsaa
sammen med den usande Forestilling om et for sig bestaaende
Endeligt; idet denne Forestilling giver Plads for den sande Idee
om det Uendelige, bliver den ydre Nødvendighed til en \emph{indre}, til
Frihed ɔ: Bestemmelse ved sit eget Væsens Lov; Causalitetsforholdet
i Substantsen, dens Forhold til sig selv som causa sui, omfatter
ethvert Nødvendighedsforhold, ligesom vi ovenfor saae, at
Substantsbegrebet i sig indesluttede Modificationerne. I Gud, den
eneste causa libera (Eth. I, 17. Coroll. 2) faae de enkelte modi,
opfattede under Evighedens Form, i deres Enhed med ham, deres
Frihed (jfr. Ethikens 5te Bog: de libertate humana). Problemet om
det Endelige er her det samme som i forrige Capitel, de
Vanskeligheder, der her blive tilbage, følge umiddelbart af det
der Udviklede.

% Rule closing Tredie Capitel; appears single at this resolution
% (contrast the double rule closing Cap. 1 on p. 43 and again
% noted at the Cap. 2/3 boundary on p. 64 -- worth re-checking).
\bigskip
\begin{center}\rule{2.5cm}{0.4pt}\end{center}
\bigskip

\begin{center}
{\large\textbf{Fjerde Capitel.}}\\[4pt]
\textbf{Spinozas Naturbegreb.}\\[6pt]
\rule{2cm}{0.4pt}
\end{center}
\addcontentsline{toc}{chapter}{Fjerde Capitel. Spinozas Naturbegreb}
\markboth{Spinozas Naturbegreb}{Spinozas Naturbegreb}
\bigskip

Naar man hos Spinoza taler om et Naturbegreb, til Forskjel fra
Substantsbegrebet, maa det erindres, hvad der i andet Capitel
blev udviklet, at denne Forskjel er en forsvindende, en evig
ophævet. Substantsen som causal sætter som Virkning natura
naturata, men i denne sætter den ikke et Andet, men sig selv, som
Aarsag til den er den causa sui, Enhed af det Virkende og det
Bevirkede, og denne Enhed udtrykkes idet den samme Betegnelse
bestemmer Aarsag og Virkning, \emph{natura} naturans og \emph{natura} naturata.

% --- p. 72 ---
\setcounter{footnote}{0}%
\emph{Naturen} bliver det fælles Navn, Gud og Naturen blive Synonymer.
»Vi have viist, at Naturen ikke handler efter et Øiemed,«
saaledes viser Spinoza tilbage til det i Eth. I, 36 schol.
Udviklede, hvor dette er bevist med Hensyn til Gud. Ȯternum
illud et infinitum ens, quod \emph{Deum seu naturam} appellamus, eadem
qua existit necessitate agit« (Eth. IV. præf.) --- »Ipsa \emph{Dei sive
naturæ} potentia« (Eth. IV, 4. dem.). --- »Cognitionem unionis,
quam mens cum \emph{tota natura} habet« (De intell. emend. pag. 360).
% p.72: `cum' prints unspaced here; only `tota natura' is spaced (600 dpi).
--- »Ex \emph{natura}, sub quovis attributo considerata, infinita
% p.72: `Ex' prints unspaced here; only `natura,' is spaced (600 dpi).
sequi« (Eth. III, 2. schol. jfr. Ep. 21. 58)\footnote{jfr. Eth.
IV. 50. schol.: qui recte novit, omnia ex \emph{naturæ divinæ}
necessitate sequi, et secundum æternas naturæ leges et regulas
fieri.}. I disse Citater betegner altsaa natura det alle
Attributer, hele Tilværelsen Omfattende, sig selv Frembringende,
dets Begreb er ikke indskrænket til hvad vi kalde Natur i
Modsætning til Aanden, om end Navnet bærer Spor af, hvilken
Betydning Naturen som Natur havde for Spinozas Opfattelse.

I Naturen, opfattet som Synonym med causa sui, som Enhed af
natura naturans og natura naturata, indesluttes nu (jfr. Cap. 2)
de selvstændige Attributer og deres uendelige Modificationer. Af
de uendelig mange Attributer opfattes kun to af den menneskelige
Aand (jfr. Cap. 5), nemlig \emph{Udstrækning og Tænkning}, fordi kun
disse to indbefattes i Aandens Idee. Til Udviklingen af disse to
Attributer og deres Forhold indskrænker derfor Udviklingen af
»Naturens« Begreb sig. Den bestemte begrændsede Tanke: \emph{Ideen}, og
den bestemte begrændsede Udstrækning: \emph{Legemet}\footnote{Per corpus
intelligo modum, qui Dei essentiam, quatenus ut res extensa
consideratur, certo et determinato modo exprimit. Eth. II, def.
1.}, have deres positive Bestaaen i Tænkningens og Udstrækningens
Attributer, der begge tilhøre den uendelige Substants (Eth. II,
1. f.), men som selvstændige. Den samme Substants virker altsaa i
begge Attributer, sætter begges Modificationer, og det saaledes,
at Modificationerne af de
% --- p. 73 ---
\setcounter{footnote}{0}%
enkelte Attributer danne sideordnede og \emph{parallele}
Rækker\footnote{jfr. Ep. 4. De intell. emend. pag. 378 f., princ.
phil. Cart. I, 21.}. Hver enkelt Modifications Begreb indeslutter
kun det Attributs Begreb, hvis Modification det er, og har Gud
til Aarsag forsaavidt han betragtes alene under dette Attribut,
og ikke under noget andet. Ideerne ere saaledes ikke Afpræg af
Tingene, ikke Virkninger af Gjenstandene, og disse ere
ligesaalidt en Følge af Ideerne. Tingenes objective Væren følger,
forsaavidt de ikke ere Modificationer af Tænkningen, ikke af den
guddommelige Natur som tænkende, men af det Attribut, ved hvilket
de begribes. Parallelismen i de selvstændige Rækker er begrundet
i, at Virkningens Begreb indeslutter Aarsagens (Eth. I, axiom.
4), saaledes bliver \emph{Ordenen og Sammenhængen i Ideerne} (ɔ: i
Tænkningens Modificationer) \emph{den samme som Ordenen og
Sammenhængen i Tingene} (Eth. II, 7)\footnote{Begrebet om en
Sammenhæng og Overensstemmelse i hele Naturen udvikler Spinoza i
Ep. 15. Mod den vulgære Forestilling om en ved Forholdet til
Indbildningskraften bestemt „Orden“ i Naturen polemiserer han
derimod. Eth. I, 36 schol.}. »Ligesom den udstrakte og den
tænkende Substants saaes at være een og samme Substants, der
snart begrebes under hiint, snart under dette Attribut, saaledes
er Udstrækningens Modification og denne Modifications Idee een
og samme Ting, men udtrykt paa to Maader. --- --- Saaledes er f.
Ex. Kredsen, der existerer i Naturen, og den existerende Kredses
Idee, der ogsaa er i Gud, een og samme Ting, der bestemmes
(explicatur) ved forskjellige Attributer. --- Kredsens Idees
objective Væren (esse formale) kan nemlig kun begribes ved en
anden Tankemodus som nærmeste Aarsag, denne atter ved en anden og
saaledes i det Uendelige, saa at saalænge Tingene betragtes som
modi af Tænkningen, saa maae vi forklare hele Naturens Orden
eller Aarsagernes Sammenknytning alene ved Tænkningens Attribut;
men forsaavidt de betragtes som Udstrækningsmodi, saa maa lige
saa fuldt hele Naturens Orden forklares alene ved Udstrækningens
% --- p. 74 ---
\setcounter{footnote}{0}%
Attribut, og det Samme gjælder om de andre Attributer.« (Eth. II,
7. schol.)\footnote{Gjennem denne Parallelisme af Tænkning og
Udstrækning udtrykker sig atter det Spinozistiske Grundbegreb:
Causaliteten --- i Modsætning til Øiemedet. I Øiemedsbegrebet
ligger Tankens Prioritet, Udstrækningens Bestemmen ved den; denne
Overvægt er hos Spinoza udelukket med Øiemedsbegrebet, og i hiin
Parallelisme udpræger sig Naturnødvendighedens Virkningsform.}.

Betragte vi nu, efterat have seet den almindelige Bestemmelse af
Attributernes Forhold til hverandre indenfor Naturen i videre
Forstand, det Attribut, der constituerer hvad vi i snævrere
Forstand kalde \emph{Natur: Udstrækningens}\footnote{I Eth. I, 15.
schol. bruges Udtrykket: substantia \emph{corporea}.} Attribut, hvis
Begreb forudsættes hvor den menneskelige Aands Væsen skal
udvikles (jfr. Eth. II, 13), saa er det strax iøinefaldende, hvor
forskjellig Spinozas Opfattelse af Udstrækningens Begreb bliver
fra hans Forgænger Descartes's. Hos denne var Udstrækningen
udelukket fra Guds Væsen og sat som en Bestemmelse ved den
hvilende, passive Materie, eller rettere: Bestemmelsen Materie
gik op i Bestemmelsen Udstrækning, opfattet som hvilende,
saaledes at Bevægelse og Forskjel bringes ind ved en ny postuleret
Akt af Gud; og idet Udstrækningen opfattes som passiv, opfattes
den som delelig i det Uendelige. Hos Spinoza er Udstrækningen
opfattet som \emph{Attribut}\footnote{Om Udstrækningens Forhold til Gud
jfr. princ. phil. Cartes. I, 9. schol. Cog. met. II.}, altsaa
ikke blot som uendelig og \emph{udelelig}; men ogsaa som \emph{activ}. Den
sidste Bestemmelse er givet idet Causalitetsbestemmelsen,
hvorunder Substantsen opfattedes, umiddelbart faaer Gyldighed for
Attributerne. Hvert Attribut udtrykker en \emph{actuosa} essentia, af
hvert følger en Uendelighed (Eth. III, 2. schol. II, 3. schol.).
Hvor Tænkningen og Udstrækningen af Spinoza stilles imod
hinanden, modsættes undertiden cogitandi potentia og \emph{agendi eller
operandi potentia} (f. Ex. Eth. III, 28. dem. Ep. 45).

% --- p. 75 ---
\setcounter{footnote}{0}%
Ligesom Bestemmelsen af Udstrækningen som activ er en umiddelbar
Consequents af dens Opfattelse som et af Substantsens Attributer,
saaledes ogsaa dens \emph{Uendelighed og Udelelighed} (jfr. Eth. I, 8.
12. 13. Coroll. 15. schol.). Disse to Bestemmelser, af hvilke
Descartes kun kunde indrømme Uendeligheden, ere efter Spinozas
Opfattelse uadskillelige. Hvad der med Hensyn til Udeleligheden
bedrager os, er, »at Quantiteten af os opfattes paa en dobbelt
Maade, nemlig abstrakt\footnote{„A substantia abstractam.“ Ep.
39.} eller overfladisk, saaledes nemlig som vi imaginere den, og
som Substants, hvilket kun skeer ved Forstanden (intellectus).
Naar vi derfor have Opmærksomheden henvendt paa Quantiteten
saaledes som den er i Forestillingen, hvilket skeer hyppigt og
lettere, saa ville vi finde den endelig, delelig\footnote{Med
Hensyn til den uendelige Deling jfr. Ep. 6. De int. emend. pag.
385; princ. phil. Cart. II, 5. schol. --- Om Begrebet: det tomme
Rum jfr. Ep. 6.} og sammensat af Dele; men naar vi have
Opmærksomheden henvendt paa den saaledes som den er i Forstanden,
og begribe den som Substants, \emph{hvilket er meget vanskeligt}, saa
vil den findes at være uendelig, eneste og udelelig. Dette vil
være tilstrækkelig tydeligt for Enhver, der veed at gjøre
Forskjel paa Forestilling og Forstand, især naar der tages
Hensyn til, at \emph{Materien} overalt er den samme, og at man ikke
adskiller Dele i den undtagen forsaavidtsom vi opfatte Materien
som bestemt (affectam) paa forskjellige Maader, hvorfor dens Dele
kun adskilles modaliter, ikke realiter. Vandet f. Ex. som Vand
kunne vi forestille os deelt, saaledes at dets Dele adskilles,
men ikke forsaavidt det er den legemlige Substants; thi
forsaavidt kan det hverken deles eller adskilles. Ligeledes
bliver Vandet til og forgaaer som Vand; men forsaavidt det er
Substants, bliver det hverken til eller forgaaer« (Eth. I, 15.
schol.)\footnote{jfr. Ep. 29: Som Udtryk for Substantsen kan
Udstrækningen ikke begrændses eller deles uden at dens Begreb
ophæves. »Quare ii prorsus garriunt, ne dicam insaniunt, qui
substantiam extensam ex partibus sive corporibus ab invicem
realiter distinctis conflatam esse putant. Perinde enim est, ac
si quis ex sola additione et coacervatione multorum circulorum,
quadratum aut triangulum, aut quid aliud tota essentia diversum
conflare studeat.“ --- Paa dette Punkt bestemmer Aands-Elementet
i Substantsbegrebet Natur-Elementet, ligesom dette paa mange
andre Punkter bestemmer hint.
% This note opens with a guillemet and closes with what reads, at
% this resolution, as a raised/curly mark rather than a mirrored
% guillemet -- i.e. a second MIXED quote pair in the book (opens
% low/angled, closes high), alongside the one already logged at
% p. 148. Printed as it appears; not normalised.
}.

% --- p. 76 ---
\setcounter{footnote}{0}%
I Opfattelsen af \emph{den udstrakte Substantses} Væsen staaer altsaa
Spinoza i skarp Modsætning til Descartes. Derimod kommer han ham
nærmere ved Opfattelsen af de endelige Udstrækningsmodi (see
nedenfor). Ikke blot statuerer han, ligesom Descartes, at
Udstrækningen ingen andre reelle Modificationer modtager, end de,
der ere givne med Bevægelsen og de dermed forbundne Bestemmelser:
Figur og Stilling, saa at alle andre Bestemmelser, f. Ex.
qualitates tactiles, kun udtrykke et imaginært Forhold til vore
Sandser (Ep. 6), men, efter sin Opfattelse af det Endelige, maa
han sætte enhver Bestemmethed til Bevægelse i den enkelte
Udstrækningsmodus som fremkaldt \emph{udenfra}, om end det Bevægedes
egen Natur bliver en medvirkende Factor (jfr. Eth. I, 28 og II,
13. schol.), og han maa lade de udenfra fremkaldte Forandringer i
Legemernes Verden foregaae alene efter Mechanikens Love (Ep. 9).
Men denne Overensstemmelse taber atter sin Grundvold idet det
Endelige, Betingelsen for det Mechaniske, forsvinder.

Udtrykkene: \emph{Udstrækning}, \emph{udstrakt Substants} (legemlig
Substants), \emph{Quantitet}, \emph{Materie}, bruger Spinoza som
Synonymer, og forsaavidtsom Bestemmelsen Udstrækning nærmest
fører hen til Forestillingen om det anskuede Hvilende, synes
Bestemmelsen \emph{Materie} at være en hensigtsmæssigere Betegnelse. I
et af hans sidste Breve (Ep. 72, fra 15de Juli 1676)\footnote{jfr.
det ovenfor anførte Citat af Eth. I, 15. schol.} synes han ogsaa
at være tilbøielig til at substituere denne sidste Betegnelse for
hiin, vel netop for at faae det active Moment
udtrykt\footnote{At Spinoza ikke har optaget materia som Attribut
istedenfor extensio, kan vel ogsaa have sin Grund i Hensyn til
den almindelige Gudsforestilling.}. Spinozas Yttringer
% --- p. 77 ---
\setcounter{footnote}{0}%
i dette Brev ere fremkaldte ved en Forespørgsel af en Ven: (Ep.
69) hvorledes de ved Bevægelse og Figur bestemte Legemers
Existents a priori kunde udledes af Udstrækningens Begreb.
Spinoza indskrænker først sit Svar (Ep. 70) til Udstrækningen,
saaledes som Descartes havde opfattet den, nemlig som en \emph{moles
quiescens}, og viser, at det ved denne Opfattelse bliver absolut
umuligt at deducere Legemernes Existents apriorisk, da den
hvilende Materie vil vedblive at være i Hvile naar den ikke
udenfra sættes i Bevægelse af en stærkere Aarsag. Da
Spørgsmaalet gjentages (Ep. 71) og han opfordres til at udtale
sin egen Anskuelse, svarer han undvigende (Ep. 72):
% No opening guillemet before "quod petis" on the image -- checked
% at the point where a » would be expected, after the colon; both
% OCR witnesses likewise show none. The passage closes with a
% guillemet six lines below ("licuit.«"). Transcribed as printed:
% a dropped opening quotation mark, structural, logged here rather
% than silently supplied (cf. the dropped closing « at p. 45).
quod petis,
an ex solo extensionis conceptu rerum varietas a priori possit
demonstrari, credo, me jam satis clare ostendisse, id impossibile
esse, ideoque materiam a Cartesio male definiri per extensionem;
sed eam necessario debere explicari per attributum, quod æternam
et infinitam essentiam exprimat. Sed de his forsan aliquando, si
vita suppetit, clarius tecum agam. Nam hucusque nihil de his
ordine disponere mihi licuit.« --- Det Undvigende i Spinozas Svar
har vel sin Grund i, at det Spørgsmaal, der blev stillet, berørte
de to Punkter af hans Lære, der vare mest fremmede for hans
Samtids Opfattelse og den haardeste Anstødssteen for den, nemlig:
\emph{Udstrækningens ɔ: Materiens Enhed med Gud, og de endelige
Tings Ikke-Realitet}. Seet fra Substantsen og i Substantsen, og
det er: seet i sin Sandhed, er det Endelige ophævet som Endeligt,
der kan derfor ikke blive Tale om a priori at deducere det af den
uendelige, udelelige Udstrækning; thi netop for den aprioriske
Erkjendelse, der er Erkjendelsen af Substantsen, existerer det
ikke\footnote{I Spinozas Svar maa der vel derfor ogsaa lægges
Eftertryk paa Ordene: \emph{rerum varietas}. Tænkes nemlig de enkelte
Ting som virkelig forskjellige, altsaa som virkelig bestaaende
enkelte, saa kunne de naturligviis ikke deduceres, da de som
saadanne ere usande og kun ere saadanne ved noget Negativt.}. Men
% --- p. 78 ---
\setcounter{footnote}{0}%
netop idet den af Descartes søgte Udvei: en personlig Guds
vilkaarlige Indgriben, var afskaaren, maatte Problemet stille sig
skarpere frem, de Vanskeligheder, der i Cap. 2 fremkom ved det
almindelige Spørgsmaal om det Endelige, maatte her gjentage sig
indenfor det enkelte Attribut, og de der angivne
Mellembestemmelser maatte her søges anvendte.

Som Udstrækningens \emph{uendelige modi}\footnote{In scrutandis rebus
naturalibus ante omnia investigare conamur res maxime universales
et toti naturæ communes, videlicet motum et quietem, eorumque
leges et regulas, quas natura semper observat et per quas
continuo agit, et ex his gradatim ad alia minus universalia
procedimus. tract. theol.-pol. pag. 88. --- Corpus humanum non
est absolute, sed tantum certo modo \emph{secundum leges naturæ
extensæ per motum et quietem} determinata extensio. Princ. phil.
Cart. præf. editoris.} bestemmes nu \emph{Bevægelse og
Hvile}\footnote{Dette: Bevægelse og Hvile, maa tages copulativt,
saaledes at det bliver enstydigt med: Forholdet mellem Bevægelse
og Hvile, der saa atter finder sit Udtryk i facies totius
universi. Hvile kan, efter Spinozas Opfattelse af Substantsen,
kun blive en relativ Bestemmelse, en Quantitetsforskjel indenfor
Bevægelsen. Kun i de enkelte Legemer kan der tales om en
Bevægelse for sig, og en Hvile for sig, men denne Adskillelse
falder bort med de endelige Legemers reelle Existents.} (Ep. 66,
jfr. Eth. I, 32. Coroll. 2.); som disses Produkt \emph{facies totius
universi}. --- De \emph{endelige} Modi ere \emph{Legemerne}. De \emph{deduceres}
ikke af \emph{hine}; thi en saadan Deduction er, efter hvad der i Cap. 2
blev udviklet, umulig, og erstattes ikke ved det Postulat, at
eftersom Universets Natur er absolut uendelig, saa maa »denne
absolut mægtige Natur bestemme sine Dele paa uendelige Maader og
tvinge dem til at modtage uendelige
Forandringer«\footnote{Ab hac infinitæ potentiæ natura eius
partes infinitis modis moderantur et infinitas variationes pati
coguntur. Ep. 15.}; thi netop \emph{Begrebet variatio} kan ikke
udledes af det Uendelige. De endelige \emph{Legemer forudsættes
hypothetisk}, saaledes at Opfattelsen af dem
% --- p. 79 ---
\setcounter{footnote}{0}%
som enkelte holdes svævende (Eth. II. def. 7), og føres saa
\emph{tilbage} til de uendelige Modificationer\footnote{Ogsaa i
Causalitetsbegrebets Overvægt ligger Umuligheden af at fastholde
et bestaaende Enkelt; thi Virkningen sondrer sig ikke fra
Aarsagen; der er intet Individuationsprincip.}.

I Ethikens 2den Deel, 13de Sætning, schol. er, som Grundlag for
den efterfølgende Udvikling af de ethiske Grundbegreber, \emph{Physikens
almindelige Sætninger} angivne.
% p.79: the compositor's Sperrsatz run continues into `an-' at the line
% break and stops mid-word; `givne' resumes in plain roman. Transcribed
% as printed -- \emph{} not extended into the word (600 dpi).
»De usammensatte Legemer (corpora
simplicissima)\footnote{Begrebet: corpora simplicissima, er
relativt og vilkaarligt; thi idet Legemerne opfattes som
begrændsede, enkelte, ere de sondrede fra Substantsen, og maae
opfattes som delelige i det Uendelige (Ep. 29). Netop ved denne
Modsigelse mod Spinozas bestemt udviklede Opfattelse viser hiin
Bestemmelse sig som hypothetisk.} ere enten i Hvile eller i
Bevægelse, og snart i hurtigere, snart i langsommere Bevægelse.
Kun derved, ikke ved Forskjel i Substants, adskilles de fra
hverandre; de stemme alle overens deri, at de indeslutte det
samme Attributs Begreb og bestemmes ved Bevægelse og Hvile.«
Idet de opfattes som enkelte, maa Bestemmelsen til Bevægelse og
Hvile udgaae fra et andet Legeme, der atter maa bestemmes ved et
tredie og saaledes i det Uendelige (jfr. Eth. I, 28. og ovenfor
S. 70). Den modtagne Bestemthed bevare de saalænge indtil en
modsat sættes i dem ved et andet Legeme, og enhver saadan
Affection er bestemt baade ved det virkende og det paavirkede
Legemes Natur. Hermed ere Bestemmelserne angivne for de som
usammensatte (simplicissima) opfattede Legemer, der kun adskilles
ved Bevægelse og Hvile, Hurtighed og Langsomhed; for de
\emph{sammensatte} Legemer (composita) fordres nærmere Bestemmelser.
»Naar nemlig flere Legemer af samme eller forskjellig Størrelse
saaledes sammenholdes af de andre, at de gjensidigen slutte
sammen (ut invicem incumbant), eller naar de, bevægede med samme
eller forskjellig Hurtigheds Grad, i et bestemt Forhold meddele
hverandre deres Bevægelser, saa kunne disse Legemer kaldes
\emph{indbyrdes forenede}, og siges alle tilsammen at udgjøre \emph{eet
Legeme eller Individuum}, der ved denne
% --- p. 80 ---
\setcounter{footnote}{0}%
Forening adskilles fra de øvrige.« Individualitetens Begreb er
saaledes ikke en indre Bestemmelse fra et blivende Centrum; thi
et saadant er gjort umuligt ved Opfattelsen af det Endelige; ---
men en udenfra givet, tilfældig og relativ Bestemmelse. Dets
Relativitet fremgaaer nu ogsaa af de Udvidelser, der gives det,
og som netop paa Mechanismens Standpunkt opløse det paa de
underordnede Trin\footnote{Ogsaa Spinozas Opfattelse af Grændsen
som blot negativ maatte gjøre Individualitetsbegrebet relativt.}.
Deels udvides det, hvad det af de meest enkelte Dele sammensatte
Legeme angaaer, idet Individet betragtes som det
samme\footnote{retinebit suam naturam absque ulla formæ (ɔ:
Væsenets) mutatione.} om ogsaa nogle af dets Dele udsondres, naar
kun ligesaamange andre af samme Natur komme i Stedet; ligeledes
om ogsaa Delene blive større eller mindre naar de kun under
Forøgelsen eller Formindskelsen bevare det samme Forhold af
Bevægelse og Hvile; ligeledes, om ogsaa nogle af Delene forandre
Bevægelsens Retning, naar de vedblive at bevæge sig og
gjensidigen at meddele hverandre deres Bevægelse paa samme Maade
som før, og endelig, om ogsaa det sammensatte Individuum som
Helhed bevæges eller hviler, eller faaer denne eller hiin
Retning, naar kun hver af dets Dele beholder sin Bevægelse, og,
ligesom tidligere, meddeler den til de andre. Deels udvides
Individets Begreb til det af flere forskjellige Individer
Sammensatte, hvis Affectioner altsaa kunne blive saameget desto
mangfoldigere uden Skade for Individets Identitet. Og idet
saadanne nye Sammensætninger atter kunne sammensættes i det
Uendelige, naaes der til at indsee, »at \emph{hele Naturen er eet
Individuum}\footnote{jfr. Cog. metaph. II, c. 7.}, hvis Dele ɔ:
alle Legemer, vexle (variant) paa uendelig mange Maader uden
nogen Forandring i det hele Individuum.«

Retningen i denne Udvikling er altsaa den samme som i det
almindelige Problem om det Endelige. De enkelte Legemer udledes
ikke af den uendelige Udstrækning, men
% --- p. 81 ---
\setcounter{footnote}{0}%
idet der foreløbig gaaes ud fra den mod Grundprinciperne
stridende Forestilling om \emph{enkelte} Legemer, gives der
Individualitetens Begreb en Elasticitet, der først gjør
Bestemmelsen relativ, og endelig bringer den til at forsvinde
idet den falder sammen med Uendeligheden. Men under denne
tilbagegaaende Bevægelse, denne venøse Strømning, der fører Alt
tilbage til Centrum, uden at en Arteriestrømning atter fører det
ud derfra, er der dog en Bestræbelse for, ved selve det
Forsvindende at give det Uendelige en Fylde, og for Anskuelsen
danner der sig et flygtigt Billede af et uendeligt, i sig evigt
bevæget, ved evige Love ordnet og alt Endeligt omfattende
Naturens Liv, i hvilket Organismens Tanke ligesom afspeiler
sig\footnote{Anskuelsen af evige Naturlove, der beherske en
Forandring og binde den under Enheden, maatte, consequent
gjennemført, føre til Anskuelsen af en gjensidig Meddelen af en
oprindelig Bevægelse, et Bevægelsens Kredsløb, altsaa af en
Organisme; men Organismens Begreb er fremmed for Spinoza, og
allerede en gjensidig Virken kan ikke tænkes i Rækken af
Aarsager og Virkninger. Men antydet er hint Begreb, og kun ved
det faaer Tanken et Indhold.}. Men dette Billede er et
forsvindende; thi det Endelige er ikke blot i Forsvinden, men er
evigt forsvundet\footnote{Om denne Naturopfattelses Forhold til
den Leibnitziske Monadelære jfr. Cap. 10.}.

Kun under Formen af Laanesætninger til Ethiken har Spinoza
udviklet de almindeligste Sætninger af sin Physik (»generalia in
physicis«). En mere detailleret Udførelse var afskaaret ved hans
Princip, og saameget end Spinoza har beskjæftiget sig med Studiet
af Naturen, saa kunde dog hans dybere Interesse ikke hvile paa
Naturen i dens concrete Fylde\footnote{Opfattelsen nemlig af
Naturens concrete Gjenstande i deres particulares differentiæ
faaer væsentlig sin Interesse derved, at disse vise sig at kunne
reduceres til de almindelige mechaniske Principer (Ep. 6).}, der
forsvandt i Substantsens Skjær, ligesom efter Oldtidens Tro det
Endelige maatte gaae tilgrunde, naar det havde seet det
Uendelige; --- den kunde kun hvile paa Naturens abstracte
metaphysiske Begreb. Og dog bliver det
% =====================================================================
%  FEMTE CAPITEL  (printed pp. 82--108 = PDF 92--118; the Indhold gives
%  82--108, and the chapter does run onto p. 108 -- there is no closing rule
%  on p. 107, checked on the image.)
%  Opens PART WAY DOWN p. 82, so there is NO \chapter* here: the fragment
%  below begins with the tail of Cap. 4 at the head of p. 82, then the
%  ornamental rule closing it, then the head of Cap. 5 inline. The agent
%  emits the head block; BATCH-AGENT.md §Heads gives the exact markup.
%  THE ENUMERATIONS ON pp. 95--96, corrected after two independent readings
%  of the page. RECON §8 and the first draft of the batch brief said the run
%  `1) 2) 3) 4)' there is body text -- Spinoza's four modes of perception --
%  and warned against setting it as footnotes. That conflated two different
%  lists. What the page actually carries:
%    -- p. 95 has a genuine IN-TEXT enumeration of THREE items, in full body
%       type; and
%    -- the FOUR-item list of the modi percipiendi is set in the smaller note
%       fount below the rule, as part of footnote 3, which begins on p. 95
%       and continues into p. 96's note block.
%  The file sets both correctly. p. 96 DOES carry a rule above its notes,
%  standing unusually high on the page because the note is so long; an
%  earlier comment denying it has been removed.
% =====================================================================
% --- p. 82 ---
psychologisk og dialektisk interessant, hos ham, som i det
hele hos hiin Periodes idealistiske Tænkere, at see, hvorledes
Naturen, der trænges tilbage af Tanken, og kun opfattes i
dennes Form, fordi Tanken overalt søger sig selv, vedbliver
ligesom at friste, saa at ikke blot Tænkeren drages hen til
Studiet af den, men ogsaa Naturens, om end forflygtigede
Bestemmelser, ubevidst lægge sig ind under Opfattelsen af
Tankens Væsen. Hos Spinoza have vi allerede seet Prøver
derpaa i Bestemmelsen af Substantsbegrebet, vi ville see det
endnu tydeligere i Bestemmelsen af den menneskelige Aand.

De almindelige Sætninger af Physiken fulgte som simple
Consequentser af det almindelige metaphysiske Grundlag og
det hypothetisk tagne Udgangspunkt. Den nærmere Gjennemførelse
deraf, hvis en saadan var bleven mulig, vilde kun
have dannet en Parallel til Ethiken, og derfor ikke givet mere
philosophisk Udbytte end vi allerede have i denne. Ethiken
egnede sig, som vi ville see, til en saadan nærmere Udførelse;
i den ville vi derfor finde det egentlige Complement til det
abstract metaphysiske Grundlag.

% The rule closing Fjerde Capitel reads as a DOUBLE rule at 600 dpi
% (cf. the double rule at p. 43; contrast the single rule closing
% Tredie Capitel, batch 6). Recorded per RESUME-NOTES's request to
% watch this. Markup below follows BATCH-AGENT.md's fixed template
% regardless of single/double, as batch 6 did.
\bigskip
\begin{center}\rule{2.5cm}{0.4pt}\end{center}
\bigskip

\begin{center}
{\large\textbf{Femte Capitel.}}\\[4pt]
\textbf{Den menneskelige Aand.}\\[6pt]
\rule{2cm}{0.4pt}
\end{center}
\addcontentsline{toc}{chapter}{Femte Capitel. Den menneskelige Aand}
\markboth{Den menneskelige Aand}{Den menneskelige Aand}
\bigskip

Gjennem det andet af de to for Mennesket bekjendte
Attributer, \emph{Tænkningens Attribut}, naae vi til Bestemmelsen
af \emph{den menneskelige Aand}. Tænkningen er os umiddelbar
vis som Guds Attribut; thi hver enkelt Tanke indeslutter dens
Begreb (jfr. Cap. 4). Gud er \emph{res cogitans, ens cogitans
infinitum}. Tænkningens Attributs \emph{uendelige Modification
er intellectus absolute in\-%
% --- p. 83 ---
\setcounter{footnote}{0}%
finitus} (Ep. 66)\footnote{Denne betegnes ogsaa som intellectus
\emph{actu} infinitus (Eth. I, 30), ikke forat adskille den fra en
intellectus potentiâ, hvilken Spinoza ikke anerkjender; men
forat udtrykke, at denne intellectus er liig med selve intellectio
(Eth. I, 31. schol.).}. Som Modus hører intellectus, ogsaa
som uendelig, til natura naturata, ikke til natura naturans;
den er ikke selve den absolute cogitatio; men dennes uendelige
Virkning (Eth. I, 31 f.). Som eenstydig med intellectus
infinitus\footnote{Forskjellen mellem intellectus og idea er kun
en Navneforskjel, ligesom i Endelighedens Sphære Forskjellen
mellem mens humana og den idea, der constituerer Menneskets
Væsen. Naar der tales om intellectus infinitus Dei, bliver
dette Udtryk fuldkommen parallelt med idea Dei.} staaer idea
Dei. Den følger, ligesom hiin, af »Attributets absolute Natur«
(Eth. I, 21 dem.), den udtrykker Guds Væsen og Alt, hvad der
med Nødvendighed følger af dette (Eth. II, 3), og ved Guds
Væsens Enhed er den een (Eth. II, 4)\footnote{Cog. metaph. II,
c. 7: idea dei, per quam omniscius vocatur, unica et
simplicissima est. Nam revera deus nulla alia ratione vocatur
omniscius, nisi quia habet ideam sui ipsius, quæ idea sive
cogitatio simul semper cum Deo exstitit; nihil enim est præter
eius essentiam, nec illa alio modo potuit esse.} og uendelig
(Eth. II. 8). --- Tænkningens \emph{endelige} modi ere \emph{Ideerne}: per
\emph{ideam} intelligo mentis conceptum, quem mens format propterea
quod res est cogitans (Eth. II. def. 3)\footnote{Ideen kaldes
mentis \emph{conceptus}, ikke mentis \emph{perceptio}, forat betegne, at
Aanden ikke forholder sig passivt til Objectet, men er activ.}.
I Opfattelsen af Ideen og dens Forhold til intellectus infinitus
og idea dei maa Bestemmelsen af Udstrækningen og dens modi
tjene til Veiledning. Ideerne anskuer Spinoza som begrændsende
hverandre ligesom Legemerne (Eth. I, def. 2), han taler om
\emph{Sammensætning} af Ideer, ligesom om Sammensætning af
corpora simplicissima, og idet Begrændsningen i Ideernes
Sphære ligesom i Legemernes kun opfattes som Negation, bliver
Ideens individuelle Enhed, ligesom Individets Begreb i det
Legemlige, et relativt Begreb, fra hvilket der, idet det
% the `6 *' at the foot of p. 83 is the signature mark of the
% gathering and is not transcribed.
% --- p. 84 ---
\setcounter{footnote}{0}%
udvides ved den negative Begrændsnings Forsvinden, naaes til
det universelle Individuum i Ideernes Sphære, (svarende til det
universelle Individuum i Udstrækningens) det er: til den
uendelige idea Dei\footnote{Da intellectus og idea ere
eensbetydende, udtrykkes det Samme naar der siges, at alle
evige Tankemodi som Enhed constituere Guds evige og uendelige
intellectus (Eth. V, 40. schol.).}, i hvilken al Begrændsning
er hævet, og som, paa samme Maade som dens Analogie i
Udstrækningens Sphære, maa opfattes som uforandret under
Delenes uendelige Vexlen. Naar der om denne idea Dei siges, at
den er i Gud (Eth. II, 3), saa er det i samme Betydning, i
hvilken alt det, der er, er i Gud (Eth. I, 15. 25. Coroll.);
den har sin positive Bestaaen i ham som natura naturata i
natura naturans\footnote{I Ethiken V, 30. dem., hvor der
henvises til Eth. II, 3, synes idea Dei at tillægges Gud som
absolute infinitus, men ogsaa paa dette Sted er det Gud som
causatum, der har denne Idee om sig som causa, og Substantsen
som causatum bliver liig med natura naturata i dens
Uendelighed. Imidlertid er Forskjellen mellem idea Dei, og
cogitatio i dens umiddelbare, uendelige Vished om sig selv
ifærd med at forsvinde, og dette er ogsaa Tilfældet hvor der
uden nærmere Bestemmelse tales om absolutus intellectus Dei
eller intellectus divinus (Eth. II, 11. Coroll. V, 40. schol.
I, 33. schol. 2. II, 40. dem.). --- Som et Motiv for Spinoza
til at værge sig imod, at intellectus tillægges Gud som natura
naturans, kan man vel antage Bestræbelsen forat afskjære de
med hiint Begreb følgende anthropomorphistiske Forestillinger
(jfr. Ep. 58. pag. 570 og Beviisførelsen i Eth. I, 33. schol.
2. og I, 17. schol., hvor man ikke maa oversee, at der
argumenteres ex concessis).}.

Forestillede som \emph{enkelte} staae Ideerne parallele med de
enkelte Legemer, og danne en Række, parallel med Tingenes
Række (jfr. Cap. 4). Guds cogitandi potentia er saaledes liig
med hans actualis agendi potentia ɔ: hvad der følger objectivt
(formaliter) af \emph{Guds uendelige Natur}, det følger
subjectivt (objective) i samme Orden og med samme Sammenhæng
i Gud af \emph{Guds Idee} (Eth. II, 7. Coroll. jfr. Ep. 15; de
emend. intell. pag. 368)\footnote{Derfor maae ogsaa Ideerne om
de ikke existerende endelige Ting eller modi indbefattes paa
samme Maade i Guds uendelige Idee som disse Tings objective
Væsen indbefattes i Guds Attributer; hvoraf altsaa følger, at
saalænge de enkelte Ting ikke existere uden forsaavidt de
indbefattes i Guds Attributer, saa existerer deres subjective
Væsen (esse objectivum) eller deres Idee ikke uden
forsaavidtsom Guds uendelige Idee existerer; og naar de
endelige Ting siges at existere, ikke blot forsaavidt de
indbefattes i Guds Attributer, men ogsaa forsaavidt de siges
at vedvare (durare), saa indeslutte deres Ideer ogsaa
Existents, ved hvilken de siges at have Vedvaren (Eth. II.
8).}. Til
% --- p. 85 ---
\setcounter{footnote}{0}%
enhver Ting svarer saaledes en Idee, der subjectivt udtrykker
hvad der objectivt er i Tingen. Den er Tingens idea, dens
mens, og ethvert Naturindividuum har en saadan mens, de ere
alle besjælede, skjøndt i forskjellig Grad (Eth. II, 13.
schol.), idet Ideens Fortrinlighed ɔ: dens større Realitet, er
afhængig af Gjenstandens Grad af Realitet.\footnote{Jo mere
nemlig Ideens Gjenstand paa een Gang kan virke og lide, desto
mere skikket er Aanden til paa een Gang at fornemme
(percipere) flere Ting, og jo mere Gjenstandens Handlinger ere
afhængige af den selv alene, desto mere skikket er Aanden til
at begribe tydeligt (Eth. II, 13. schol.).}

Paa dette Punkt, hvor Forudsætningerne ere givne til at
bestemme den menneskelige Aands Væsen, kommer Problemet om
Forholdet mellem de to Attributer, der constituere dette, og
de uendelig mange Attributer, atter frem. Af Parallelismen
mellem Attributerne maatte det synes, at man kunde drage den
Consequents, at enhver Ting var bestemt som Modus af alle
Attributer, og at Aanden opfattede dem alle, ikke blot
Udstrækningens Modus, Legemet (jfr. Ep. 67). Men Spinoza
undgaaer denne sidste Consequents, idet han siger: quamvis
unaquæque res infinitis modis expressa sit in infinito Dei
intellectu, illæ tamen infinitæ ideæ, \emph{quibus exprimitur,
unam eandemque rei singularis mentem constituere
nequeunt, sed infinitas; quandoquidem unaquæque harum
infinitarum idearum nullam connexionem cum invicem habeat}
(Ep. 68). Men idet alle de andre Attributers modi opfattes
gjennem cogitatio, hvis modus Ideen er, og idet Ideen bliver
en idea
% --- p. 86 ---
\setcounter{footnote}{0}%
ideæ ɔ: idea cogitationis, kommer man enten til at spalte
Tænkningens Attribut i et uendeligt Antal Attributer, eller
til at sætte de uendelige Attributer i Enhed i Tænkningens
Attribut, eller til at lade dette forsvinde i de andre; i
ethvert Tilfælde undgaaer man ikke Strid med Grundanskuelsen.

Ved Bestemmelserne af Ideernes Forhold til Tingene, er
\emph{den menneskelige Aand} ikke sat som qualitativ forskjellig
fra de andre Tings mentes. Den er en bestemt Modification af
Tænkningens Attribut; det, der constituerer dens actuelle
Væren, er Ideen om en actuelt existerende enkelt Ting (Eth.
II, 11), og denne Ting er \emph{Legemet} ɔ: en \emph{actuelt}
existerende bestemt Udstrækningsmodus (Eth. II, 13). Dette er
givet derved, at vi have Ideer om Legemets Affectioner,
hvilket ikke kunde finde Sted naar ikke Legemet var Aandens
Object. Mennesket bestaaer saaledes af Aand og Legeme, altsaa
af parallele Modificationer af de to Attributer, og
Betydningen af Aandens og Legemets Forening (unio mentis et
corporis) er klar af det Foregaaende. Aanden er idea corporis;
men idet Ideen, om end paa en begrændset Maade, udtrykker
Tænkningens actuosa essentia, maa den opfattes, ikke som et
Afbillede af Legemet, men som activ, og idea corporis bliver
saaledes til en \emph{idea ideæ} (idea mentis), der paa samme Maade
er forenet med Aanden som Aanden er forenet med Legemet.
Legemet og Legemets Idee var den samme Ting, opfattet under to
forskjellige Attributer, Ideen og Ideens Idee den samme Ting,
opfattet under det samme Attribut, Tænkningens. Ideens Idee er
Ideens Væsen (forma ideæ) forsaavidt den betragtes som en
modus cogitandi uden Hensyn til Gjenstanden (Eth. II, 20 f.).
Gaves der kun en idea corporis, vilde Tænkningen være passiv i
Forholdet; Tænkningens Activitet betinger Ideens Idee, den
bestemte Idees Viden af sig selv.\footnote{Sammenligningen
mellem Spinozas Opfattelse af idea mentis og den Fichte'ske
allerede tidligere forberedte Opfattelse af Selvbevidsthedens
Synthese er instructiv (jfr. Cap. 10). Spinozas idea mentis er
ikke den abstracte Selvbevidsthed (Jeg = Jeg), men Aandens
\emph{concrete umiddelbare Selvvished} i sin Erkjenden ---
\emph{Selvfølelse} kunde man sige, hvis ikke dette Udtryk
indeholdt en formel Modsigelse, hvilken Spinoza dog selv
forskylder naar han bruger Udtrykket sentire for percipere.}
% --- p. 87 ---
\setcounter{footnote}{0}%
Opfattelsen af Aanden som idea corporis,\footnote{Eth. II, 19:
mens humana est ipsa idea sive cognitio corporis humani. ---
Eth. III, 57. schol.: idea seu anima eiusdem individui. --- Om
Aandens Forhold til Legemet jfr. Ep. 15. 68. De intell. emend.
pag. 368. Eth. V. præf.} og den derved betingede Overførelse
af Udstrækningens Forhold, af Forestillingen om Begrændsning
og Sammensætning\footnote{Den Idee, der udgjør den
menneskelige Aands objective Væsen, er ikke enkelt, men
sammensat af mangfoldige Ideer, idet det menneskelige Legeme
er sammensat af mange Individer, der hvert for sig atter ere
sammensatte (Eth. II, 15. 13. post. 1.).} paa Tænkningens
modi, Ideerne, behersker nu Spinozas Lære om den menneskelige
Aand som erkjendende og villende. Det er kun ved bestandig at
erindre Parallelismen mellem de to Attributer, at man forstaaer
Spinozas Erkjendelseslære, og det er Anskuelsen af den, der
ligger til Grund for den for Erkjendelseslæren væsentlige
Adskillelse mellem den \emph{adæquate og uadæquate} Idee.
»Den menneskelige Aand er en Deel af Guds uendelige
intellectus, og naar man siger, at den menneskelige Aand
begriber (percipit) dette eller hiint, saa siger man intet
Andet end at Gud, ikke forsaavidt han er uendelig, men
forsaavidt han bestemmes (explicatur) ved den menneskelige
Aands Natur, eller forsaavidt han constituerer den
menneskelige Aands Væsen, har denne eller hiin Idee; naar vi
derimod sige, at Gud har denne eller hiin Idee, ikke forsaavidt
han blot constituerer den menneskelige Aands Natur; men
forsaavidt han sammen med den menneskelige Aands tillige har
en anden Tings Idee, saa sige vi, at den menneskelige Aand
begriber en Ting deelviis eller uadæquat« (Eth. II, 11
Coroll.).\footnote{Eth. III, 1. dem. [ideæ] quæ --- inadæquatæ
sunt in mente, sunt etiam in Deo adæquatæ, non quatenus
eiusdem solummodo mentis essentiam, sed etiam quatenus
aliarum rerum mentes in se continet. --- De intell. emend.
pag. 380: Certum est, ideas inadæquatas ex eo tantum in nobis
oriri, quod pars sumus alicuius entis cogitantis, cuius
quædam cogitationes ex toto, quædam ex parte tantum nostram
mentem constituunt. --- Den menneskelige Aand er sammensat af
adæquate og uadæquate Ideer. Eth. III, 3. dem.}
% --- p. 88 ---
\setcounter{footnote}{0}%
Ved Udviklingen af den menneskelige Aand som erkjendende
gjennemføres den samme Opfattelse af det Endelige, som vi
mødte ved Udviklingen af Udstrækningens Begreb og allerede
tidligere ved Udviklingen af det almindelige metaphysiske
Grundlag, og den samme Stræben efter at føre det Endelige
tilbage til det Uendelige gjentager sig. Ligesom vi der saae,
at Endeligheden kunde opfattes paa en dobbelt Maade: efter sin
Sandhed, som forsvunden i det Uendelige, og forvirret, usandt,
som bestaaende i sig og for sig, saaledes viser det sig ogsaa
her, og denne dobbelte Opfatningsmaade gjør den væsentlige
Adskillelse mellem \emph{Erkjendelsens Arter}.

Gaaer man ud fra den hypothetiske Forudsætning af et
existerende Endeligt i dets Mangfoldighed, saa fører dette til
den Erkjendelsesart, Spinoza betegner som Imaginationen. Det
menneskelige Legeme paavirkes paa mangfoldige Maader af de
andre Legemer, og Aanden udtrykker disse Affectioner. Idet
Legemets Affectioner bestemmes baade ved det afficerende og
det afficerede Legemes Natur (jfr. Eth. II, 13. Lemma 3), saa
indeslutte deres Ideer begge Legemers Natur. »De Affectioner
af det menneskelige Legeme, hvis Ideer forestille os de ydre
Legemer som nærværende, kalde vi Tingenes Billeder (imagines),
skjøndt de ikke gjengive Tingenes Former, og naar Aanden paa
denne Maade betragter Legemerne, kalde vi det at
\emph{imaginere}« (Eth. II, 17. schol.).\footnote{Eth. III,
56. dem.: quatenus imaginamur, sive quatenus afficimur
affectu, qui naturam nostri corporis et naturam corporis
externi involvit; jfr. de intell. emend. pag. 384. Eth. II,
49. schol.: Verborum et imaginum essentia a solis motibus
corporeis constituitur, qui cogitationis conceptum minime
involvunt.} Men af Imaginationens Oprindelse følger, at
Ideerne om de ydre Ting mere udtrykke vort Le-
% --- p. 89 ---
\setcounter{footnote}{0}%
gemes Tilstand end deres Natur (Eth. II, 16. Coroll. 2 jfr. 25
f.), og at de udsætte os for Vildfarelser idet Legemet bevarer
Affectionen, og altsaa Aanden Ideen om det fremmede Legeme som
nærværende, indtil en ny Affection udelukker hiin (Eth. II,
17). Af den ovenfor givne Bestemmelse af de adæquate og
uadæquate Ideer følger det, at der gjennem Imaginationen, der
opfatter Tingene som enkelte og forskjellige, ikke kan naaes
til nogen adæquat Erkjendelse, idet dens Ideer alle maae
indeholde hvad der ikke begribes ud af den menneskelige Aand
alene. Idet nemlig det Endelige opfattes som Endeligt, bliver
det et Led i den uendelige fortløbende Kjæde af endelige
Aarsager og Virkninger, og som Led viser det ud over sig selv
og fordrer, forat begribes, et foregaaende Led som Aarsag, og
derved bliver det umuligt, at den som enkelt fastholdte og
isolerede Tings Idee kan være adæquat; thi det vilde
forudsætte en Løsrivelse fra Rækken som selvstændig, en
Begriben alene ved selve det endelige Leds Natur. Skjøndt
derfor Aanden er idea corporis, cognitio corporis, saa er dog
denne Idee, betragtet som enkelt, ikke adæquat ɔ: klar og
tydelig, og ligesaalidt er Ideen om den menneskelige Aand det
(Eth. II, 28. schol.)\footnote{Ved denne Modsigelse: at
Aanden, der efter sit Væsen er cognitio corporis (et mentis),
bliver ikke-erkjendende, viser Imaginationen og den
tilsvarende Opfattelse af Endeligheden sig som usand, som
ikke-værende.}. Aanden kjender nemlig kun det menneskelige
Legeme og veed det kun som existerende gjennem dets
Affectioner; thi medens selve det menneskelige Legeme som
Virkning ikke begribes ved sig selv, men ved sin Aarsag, saa
indeslutte Affectionerne, der følge af dets Natur, Legemets
Begreb, og gjennem dem erkjender derfor Aanden Legemet, og
gjennem Ideerne om dem sig selv. Men idet disse Affectioners
Ideer indeslutte baade det egne og det fremmede Legemes
Natur,\footnote{Derfor kan Begrebet imaginatio anvendes paa
vor Erkjendelse af \emph{os selv}.} bliver lige saa lidt
Erkjendelsen af hiint som af dette ved den menne\-%
% --- p. 90 ---
\setcounter{footnote}{0}%
skelige Aand alene adæquat, og altsaa ikke heller Aandens
Erkjendelse af sig selv ved sig selv som enkelt. Affectionernes
Ideer, refererede til den menneskelige Aand alene, ere ligesom
Consequentser uden Præmisser, \emph{ideæ confusæ}. De ere
nemlig i Gud, ikke forsaavidt han betragtes som bestemt ved
den menneskelige Aands Idee alene; men forsaavidt han
betragtes som bestemt ved flere andre Ideer.\footnote{I dette
Forhold ligger Muligheden af den abstracte Opfattelse af modi
som adskilte fra deres Aarsag.} Den samme Consequents følger
med Hensyn til det menneskelige Legeme af det Relative i
Begrebet Individuum, idet Legemets Dele bevare en
Selvstændighed som særskilte Individer, af hvilke hvert maa
begribes ved sin særegne, fra det menneskelige Legeme
forskjellige Aarsag (jfr. Eth. II, 19---29). Mangelen paa
adæquat Erkjendelse udtrykker sig ogsaa som mangelfuld
Erkjendelse af vort og de ydre Legemers Varighed.\footnote{Begrebet
\emph{duratio} er, i sin Forskjel fra Evigheden, et relativt Begreb,
der ingen Betydning har for den sande Erkjendelse (jfr. Cap.
8).} Idet nemlig deres Existents som enkelte ikke er afhængig
af deres Væsen, men af »communis naturæ ordo et rerum
constitutio«, hvilken vi ikke kunne overskue i dens
Uendelighed, bliver Erkjendelsen af Existentsen i høi Grad
ufuldstændig, og vi opfatte Tingene som \emph{tilfældige og
forgængelige}, hvilke Bestemmelser altsaa kun udtrykke en
Mangel ved vor Tanke, ikke en Egenskab ved Tingen.

Gjennem Imaginationen, der hviler paa den usande Opfattelse af
det Endelige som bestaaende, naaes der altsaa ikke til en sand
Erkjendelse. Den fremkaldes ved det tilfældige Sammenstød
mellem Tingene (rerum occursu fortuito); Aanden bestemmes
udenfra, og sammenkjæder sine Ideer, ikke efter deres evige og
nødvendige Sammenhæng, efter Fornuftens i alle Mennesker lige
Orden, ved hvilken den begriber Tingene ved deres første
Aarsager; men ved \emph{Erindringen} om den tilfældige Sammenkjædning
af Le-
% --- p. 91 ---
\setcounter{footnote}{0}%
gemets Affectioner. Dens Sammenknytning af Ideerne er, hvad
vi, efter en moderne Terminologie, vilde kalde
Ideeassociationens (Eth. II, 18. schol. 29, schol.).
Imaginationen paavirkes kun af det Enkelte (De intell. emend.
pag. 383), den opfatter det Uendelige som endeligt, begrændset
og deleligt (Ep. 29), er derfor uskikket til at opfatte det
Guddommelige,\footnote{Ep. 60: Deum enim non imaginari, sed
quidem intelligere possumus; jfr. Eth. II, 47. schol.} og da
alle Ideer kun have deres Sandhed ved Gud, blive \emph{alle}
Imaginationens Ideer\footnote{jfr. Eth. III, 56. dem.}
uadæquate, lemlæstede og forvirrede (mutilatæ et confusæ), den
bliver Kilden til ideæ fictæ, falsæ og dubiæ (jfr. de intell.
emend.), af hvilke Vildfarelserne opstaae, samt til entia
imaginationis, der ikke udtrykke Tingenes Væsen, men deres
usande Forhold til os som enkelte (see nedenfor).

Men ved Imaginationens Usandhed er ikke en sand Erkjendelse
udelukket; thi Imaginationen blev usand ved Usandheden i det
Endelige; men til det Endelige i dets sande Opfattelse,
opfattet saaledes, at Begrændsningen er et blot Negativt, og
det Endelige værende i det Uendelige, svarer den sande
Erkjendelsesart, hvis Grundlag denne Opfattelse er. Det Usande
i den forrige Erkjendelsesart var nemlig ikke noget
Positivt:\footnote{jfr. de intell. emend. pag. 392; princ.
phil. Cart. I, 15.} der vilde da i Gud være en positiv falsk
Tænkningsmodus, hvilket er umuligt; men det bestod i den
Privation af Erkjendelsen, som laae i de uadæquate,
ufuldstændige og forvirrede Ideer. Idet den menneskelige Aand,
hvis Individualitet er relativ, kun udgjør ligesom et Fragment
af en mere omfattende Idee, en Deel i en Sammensætning, er
dens Erkjendelse ufuldstændig, ligesom lemlæstet; idet
Begrændsningen hæves og Helheden sættes, ɔ: \emph{idet Ideen føres
tilbage til Gud,} er den sand\footnote{Idet alle Ideer, førte
tilbage til Gud, ere sande ɔ: adæquate, og kun sees som
uadæquate forsaavidt de refereres til en enkelt Tings Aand,
behersker den samme Nødvendighedens Lov baade de adæquate og
uadæquate Ideer i deres Consequentser. Eth. II, 36.} ɔ:
stemmer den fuldkommen overeens med sit
% --- p. 92 ---
\setcounter{footnote}{0}%
ideatum, og sand er enhver Idee, der er fuldstændig i os. Men
det, som er fælles for alle Ting, og som ligemeget er i Delene
og i det Hele,\footnote{I Legemernes Sphære var det Fælles
det, at de alle henhørte til samme Attribut, og at de havde
Mulighed til Bevægelse --- hurtigere eller langsommere --- og
Hvile. Eth. II, 13. Lemma 3. Betydningen af „det Fælles“ for
de andre Attributer fremgaaer heraf.} det constituerer ikke
exclusivt nogen enkelt Tings Væsen, og kan derfor kun opfattes
adæquat. Deraf bliver Consequentsen, at der gives visse Ideer
eller Begreber, der ere fælles for alle Mennesker, og i alle
ere som adæquate ɔ: klare og tydelige (Eth. II. 38.
Coroll.).\footnote{Per ideam \emph{adæquatam} intelligo ideam,
quæ \emph{quatenus in se, sine relatione ad obiectum consideratur,}
omnes veræ ideæ proprietates sive denominationes intrinsecas
habet. Eth. II. def. 4. Den sande Idees Særkjende var
Overeensstemmelsen med dens ideatum (jfr. Eth. 1. axiom. 6.
Ep. 64) og dette er en denominatio extrinseca. --- Som
Synonym med ideæ adæquatæ bruges: claræ et distinctæ, f. Ex.
Eth. II, 28; jfr. Eth. II, 13. schol.: adæquate sive
distincte.} Det er Ideerne om det, der paa een Gang er
eiendommeligt og fælles for det menneskelige Legeme og de
ydre Legemer, og som ligemeget er i den enkelte Deel og i det
Hele. Og idet et saadant Fælles er Grundlaget for en adæquat
Erkjendelse, saa maa Aanden blive saameget desto mere skikket
til et begribe adæquat, jo flere Ting Legemet har tilfælles
med andre Legemer (Eth. II, 39. Coroll.). De adæquate Ideer
blive nu Kilden til den sande Erkjendelse; thi de Ideer, der
følge af dem, maae ogsaa være adæquate.\footnote{Om
Consequentsens Betydning for Sandheden jfr. de int. emend.
pag. 375 f.} Den fra dem udgaaende Erkjendelsesart er
\emph{Forstandens} (intellectus). Dens Eiendommelige er, ikke
blot at være \emph{sand}, men at have Visheden om sin Sandhed,
idet Ideen som activ ikke er noget Stumt som et Billede paa en
Tavle, men den levende Tænkningens Act, der bekræfter sig selv
og umiddelbar giver Visheden (Eth. II, 43.
% --- p. 93 ---
\setcounter{footnote}{0}%
schol.).\footnote{Intelligere er liig certum esse, idea liig
ipsum intelligere. Eth. II, 43, schol. Visheden kan derfor
sættes som eet med selve essentia obiectiva (Ideen), der ogsaa
bliver enstydig med Sandheden. De intell. emend. pag. 366 ff.
Intellectus bliver det Klareste af Alt (Eth. I, 31. schol.);
Sandheden sit eget Beviis (Ep. 74. De intell. emend. 378.
382). Visheden som noget Positivt er forskjellig fra det
Negative: ikke at tvivle: per certitudinem quid positivum
intelligimus, non vero dubitationis privationem. Eth. II, 49.
schol. Ideen som Idee indeholder Affirmation og Negation
(ibid). Idet Fornuftens Erkjendelse indeholder Visheden om sin
Sandhed, udelukkes \emph{Sandsynligheden} fra dens Gebet (jfr.
Ep. 15. 60); ligeledes sættes den som uafhængig af
\emph{Erfaringen}, der kun refererer sig til den endelige
\emph{Existents} (nam experientia nullas rerum essentias
docet; sed summum, quod efficere potest, est, mentem nostram
determinare, ut circa certas tantum rerum essentias cogitet.
Ep. 28), og aldrig kan blive et Modbeviis mod hvad vi begribe
klart og tydeligt (princ. phil. Cartes. II, 6. schol.
Slutning.)} Idet den opfatter Tingene sandt (vere), opfatter
den dem ikke, paa Indbildningskraftens Viis, som tilfældige,
men som nødvendige, og idet Tingenes Nødvendighed er selve
Guds evige Naturs Nødvendighed, opfatter den dem under en
\emph{Evighedens} Form (sub quadam æternitatis specie. Eth. II, 44.
Coroll. 2).\footnote{Sub specie æternitatis et
\emph{necessitatis}. Eth. IV. 62. dem. --- Eth. V, 29. schol.:
res duobus modis a nobis ut actuales concipiuntur, vel
quatenus easdem cum relatione ad certum tempus et locum
existere, vel quatenus ipsas in Deo contineri et ex naturæ
divinæ necessitate consequi concipimus. Quæ autem hoc secundo
% printed "antem" for "autem" -- not a Latin word, differs by a
% single n/u-class sort; sense-reading per the book's standing
% JBIG2 rule (JBIG2-TEST.md), doubt logged, no defect claimed.
modo ut veræ seu reales concipiuntur, eas sub æternitatis
specie concipimus, et earum ideæ æternam et infinitam Dei
essentiam involvunt. --- Eth. V, 30. dem.: res --- sub specie
æternitatis concipere est res concipere quatenus per Dei
essentiam ut entia realia concipiuntur, sive quatenus per Dei
essentiam involvunt existentiam.} Dette ligger ogsaa deri, at
Fornuftens Grundlag (rationis fundamenta) ere Begreber, der
udtrykke det, som er fælles for alle Ting, og som altsaa, idet
de ikke udtrykke nogen enkelt Tings Væsen, maae opfattes uden
noget Forhold til Tiden, under Evighedens Form. Men denne
sande Opfattelse af de enkelte existerende Ting indeslutter
nødvendigen \emph{Ideen om Guds evige og uendelige Væsen}. Tingenes
Existents er
% --- p. 94 ---
\setcounter{footnote}{0}%
nemlig ikke tænkt adskilt fra dens Aarsag som Varighed; men
\emph{efter sin Natur (»loquor de ipsa natura existentiæ«.} Eth. II,
45. schol.), som fremgaaende af Guds Naturs evige
Nødvendighed, af hvilken der følger Uendeligt i uendelige
Modificationer. Tingene ere altsaa tænkte i deres Existents i
Gud; thi om end hver enkelt Ting bestemmes af en anden enkelt
Ting til at existere paa en vis Maade, saa er dog den Kraft,
med hvilken de vedblive i deres Existents, en Følge af Guds
Naturs evige Nødvendighed.\footnote{Tilbageførelsen af den
enkelte Tings Idee til Gud skeer derfor ikke ved en aldrig
afsluttet Tilnærmen gjennem Rækken af de endelige Ting; men
ved, fra Uendelighedens Synspunkt, ligesom med eet Blik at see
Tingene i Gud. I denne Opfattelse bliver Væsen og Existents
(\emph{ipsa existentia, ipsa natura existentiæ}) satte i
Enhed ligesom de ere det i Substantsen.} Idet enhver Tings
Idee indeslutter Guds Idee, er denne nødvendigviis adæquat og
fuldkommen, og \emph{den menneskelige Aand har saaledes en
adæquat Erkjendelse af Guds evige og uendelige
Væsen}\footnote{Det, at den menneskelige Aand kun erkjender to
af Guds uendelig mange Attributer, forhindrer den kun i, at
erkjende Gud fuldstændig (omnino cognoscere), ikke i at
erkjende ham klart og tydeligt, da hvert Attribut erkjendes
fuldstændig ved sig. jfr. Ep. 60.} (Eth. II, 46 f.), og dette
uendelige Væsen og dets Evighed maa være bekjendt for Alle
(Eth. II, 47. schol.).\footnote{jfr. Opregningen af Fornuftens
Kræfter og Egenskaber i Afhandlingen de emend. intell. pag.
390 ff.}

Fra Imaginationens usande og uadæquate Erkjendelse adskiller
altsaa Fornuftens Erkjendelse sig som sand og
adæquat,\footnote{Angaaende det almindelige Forhold mellem
intellectus og imaginatio jfr. Eth. IV. 59. schol. De intell.
emend. pag. 384 f. Ep. 29. 60.} og denne Adskillelse er
begrundet i den forskjellige Opfattelse af det Endelige, der i
hiin var abstract, partiel, under Tilfældighedens Form og
derfor forskjellig for de Forskjellige; i denne i Eenhed med
Gudsideen, under Nødvendighedens og Evighedens Form, og eens
for Alle. Men indenfor den adæquate Erkjendelse adskiller
Spinoza atter en dobbelt Form: \emph{Fornufterkjendelsen}
(ratio), der
% --- p. 95 ---
\setcounter{footnote}{0}%
gaaer ud fra notiones communes, og giver adæquate Ideer om
Tingenes Egenskaber, og \emph{den intuitive Viden}
% printed "intnitive" for "intuitive" -- not a Danish/Latin word,
% differs by a single n/u-class sort; sense-reading per the book's
% standing JBIG2 rule (JBIG2-TEST.md), doubt logged, no defect claimed.
(scientia intuitiva), der gaaer frem fra den adæquate Idee om
Guds Attributers objective Væsen til den adæquate Erkjendelse af
Tingenes Væsen.

De forskjellige Erkjendelsesarter blive altsaa følgende tre: 1)
\emph{Den første Erkjendelsesart: Meningen} (opinio) eller
\emph{Imaginationen}, som dels har sit Udgangspunkt fra de enkelte
Ting, der gjennem Sandserne fremstille sig for os ufuldstændigt,
forvirret og uden Fornuftens Orden --- sine ordine ad intellectum
(dette kalder Spinoza cognitio \emph{ab experientia vaga}); --- dels fra
Tegn, f. Ex. derfra, at der hos os gjennem Ord, som vi høre eller
læse, ved Ideeassociationen fremkaldes et Billede af
Tingene.\footnote{jfr. tract. theol.-pol. pag. 185: \emph{intellectio
pura mente extra verba et imagines}.} --- 2) Den \emph{anden}
Erkjendelsesart, eller Fornuften (ratio); --- 3) Den \emph{tredie}
Erkjendelsesart, \emph{den intuitive Viden}. Som et Exempel (der dog
efter sin Natur maa være uadæquat) til at oplyse Forskjellen, anfører
han den Opgave: af tre givne Tal at finde det fjerde Forholdstal.
»Kjøbmændene have ingen Tvivl om,\footnote{Man erindre Forskjellen
mellem: ikke at tvivle, og: at have Vished (see ovenfor S. 93 Anm.
1).} at man skal multiplicere det andet med det tredie og dividere
Summen med det første, fordi de ikke have glemt hvad de uden noget
Beviis have hørt af deres Lærere, eller fordi de hyppigt have gjort
denne Erfaring ved de simpleste Tilfælde (1ste Art). Men det kan
ogsaa skee i Kraft af Euklids Sætning, ifølge de proportionale Tals
fælles Egenskab (2den Art). Men med de simpleste Tal er Intet af
dette nødvendigt. Naar nemlig Tallene 1, 2 og 3 ere givne, saa seer
Enhver umiddelbart, at det fjerde Tal er 6, og det meget klarere,
fordi vi slutte os til det fjerde af selve det Forhold, som vi med
eet Blik see, at det første har til det andet« (Eth. II, 40. schol.
2).\footnote{I Afhandlingen de intellectus emendatione (pag. 362 f.)
er Inddelingen i Formen noget forskjellig; men danner dog en
Parallel til den i Ethiken givne. Spinoza opregner der 4 modi
percipiendi: 1) \emph{perceptio ex auditu aut ex aliquo signo}; 2)
\emph{ab experientia vaga}; 3) ubi essentia rei ex alia re
concluditur, sed non adæquate f. Ex. naar vi slutte fra en Virkning
til Aarsagen, eller naar der sluttes fra noget Universelt, med
hvilket der altid følger en vis Egenskab; 4) naar Tingen begribes ved
selve sit Væsen eller ved Erkjendelsen af den nærmeste Aarsag. De to
første af disse Erkjendelsesarter svare til hvad der sildigere er
sammenfattet under Betegnelsen: imaginatio. Den tredie skulde
nærmest svare til \emph{ratio}, men der kommer nogen Usikkerhed ind,
idet Spinoza ikke har skjelnet mellem notiones communes og de
abstracte notiones universales (see nedenfor), og derfor ikke
betragter denne, men den 4de Art, som adæquat og sikkret mod
Vilfarelse. Som Exempler paa den tredie Erkjendelsesart anfører han:
naar vi klart have opfattet, at vi fornemme dette bestemte Legeme og
intet andet, saa slutte vi med Klarhed, at Sjælen er forenet med
Legemet, hvilken Forening er Aarsag til en saadan Fornemmelse; men af
hvad Natur hiin Foreningens Væsen er, det begribe vi ikke fuldkomment
deraf, idet vi nemlig ved Foreningen endnu ikke forstaae Andet end
Fornemmelsen af Virkningen, hvorfra Slutningen droges til den
ubekjendte Aarsag. Et andet Exempel er: naar jeg kjender Synets
Natur, og veed, at det har den Egenskab, at vi efter den forskjellige
Afstand see den samme Ting som større eller mindre, saa slutter jeg
deraf, at Solen er større end den viser sig; men dette er kun en
Egenskab (proprium) ikke rei essentia particularis. Denne
Erkjendelsesart er saaledes ikke i og for sig tilstrækkelig idet dens
Resultat er abstract. --- „Ved \emph{selve Tingens Væsen} begribes
Tingen derimod, naar jeg deraf, at jeg veed Noget, veed hvad det er
at vide, eller naar jeg deraf, at jeg har erkjendt Sjælens Væsen,
veed at den er forenet med Legemet. Ved den samme Erkjendelse veed
jeg, at 2 og 3 er 5, og at to Linier, der ere parallelle med en og
samme tredie, ere indbyrdes parallele.“}
% Footnote 3 (above) is printed as ONE continuous note: it opens on
% p. 95 and its text -- INCLUDING the whole "4 modi percipiendi"
% enumeration -- runs on into the note block of p. 96, ending there
% with "indbyrdes parallele." This is footnote text throughout, in
% visibly smaller type than the body above the rule, set BELOW a
% footnote rule that appears after only two short body paragraphs
% ("For bestemtere..." / "Vi saae... notiones communes¹) og at det
% var disse,") -- confirmed at 600 dpi (audit of 3 Sept 2026). The
% rule sits higher up the page than usual only because the note block
% it introduces is so long (this overflow plus p. 96's own footnote
% 1) that it fills nearly the rest of the page; it is NOT absent. An
% earlier note here claiming "no rule (tight page, per recon)" was
% checked against the image and was wrong; corrected.
% This footnote is NOT a second in-text enumeration; the only in-text
% 1)2)3) list in this stretch is the three-part "Erkjendelsesarter"
% list earlier on this page. (BATCH-AGENT.md's note that pp. 95-96
% carry an in-text "1)2)3)4)" enumeration appears to describe this
% footnote's numbered list, not body text; confirmed by direct
% inspection at 600 dpi -- the "4 modi percipiendi" list is set in the
% smaller footnote type, below the rule described above, not in the
% body face. Flagged here so the claim isn't repeated.)

% --- p. 96 ---
\setcounter{footnote}{0}%
For bestemtere at opfatte Forholdet mellem de tre Erkjendelsesarter,
er det nødvendigt at gaae noget nærmere ind paa den Spinozistiske
Terminologie og paa hans Forhold til den tidligere Metaphysiks
Begreber.

Vi saae, at den fornuftige Erkjendelse tog sit Udgangspunkt fra
notiones communes\footnote{Eth. II, 40. schol. 1: notionum, quæ
communes vocantur, quæque ratiocinii nostri fundamenta sunt. jfr.
tract. theol.-polit. pag. 430 (ed. Pauli): ad notiones quasdam
\emph{simplicissimas}, quas communes vocant; sammesteds pag. 88: res maxime
universales et \emph{toti naturæ communes} (der tales om Bevægelse og
Hvile).} og at det var disse,
% Footnote 1 (above) is likewise printed as one note running from its
% mark here on p. 96 on into the note block of p. 97, ending there
% with "(der tales om Bevægelse og Hvile)." before p. 97's own note 1
% begins. Whole note kept at its mark, per BATCH-AGENT.md. Sentence
% runs on across the boundary ("...og at det var disse," + "der gav
% den væsentlige..."), so no blank line here.
% --- p. 97 ---
\setcounter{footnote}{0}%
der gav den væsentlige Adskillelsesgrund mellem den første og anden
Erkjendelsesart. Det bliver derfor af Vigtighed at bestemme dette
Begreb, navnlig at adskille det fra de Imaginationen tilhørende
notiones universales, notiones transscendentales, generalia, entia
rationis o. desl. Naar vi erindre, at notiones communes ere
adæquate, at de udtrykke Tingene efter deres Sandhed (vere), og
tillige mindes Sætningen, at Begrændsning kun er Negation, saa sees
det let, at den egentlige Betydning af: \emph{det, der er fælles},
ikke kan være: det for flere virkelig adskilte Ting fælles
Abstracte; thi efter deres Sandhed existere Tingene ikke som enkelte;
men at Udtrykket maa betegne den gjennem \emph{de mange
tilsyneladende enkelte Ting gaaende positive
Enhed}.\footnote{Som Parallel kan her erindres om Geulincx's
Betegnelse af Gud og Legemverdenen som res singulares, udenfor
hvilke der kun gaves Tilværelsesmodi. Man seer let, at for Spinoza
Bestemmelserne Realisme og Nominalisme i Scholastikens Forstand ikke
faae nogen Betydning; ikke heller ved Sætninger som: universalium
--- --- quæ non sunt, nec ullam habent præter singularium essentiam.
Cog. metaph. II, c. 7.} I det Hele maae vi erindre, at Tallet,
Flerheden, ikke har nogen reel Betydning for Spinoza, idet
Substantsen ikke bestemmes derved; de Consequentser, der følge af
Forudsætningen om Tallet som en Realitet, maae derfor forsvinde.
Notiones communes blive saaledes Ideerne om det positive Ene, og de
ere derfor ogsaa i den Betydning communes, at de ere fælles for alle
de enkelte Mennesker; thi ogsaa disses Existents som enkelte er jo
kun et Skin, de ere kun Dele af den sig selv vidende Idee. Det
subjective og det objective Moment er i Enhed i dette Begreb (jfr.
Eth. II, 38. Coroll. med prop. 37, 38 og 47, dem.).

Fra notiones communes adskilles først og fremmest \emph{termini
transscendentales} (som ens, res, aliquid), og
% sentence continues on p. 98 ("...aliquid), og de almindeligen
% saakaldte..."), so no blank line here.
% --- p. 98 ---
\setcounter{footnote}{0}%
de almindeligen saakaldte\footnote{Ved Begyndelsen af Schol. 2 er
nemlig notio universalis brugt i en videre Forstand. --- Som
Synonymer bruger Spinoza: entia metaphysica (Eth. II, 48. schol.),
ideæ universales (Ep. 50) og notiones abstractæ; den ved hine
Begreber opnaaede Erkjendelse er en cognitio abstracta vel
universalis (Eth. IV. 62 schol.). Af Guds Væsen kunne vi, da han
ikke bestemmes ved Flerhed, ikke danne en universalis idea (Ep.
50).} \emph{notiones universales} (som Menneske, Hest, o. s. v.). Hine have
deres Oprindelse deraf, at det menneskelige Legeme, som begrændset,
kun er modtageligt for et vist Antal Billeder (Indtryk) i tydelig
Form, saaledes at Billederne, naar dette Antal overskrides, begynde
at forvirres, og naar det i høi Grad overskrides, fuldstændig
sammenblandes. Men den menneskelige Aand kan kun tydelig forestille
sig saamange Legemer paa een Gang som der kan være tydelige Billeder
samtidig i Legemet. Naar Billederne i Legemet fuldstændig
sammenblandes, vil Aanden forestille sig dem forvirret uden nogen
Adskillelse, og sammenfatte dem under eet Begreb, som res, aliquid.
--- En lignende Oprindelse have de saakaldte notiones universales.
Der dannes nemlig i det menneskelige Legeme paa een Gang saamange
Billeder f. Ex. af Mennesker, at de vel ikke fuldstændig overgaae
Forestillingens Kraft, men dog saameget, at Aanden ikke kan
forestille sig de enkelte smaa Forskjelligheder og deres bestemte
Antal, men kun forestiller sig det tydeligt, i hvilket de alle,
forsaavidt Legemet afficeres af dem, stemme overens, og dette
betegner man ved det universelle Navn, og udsiger dette om de
uendelig mange enkelte. Men idet disse Begreber have en saadan
Oprindelse, dannes de ikke paa samme Maade af Alle, men ere
forskjellige hos den Enkelte, i Forhold til det, af hvilket Legemet
hyppigst er blevet afficeret, og som Aanden lettest forestiller sig
eller erindrer, (saaledes tænker ved Navnet Menneske det Ene nærmest
paa den opreiste Skikkelse, den Anden paa, at Mennesket er et animal
bipes, risibile, sine plumis, rationale o. s. v.),
% sentence continues on p. 99 ("...o. s. v.), saaledes at i
% Grunden..."), so no blank line here.
% --- p. 99 ---
\setcounter{footnote}{0}%
saaledes at i Grunden disse universales notiones ere \emph{universales
imagines}, som Enhver danner efter sit Legemes Disposition (Eth. II.
40, schol. 1.).

Dernæst adskilles fra notiones communes, idet disse udtrykke et
reelt Existerende, de modi imaginandi,\footnote{Entia non rationis,
sed imaginationis. Eth. I, 36. schol. Dog bruges ikke sjelden
Udtrykket entia rationis, ligesom ogsaa Spinoza som Modsætning til
det, der er i Tingene, kan stille det, der blot \emph{in intellectu}
(de intell. emend. pag. 386).} der ingen tilsvarende Existents have,
men kun have Betydning som Hjælpemidler for Imaginationen til at
fastholde, bestemme og forestille Tingene; saaledes Bestemmelserne:
\emph{Art, Slægt; Tid, Tal, Maal, Modsætning, Orden, Overensstemmelse,
Forskjel o. desl.}, samt de Betegnelser af et Negativt, der udtrykke
en »forvirret Affirmation«, f. Ex. \emph{Ende, Grændse,
Mørke};\footnote{Omvendt er det uendelige Positive, der kun er in
intellectu, ikke in imaginatione, betegnet negativt, f. Ex.
incorporeum, infinitum.} ligeledes de Bestemmelser, der kun betegne
en Mangel ved vor Erkjendelse, som: det Mulige, det Tilfældige (see
ovenfor); Fictioner som: accidentia realia, formæ substantiales,
qualitates occultæ, og endelig saadanne Bestemmelser, der kun
udtrykke et Forhold til os som enkelte, ikke en Bestemmelse af
Tingenes Væsen,\footnote{Ep. 6. Notiones ex vulgi usu factæ, vel quæ
naturam explicant, non ut in se est, sed prout ad sensum humanum
refertur.} som \emph{Orden-Forvirring, Godt-Ondt, Skjønt-Hæsligt o. desl.}
(jfr. Cap. 6).

Imaginationens \emph{almindelige} Bestemmelser udtrykke saaledes ikke et
indholdsrigt Existerende, men kun en uvirkelig Abstraction, et
generaliseret Enkelt, hvis Indhold er udvisket, eller de udsige kun
selve Imaginationens Tilstand, ikke Tingenes Natur. De ere endvidere
hos hver Enkelt forskjelligt modificerede, da det, der ligger til
Grund for dem, er det Enkelte i dets Forskjellighed.

Ved det forskjellige Udgangspunkt bliver Fornufterkjendelsens
Charakteer en anden. Notiones communes udtrykke,
% sentence continues on p. 100 ("...udtrykke, opfattede fra
% Grundbegrebet af..."), so no blank line here.
% --- p. 100 ---
\setcounter{footnote}{0}%
opfattede fra Grundbegrebet af, ikke en forvirret Abstraction, men
det Existerende i dets fyldige Enhed; de udtrykke tillige et Reelt,
og, idet i Tænkningen det Enkeltes uvirkelige Adskillelse er hævet,
udtrykke de dette Reelle paa samme oprindelige Maade i Alle, og
blive almeengyldige Axiomer.\footnote{„Opinionibus et documentis
humano generi universalibus, hoc est notionibus \emph{communibus et
veris};“ (tract. theol.-pol. pag. 50). --- Ex axiomatis
intellectualibus per se notis (sammest. pag. 62). --- Om Forholdet
mellem Definition og Axiom jfr. Ep. 26 f.}

Gaaer man ud fra denne Bestemmelse af notiones communes, der er den
efter Spinozas Grundanskuelse consequente, har man deri en væsentlig
Adskillelse fra Imaginationens Bestemmelser, og den intellectuelle
Erkjendelse kan simpelthen modsættes Imaginationens som den der
hviler paa de sande Ideer, hvis Række fremgaaer i Aanden i
Overensstemmelse med de reale Objecters Række, i Sammenkjædning af
Aarsag og Virkning, og det saaledes, at Aanden er et efter bestemte
Love handlende automaton spirituale.\footnote{De intell. emend. pag.
384. Udtrykket automaton spir. finder allerede hos Descartes
(passions de l'âme Art. 16).} Men Vanskeligheden kommer frem idet
Forholdet mellem den anden og tredie Erkjendelsesart skal bestemmes,
og Bestemmelsen af dette Forhold virker tilbage paa Opfattelsen af
Forholdet mellem ratio og imaginatio og gjør atter Adskillelsen
svævende. Det vil vise sig, at den dybere Grund til de
Vanskeligheder, som her opstaae, er den almindelige Anstødssteen for
Spinozas Philosophie: Opfattelsen af det Endelige.

Fornuftens Erkjendelse (ratio) sættes som adæquat og
sand;\footnote{Cognitio \emph{secundi} et tertii [generis] est
% only "secundi" is letterspaced in the print, not "tertii" -- a
% partial-run case like p. 71's "n a t u r a  naturans"; transcribed
% as printed, not regularised.
necessario vera. Eth. II, 41. jfr. V. 28. dem. I Ethiken IV, 31
bruges Betegnelsen ratio saaledes, at den omfatter baade den anden
og tredie Erkjendelsesart.} den opfatter Tingene ikke som tilfældige
og usandt enkelte, men som nødvendige og under Evighedens Form. Den
fører til Gudserkjendelse, og Gudsideen, den giver, er adæquat og
fuldkommen. Men dog bliver den tredie Er\-%
% word breaks across the page boundary ("Er-kjendelsesart"); no
% blank line, no space -- join directly.
% --- p. 101 ---
\setcounter{footnote}{0}%
kjendelsesart, den intuitive Erkjendelse, sat som et Høiere, der
giver en fortrinligere Art Erkjendelse (Eth. V, 36. schol.). Skal
Grunden til dette Fortrin angives, saa angives den saaledes, at i
hiin \emph{universelle Erkjendelse} (cognitio universalis) Beviset for
Tingenes Afhængighed af Gud, baade med Hensyn til Væsen og
Existents, fører \emph{generaliter} og anvendes paa den enkelte
Ting, medens det i den intuitive Erkjendelse føres af selve den
enkelte Tings Væsen. I denne Opfattelse faae da notiones communes
Betydningen af Begreber om det for de eensartede Ting Fælles, der
udtrykke, ikke Tingenes individuelle Væsen, men en
proprietas.\footnote{jfr. Eth. II, 40. schol. 2; 39, dem. V, 12,
dem. og Bestemmelsen af de forskjellige Erkjendelsesarter i de
intell. emend.} De komme saaledes til at udtrykke et Abstract, og
bringes derfor sammen med axiomata intellectualia, der ere sande,
men abstracte.\footnote{De intell. emend. pag. 380; jfr. pag. 386:
ex axiomatis solis universalibus non potest intellectus ad
singularia descendere, quandoquidem axiomata ad infinita se
extendunt, nec intellectum magis ad unum quam ad aliud singulare
contemplandum determinant.} Idet den anden Erkjendelsesart kaldes
universel (»cognitio universalis, quam secundi generis esse dixi«.
Eth. V, 36. schol.) betegnes den med det samme Udtryk, der brugtes
for Imaginationens Bestemmelser; dens Begreber kaldes, ligesom
Imaginationens almindelige Forestillinger, notiones universales
(Eth. II, 40. schol. 2); og idet den universelle Erkjendelse
opfattes som generel og abstract (de intell. emend. S. 381), men det
Generelle og Abstracte fører ind paa Imaginationens
Gebet,\footnote{jfr. De intell. emend. pag. 363 Anm. h., 372, 376.
Pag. 383 f. føres den abstracte, universelle Erkjendelse
udtrykkeligt tilbage til Imaginationen.} blive Grændserne mellem
dette og Fornuftens usikkre.

I Modsætning til denne abstracte og universelle Erkjendelse gaaer
den \emph{tredie} Erkjendelsesart ud fra Gudsbegrebet (natura, fons
et origo naturæ, ens unicum, infinitum, quod est omne esse, et
præter quod non datur esse), der ikke kan opfattes abstract eller
universelt, og ud fra den adæquate
% sentence continues on p. 102 ("...den adæquate Erkjendelse af
% Attributerne..."), so no blank line here.
% --- p. 102 ---
\setcounter{footnote}{0}%
Erkjendelse af Attributerne udledes den adæquate Erkjendelse af
Tingenes Væsen, der følger af den guddommelige Naturs Nødvendighed.
Den opfatter Ideerne i deres Sammenkjædning som een Idee, saaledes
at Aanden subjectivt udtrykker Naturens objective Væsen, dens
Erkjendelse af Virkningerne bliver kun en fyldigere Erkjendelse af
Aarsagen. Den udleder Intet fra det Abstracte, men fra »essentia
particularis affirmativa«,\footnote{De intell. emend. pag. 388:
Optima conclusio erit depromenda ab essentia aliqua particulari
affirmativa. Quo enim specialior est idea, eo distinctior, ac
proinde clarior est. Unde cognitio particularium quam maxime nobis
quærenda est.} et »ens physicum et reale«, eller fra »en sand og
lovmæssig Definition« ɔ: en saadan, der ikke blot angiver Egenskaber
ved Tingen, men Tingens intima essentia, af hvilken alle dens
Egenskaber kunne udledes, og som gjør dette i en affirmativ Form.
Den Række, Erkjendelsen danner i Ideerne, er »Rækken af de faste og
evige Ting«, ikke den uoverskuelige »Række af enkelte, foranderlige
Ting«, da disse Tings Existents ikke har nogen Sammenhæng med deres
Væsen, eller, med andre Ord, ikke er en »evig Sandhed« (æterna
veritas). Disse »blivende og evige Ting«, der alle ere simul natura,
ere, skjøndt enkelte, dog paa Grund af deres Allestedsnærværelse
(ubique præsentia) og altomfattende Magt (latissima potentia), for
os ligesom universelle eller almindelige Definitioner (genera
definitionum) for de enkelte og foranderlige
Ting,\footnote{Saaledes „udstrække Naturlovene sig til uendelig
mange Ting, og opfattes af os under en Evighedens Form“, tract.
theol.-polit. pag. 72.} og nærmeste Aarsager til Alt; i dem ere, som
i de sande Lovbøger, de Love indskrevne, efter hvilke alle enkelte
Ting blive til og bestemmes; ja disse foranderlige og enkelte Ting
ere saa inderligt og væsentligt (intime et essentialiter) afhængige
af hine blivende, at de hverken kunne være eller begribes uden dem
(de intell. emend. 388 ff.).\footnote{Angaaende Opfattelsen af
Tingenes Existents ved Guds Væsen jfr. Eth. V, 30, dem. 36, schol.}

% --- p. 103 ---
\setcounter{footnote}{0}%
Sammenstiller man imidlertid hvad der saaledes er sat som
eiendommeligt for den tredie Erkjendelsesart, med det, der er
eiendommeligt for den anden, bliver Adskillelsen vanskelig. Ogsaa
denne skal kunne give klare og adæquate Ideer, og af disse udlede
andre klare og adæquate. Men idet den udleder Ideer af andre,
erkjender den Virkningen af Aarsagen, og dette kan ikke blot
opfattes saaledes, at ratio indenfor de enkelte Ideers
Sammenkjædning skulde kunne uddrage adæqvate Ideer, men ikke sætte
dem i Forbindelse med Gud, da dette vilde stride imod Begrebet af
adæquat Opfattelse; thi den opfatter Tingene som de ere efter deres
Væsen, og efter deres Væsen ere de i Gud\footnote{jfr. Eth. V, 29.
schol. 30, dem.}. Erkjendelsen er netop kun adæquat idet Virkningen
føres tilbage til Gud, der er alle Tings causa absolute proxima
(Eth. I, 28), i hvem de indeholdes, og af hvis Naturs Nødvendighed
de følge (Eth. V, 29. schol.). Men idet dette skeer, vilde
Forskjellen mellem ratio og scientia intuitiva forsvinde, da Gud,
der ikke kan opfattes abstract eller universelt, umiddelbart viser
sig som det virkeligt Værende i Tingen, og denne, opfattet adæquat,
altsaa under Nødvendighedens og Evighedens Form, umiddelbart føres
tilbage til den første Aarsag. --- Vilde man, hvad der kunde synes
antydet i Afhandlingen de intell. emend., ved den anden
Erkjendelsesart nærmest tænke paa Slutningen fra Virkning til
Aarsag eller fra Abstracta til Abstracta, da vilde det Første føre
til »lemlæstede og forvirrede Ideer«, det Sidste til Imaginationens
Universalia, altsaa i intet af Tilfældene til en adæquat
Erkjendelse.

Vanskeligheden opstaaer ved det Endeliges tvetydige Stilling.
Fastholdes den uendelige Substantses strænge Begreb, da kan der, som
ovenfor i Cap. 2 er udviklet, hverken umiddelbart eller middelbart
naaes til et Endeligt, altsaa ikke til en bestaaende Flerhed, men
denne er Forudsætningen for den abstracte, universelle Opfattelse.
Erkjendelsens Gjenstand vilde da være een, Opfattelsen af den i dens
% sentence continues on p. 104 ("...i dens Fylde nødvendigen
% intuitiv..."), so no blank line here.
% --- p. 104 ---
\setcounter{footnote}{0}%
Fylde nødvendigen intuitiv, da dens Væren er udelt, uendelig og
evig. Den adæquate Idee vilde kun have det Uendelige til Gjenstand,
thi den udtrykker sin Gjenstand under Evighedens og Nødvendighedens
Form\footnote{Der vilde saaledes ingen Overensstemmelse kunne blive
mellem intellectus og imaginatio, hvilken Spinoza dog indrømmer
(jfr. de intell. emend. pag. 384, Ep. 29. Eth. V, 12) skjøndt
Grændserne ikke drages og ikke kunne drages. Idet de mathematiske
Figurer kaldes entia rationis ɔ: imaginationis (Ep. 72; de intell.
emend. pag. 387) blive de et Berøringspunkt.}, og deri ligger
Uendelighed; den er Tænkningens Affirmation, men det Endeliges
Begrændsning er en Negation. Men den laante, flygtige Existents, der
maa gives det Endelige (jfr. Cap. 2), giver en Tilknytning for den
Erkjendelsesart, der tilhører Mathematiken, og som Spinozas
historiske Forhold maatte lægge ham det nær at hævde en Plads.
Denne kan han imidlertid ikke skaffe tilveie uden at sammenblande
Fornuftens og Imaginationens Bestemmelser. Imaginationens Grundlag
er det Enkelte, abstract som enkelt, der ikke existerer for
Fornuften; den abstracte Erkjendelse forudsætter et saadant Enkelt,
om end kun som forladende det; den vilde ellers ikke blive abstract,
universel, men individuel og intuitiv. Det Ubestemte kommer paa
denne Maade til at staae i et tvivlsomt Forhold til det Uendelige,
»det sande Abstracte« i et tvivlsomt Forhold til Imaginationens
usande Abstraction; men denne Usikkerhed og Modsigelse var
Consequentsen af Usikkerheden i det metaphysiske Forhold mellem det
Uendelige og det Endelige. Efter den consequente Opfattelse af
Substantsbegrebet maatte den anden Erkjendelsesart forsvinde i den
tredie; men i denne Uendelighed forsvinder den individuelle Tanke
selv i det Uendelige, Philosophien bliver uforklarlig, den
philosopherende Tanke kommer i Modsigelse med sig selv, og denne
Modsigelse tiltvinger det Endelige en Art Bestaaen, der dog ikke kan
bestemmes uden at nye Modsigelser komme frem.

Læren om \emph{Erkjendelsens} forskjellige Former er afgjørende for
Spinozas Opfattelse af den menneskelige Aand,
% sentence continues on p. 105 ("...menneskelige Aand, hvis Væsen
% bestemtes..."), so no blank line here.
% --- p. 105 ---
\setcounter{footnote}{0}%
hvis Væsen bestemtes ved Erkjendelsen, og for hvilken der intet
Individuationsprincip gaves. \emph{Ideen er}, som Tankemodus, efter sin
Natur tidligere end de andre Tankemodi\footnote{Eth. II. axiom. 3.:
Modi cogitandi, ut amor, cupiditas, vel quicunque nomine affectus
animi insigniuntur, non dantur nisi in eodem individuo detur idea
rei amatæ, desideratæ \&c. At idea dari potest, quamvis nullus alius
% closing „..." mark printed with no opening mark anywhere in this
% note -- kept as printed, not balanced (cf. p. 45's dropped closer).
detur cogitandi modus.“}. Idet den opfattes som activ, indbefatter
den \emph{Villien} i sig: »Villie og Forstand ere Eet og det Samme«
(Eth. II, 49. Coroll.). Villien (voluntas) bestemmes nemlig af
Spinoza som facultas affirmandi et negandi ɔ: som den Evne, ved
hvilken Aanden bekræfter og benægter det Sande og Falske (ikke som
den Attraa, ved hvilken Aanden attraaer eller afskyer en Ting). Men
Bestemmelsen: \emph{Evne}, er kun et ens metaphysicum sive universale,
hvilket vi danne ved Abstraction fra det Enkelte; og voluntas kommer
derfor til at forholde sig til de enkelte volitiones som lapideitas
til lapis ɔ: den er i Realiteten ikke forskjellig fra dem. Men idet
Villien saaledes bliver den enkelte Affirmation eller Negation,
bliver den liig Ideen, idet der i Aanden ingen anden Affirmation og
Negation gives end den, Ideen, som Idee, indeslutter. Saaledes er f.
Ex. den volitio, ved hvilken Aanden bekræfter, at Triangelens tre
Vinkler ere liig to Rette, ikke forskjellig fra den Affirmation, der
ligger i Triangelens Idee, hvilken Idee er Betingelsen for hiin
Affirmation, og atter ikke kan tænkes uden denne, saa at deres
Begreb falder sammen (Eth. II, 49. dem.). Den enkelte volitio bliver
altsaa liig med den enkelte Idee, voluntas med intellectus.
Opfattede som enkelte maae volitiones, ligesom enhver anden enkelt
Ting, bestemmes til Existents og Virken ved en anden enkelt Ting; de
ere altsaa nødvendige (Eth. I, 32), og den menneskelige Aand,
opfattet som en bestemt og begrændset Tankemodus, kan ligesaalidt
have en absolut Evne til at ville og ikke ville (en absolut og fri
Villie), som den under den samme Opfattelse kan have en absolut Evne
til at begribe (Eth. II, 48); den kan ikke blive fri Aarsag
% sentence continues on p. 106 ("...fri Aarsag (causa libera) til
% sine Handlinger..."), so no blank line here.
% --- p. 106 ---
\setcounter{footnote}{0}%
(causa libera) til sine Handlinger, men bestemmes udenfra med
Nødvendighed.

I Opfattelsen af Villien maa derfor Spinoza bestemt adskille sig fra
Descartes, der, idet han i Gud satte en Vilkaarlighedens Frihed,
tillige, trods sin Grundanskuelses Consequentser, indrømmede
Mennesket en lignende, og ved at bestemme den menneskelige Villie
som uendelig, altsaa videre i Omfang end den begrændsede
Erkjendelse, søgte en Forklaringsgrund for Vildfarelsen, hvorved
denne i Forhold til Gud kunde opfattes som en blot Negation og kun i
Forhold til Mennesket blev en Privation. Spinoza indrømmer vel, at
Villien har et videre Omfang end de klare og tydelige Ideer; men
ikke end perceptiones sive concipiendi facultas overhovedet, idet
ogsaa denne, paa samme Maade som Villien, \emph{successivt} kan
optage Uendeligt. Opfattelsen af Villien som uendelig udleder han af
den misforstaaede Brug af det abstracte Begreb: Evne, der anvendes
paa den; thi idet det Abstracte, Universelle, ligemeget udsiges om
een, om flere og om uendelig mange Enkeltheder, hvis Fælles det
udtrykker, kan man give det en uendelig Udstrækning ud over den
bestemte Erkjendelses Gebeet; og den samme abstracte
Opfattelse\footnote{Eth. II, 49. schol.: hic apprime venit notandum,
quam facile decipimur, quando universalia cum singularibus et
entia rationis et abstracta cum realibus confundimus.} bevirker
ligeledes, at Villien, bestemt som Affirmation, udsiges i
Almindelighed om enhver Idee uden Hensyn til dens Indhold, og om det
Sande og Falske, saaledes at vi med den \emph{samme} Villies Magt bekræfte
alle Ideer, og udsige det Sande og det Falske som sandt, medens dog
Ideerne efter deres Væsen indeholde en forskjellig Grad af Realitet,
saa at altsaa de i dem indeholdte Affirmationer ogsaa maae indeslutte
en lignende, og medens det Sande og det Falske forholde sig til
hinanden som ens til non-ens. Ogsaa Evnen til, vilkaarlig at
suspendere sin Dom, benægter Spinoza, idet det at suspendere sin Dom
ikke bliver Andet, end at vi
% sentence continues on p. 107 ("...end at vi indsee, at vi ikke
% begribe Tingen adæquat."), so no blank line here.
% --- p. 107 ---
\setcounter{footnote}{0}%
indsee, at vi ikke begribe Tingen adæquat. Selv i Imaginationen
ligger der en Affirmation, saalænge ikke en anden Forestilling
udelukker den forestillede Gjenstands Existents. Saalænge dette
ikke finder Sted, kan man ikke tvivle om sin Forestillings Realitet,
skjøndt det: \emph{ikke at tvivle}, ikke er enstydigt med det:
\emph{at have Vished}, hvilken kun den sande Idee giver. Forestiller
jeg mig f. Ex. en bevinget Hest, da udsiger (bekræfter) jeg Vinger
om Hesten, og jeg vil, efter Imaginationens Natur, betragte den
bevingede Hest som existerende saalænge jeg har denne Forestilling
alene, og først tvivle om Forestillingen eller benægte den, naar en
anden Idee ophæver denne Forestilling eller viser den som uadæquat;
men da vil min Tvivl og min Benægtelse være nødvendig. --- Den Grund
for Villiens Frihed, der hentes fra, at den er uundværlig forat
ophæve den fuldkomne Indifferentstilstand, afviser Spinoza, med hvis
Opfattelse af Tænkningens Activitet den er uforenelig, spottende:
dico, me omnino concedere, quod homo, in tali æquilibrio [som
Buridans Æsel] positus (nempe qui nihil aliud percipit, quam sitim
et famem, talem cibum et talem potum, qui æque ab eo distant) fame
et siti peribit. Si me rogant, an talis homo non potius asinus quam
homo sit æstimandus, dico me nescire, ut etiam nescio, quanti
æstimandus sit ille, qui se pensilem facit, et quanti æstimandi sint
pueri stulti, vesani \&c.

Efterat Spinoza, i Ethikens 2den Bog, har udviklet sin Lære om den
menneskelige Aands Natur, fremhæver han, imod de Indvendinger, der
maatte ligge hans Samtidige nær, Fordelene ved sin Opfattelse, og
han sætter disse i: 1) at den, ved at lære os, at vi i vor Handlen
ere ubetinget afhængige af Gud, og i vort Væsen deelagtige i den
guddommelige Natur, og det i saameget desto høiere Grad jo
fuldkomnere vi handle og jo mere vi begribe Gud, --- beroliger vort
Sind og viser os, hvori vor høieste Lykke eller Salighed bestaaer,
nemlig \emph{alene i Erkjendelsen} af Gud, ved hvilken vi føres til
blot at gjøre det, som Kjærlighed og Fromhed byder, og at betragte
selve Dyden og Af\-%
% Cap. 5 does NOT end at the foot of p. 107: the chapter runs on into
% p. 108 (RECON.md line 426, Indhold "Cap. 5 ... 82--108"; confirmed
% here by the image, which shows body text running off the page
% mid-word, no footnote block and no closing rule on p. 107). The
% word broken here is completed at the head of p. 108, in the next
% batch's fragment; join with no space, per BATCH-AGENT.md's
% word-break rule. p. 107 itself carries no footnote, as the recon
% survey predicted.


% =====================================================================
%  SJETTE CAPITEL  (printed pp. 108--131 = PDF 118--141)
%  Opens PART WAY DOWN p. 108, so no \chapter* here -- the head is emitted
%  inline inside the fragment below (BATCH-AGENT.md §Heads).
%  NB the second in-text enumeration 1) 2) 3) 4) on p. 108 -- the four uses
%  of the doctrine. Not footnotes.
%  p. 130 note 2 CONTINUES at the head of p. 131's note block.
% =====================================================================

% [text to be added: pp. 108--119]

% [text to be added: pp. 120--131]


% =====================================================================
%  SYVENDE CAPITEL  (printed pp. 132--161 = PDF 142--171)
%  Opens PART WAY DOWN p. 132, so no \chapter* here -- the head is emitted
%  inline inside the fragment below (BATCH-AGENT.md §Heads).
%  The densest footnote stretch in the book: p. 148 carries three notes and
%  a genuine mixed quotation pair »…“; p. 156 carries five.
% =====================================================================

% [text to be added: pp. 132--142]

% [text to be added: pp. 143--152]

% [text to be added: pp. 153--161]

\newpage

% =====================================================================
%  OTTENDE CAPITEL  (printed pp. 162--170 = PDF 172--180)
%  Opens at the HEAD of p. 162 -- the only mid-book chapter that does.
% =====================================================================
\chapter*{Ottende Capitel.}
\addcontentsline{toc}{chapter}{Ottende Capitel. Menneskets Frihed og Evighed}
\markboth{Menneskets Frihed og Evighed}{Menneskets Frihed og Evighed}

\begin{center}
\textbf{Menneskets Frihed og Evighed.}
\end{center}

% [text to be added: pp. 162--170]


% =====================================================================
%  NIENDE CAPITEL  (printed pp. 171--180 = PDF 181--190)
%  Opens PART WAY DOWN p. 171, so no \chapter* here -- the head is emitted
%  inline inside the fragment below (BATCH-AGENT.md §Heads). The Indhold's
%  `171' for this chapter is the damaged sort noted above.
% =====================================================================

% [text to be added: pp. 171--180]


% =====================================================================
%  TIENDE CAPITEL  (printed pp. 181--188 = PDF 191--198)
%  Opens PART WAY DOWN p. 181, so no \chapter* here -- the head is emitted
%  inline inside the fragment below (BATCH-AGENT.md §Heads).
%  The book ends on p. 188; there is no colophon, no errata leaf and no
%  advertisement. The BSB ink stamp at the foot of p. 188 is not text.
% =====================================================================

% [text to be added: pp. 181--188]

\end{document}
